% =============================================================================
%  FLUXOGRAMA DA ESTEIRA — GERADO POR scripts/gerar_fluxograma.py. NÃO EDITE À MÃO:
%  mude o código (metadados, prompts, árvore) e rode o script de novo.
%  Compilar com:  tectonic -X compile fluxograma-dos-agentes.tex
% =============================================================================
\documentclass[10pt,a4paper]{article}
\usepackage[brazil]{babel}
\usepackage[default]{sourcesanspro}
\usepackage{fontspec}
\setsansfont{Montserrat-Regular.otf}[BoldFont=Montserrat-Bold.otf,ItalicFont=Montserrat-Italic.otf,BoldItalicFont=Montserrat-BoldItalic.otf]
\setmonofont{SourceCodePro-Regular.otf}[BoldFont=SourceCodePro-Bold.otf,Scale=0.90]
\usepackage{microtype}
\usepackage{amssymb}
\usepackage[a4paper,top=2.0cm,bottom=1.8cm,left=2.6cm,right=1.6cm,headheight=14pt,footskip=0.9cm]{geometry}
\usepackage[table]{xcolor}
\usepackage{graphicx,booktabs,tabularx,array,fancyhdr,enumitem}
\usepackage{tikz}
\usetikzlibrary{positioning,arrows.meta,calc,fit,backgrounds}
\usepackage[hidelinks,pdfusetitle]{hyperref}
\definecolor{unNavy}{HTML}{0E1039}
\definecolor{unLavanda}{HTML}{8B85AC}
\definecolor{unAzul}{HTML}{7FA5E2}
\definecolor{unPapel}{HTML}{F5F5FA}
\definecolor{unLinha}{HTML}{D9D7E4}
\definecolor{unTexto}{HTML}{2B2D45}
\definecolor{sinalAmbar}{HTML}{B07A22}
\color{unTexto}
\setlength{\parindent}{0pt}
\newcommand{\cod}[1]{{\ttfamily\footnotesize #1}}
\newcommand{\arq}[1]{{\sffamily\tiny\color{unLavanda}[#1]}}
\newcommand{\campo}[1]{{\sffamily\bfseries\scriptsize\color{unNavy}#1:}\ }
\newcommand{\ciclo}[1]{{\sffamily\scriptsize\bfseries\color{sinalAmbar}#1}}
\pagestyle{fancy}\fancyhf{}
\renewcommand{\headrulewidth}{0.4pt}
\fancyhead[L]{\sffamily\scriptsize\color{unNavy}\textbf{unooria} \textbar\ Fluxograma da esteira}
\fancyhead[R]{\sffamily\scriptsize\color{unLavanda}\leftmark}
\fancyfoot[R]{\sffamily\scriptsize\color{unNavy}\thepage}
\fancyfoot[L]{\sffamily\scriptsize\color{unLavanda}Gerado do código em 24/10/2026 · complementa o Manual do Gerador de Missões}
\tikzset{
  passo/.style={rectangle,rounded corners=3pt,draw=unLinha,line width=0.6pt,fill=white,text width=14.4cm,
    align=left,inner sep=5pt,font=\scriptsize},
  pmodelo/.style={passo,draw=unLavanda,fill=unLavanda!7},
  pcodigo/.style={passo,draw=unNavy!55,fill=unPapel},
  pportao/.style={passo,draw=unAzul!80!unNavy,fill=unAzul!10,line width=0.9pt},
  psaida/.style={passo,draw=sinalAmbar,fill=sinalAmbar!7,dashed},
  fluxo/.style={-{Stealth[length=2mm]},draw=unNavy!70,line width=0.7pt},
  loopfaixa/.style={draw=sinalAmbar,line width=1.1pt},
  loopvolta/.style={-{Stealth[length=2mm]},draw=sinalAmbar,line width=0.8pt,dashed,rounded corners=2pt},
  contlabel/.style={font=\sffamily\tiny\itshape,text=sinalAmbar,inner sep=1pt},
}
\begin{document}

\thispagestyle{empty}
{\sffamily\color{unNavy}\Huge\bfseries Fluxograma da esteira}\\[0.4em]
{\sffamily\color{unLavanda}\Large os 129 passos numerados (152 blocos), na ordem em que rodam}\\[1.2em]
Cada bloco é um agente ou portão da árvore real do orquestrador: o que faz, para quê existe, o que lê do estado da sessão e o que grava --- com o arquivo da sessão em que aparece, entre colchetes. Este documento é \textbf{gerado do código} por \cod{scripts/gerar\_fluxograma.py} e complementa o \emph{Manual do Gerador de Missões}, que explica o porquê de cada regra; aqui está o mapa passo a passo.\\[1em]
\begin{tikzpicture}[node distance=0.25cm]
\node[pmodelo,text width=3.4cm] (a) {\textbf{MODELO} --- agente de linguagem (53)};
\node[pcodigo,text width=3.4cm,right=of a] (b) {\textbf{CÓDIGO} --- passo determinístico (23)};
\node[pportao,text width=3.4cm,right=of b] (c) {\textbf{PORTÃO} --- grava um veredito (63)};
\node[psaida,text width=3.4cm,right=of c] (d) {\textbf{SAÍDA DE LOOP} --- repete ou segue (13)};
\end{tikzpicture}\\[0.8em]
{\scriptsize Cada \textcolor{sinalAmbar}{\textbf{linha âmbar}} à esquerda é um ciclo de correção --- uma por nível de aninhamento ---, e a seta tracejada é a volta quando o ciclo inteiro cabe na página. O bloco diz onde o ciclo começa ($\circlearrowright$), a que ciclos pertence e para onde volta quando reprova ($\circlearrowleft$). ``de controle'' são contadores, envelopes e marcas de execução, que não são artefato.}\\[1em]
{\sffamily\bfseries\color{unNavy} As etapas}\\[0.3em]
\begin{tabularx}{\textwidth}{l r X}\toprule Etapa & Blocos & Passos \\ \midrule
Produto & 77 & 001 a 063 \\
Design & 34 & 064 a 094 \\
Dev & 38 & 095 a 126 \\
Transversal & 3 & 127 a 129 \\
\bottomrule\end{tabularx}\\[1em]
{\sffamily\bfseries\color{unNavy} Os ciclos de correção}\\[0.3em]
{\footnotesize\begin{tabularx}{\textwidth}{l c X}\toprule Loop & Teto & Passos \\ \midrule
\cod{ComiteFinalLoop} & 2 & 004 a 042 (48 blocos) \\
\cod{ComiteCanaisLoop} & 5 & 004 a 022 (24 blocos) \\
\cod{NucleoNarrativoLoop} & 2 & 004 a 013 (10 blocos) \\
\cod{SecoesApoioLoop} & 2 & 014 a 016 (3 blocos) \\
\cod{GateEstruturalLoop} & 3 & 023 a 028 (6 blocos) \\
\cod{AvaliacoesLoop} & 2 & 029a a 034 (8 blocos) \\
\cod{OnePagerPoLoop} & 2 & 035 a 037 (3 blocos) \\
\cod{SpecLoop} & 2 & 045 a 059 (19 blocos) \\
\cod{CompletudeLoop} & 2 & 045b a 045d (3 blocos) \\
\cod{OnePagerResumidoLoop} & 2 & 061 a 063 (3 blocos) \\
\cod{QaDesignLoop} & 2 & 085 a 091 (8 blocos) \\
\cod{HandoffLoop} & 2 & 100a a 103 (6 blocos) \\
\cod{BancoLoop} & 2 & 104 a 106 (3 blocos) \\
\cod{BancaLoop} & 2 & 108 a 112 (6 blocos) \\
\cod{PrdLoop} & 2 & 113 a 118 (7 blocos) \\
\cod{OnePagerPedagogicoLoop} & 2 & 120 a 123 (4 blocos) \\
\bottomrule\end{tabularx}}
\newpage
\markboth{Produto}{}
{\sffamily\bfseries\Large\color{unNavy}Produto}\\[0.3em]
\begin{center}
\begin{tikzpicture}[node distance=0.42cm]
\node[pcodigo] (n0) {{\sffamily\bfseries\small 001\hspace{0.5em}Missões já geradas (determinístico)}\hfill{\sffamily\tiny\color{white}\colorbox{unNavy}{CÓDIGO}}\\[1pt]{\tiny\ttfamily\color{unLavanda}A001VizinhasCatalogo}\\[2pt]\campo{Faz}Lê o índice de missões e entrega ao núcleo as mais próximas do ano do briefing, para a missão nova não repetir uma anterior.\\[1pt]\campo{Objetivo}Põe as missões já geradas na frente de quem está inventando uma nova.\\[1pt]\campo{Entra}\cod{nucleo\_narrativo}\,\arq{produto/00a-nucleo-narrativo.md}, \cod{briefing}\\[1pt]\campo{Sai}\cod{missoes\_vizinhas}};
\node[pmodelo, below=of n0] (n1) {{\sffamily\bfseries\small 002\hspace{0.5em}Mapa conceitual do material}\hfill{\sffamily\tiny\color{white}\colorbox{unLavanda}{MODELO}}\\[1pt]{\tiny\ttfamily\color{unLavanda}A002MapaConceitual}\\[2pt]\campo{Faz}Lê o material do bimestre antes de inventar: um conceito limiar, os conceitos núcleo e de apoio com o trecho citado, as crenças de entrada que o bimestre derruba e o vocabulário do Explicar.\\[1pt]\campo{Objetivo}Sua função é uma só: ler o material do bimestre e dizer, com fidelidade, o que ele ensina — antes que alguém invente a missão.\\[1pt]\campo{Entra}\cod{briefing}, \cod{material\_anexado}\,\arq{produto/00-material-anexado.md}\\[1pt]\campo{Sai}\cod{mapa\_conceitual}\,\arq{produto/00g-mapa-conceitual.md}};
\draw[fluxo] (n0.south) -- (n1.north);
\node[pmodelo, below=of n1] (n2) {{\sffamily\bfseries\small 003\hspace{0.5em}Ideação de mecanismos}\hfill{\sffamily\tiny\color{white}\colorbox{unLavanda}{MODELO}}\\[1pt]{\tiny\ttfamily\color{unLavanda}A003IdeacaoMecanismos}\\[2pt]\campo{Faz}Três mecanismos distintos para o conceito limiar, cada um com a anatomia do aha, pontuados na rubrica de profundidade. O núcleo desenvolve o escolhido.\\[1pt]\campo{Objetivo}Sua função: antes de alguém escrever a missão, pensar três desenhos diferentes para ela e escolher o mais profundo.\\[1pt]\campo{Entra}\cod{briefing}, \cod{mapa\_conceitual}\,\arq{produto/00g-mapa-conceitual.md}, \cod{material\_anexado}\,\arq{produto/00-material-anexado.md}, \cod{missoes\_vizinhas}\\[1pt]\campo{Sai}\cod{ideacao\_mecanismos}\,\arq{produto/00h-ideacao-mecanismos.md}};
\draw[fluxo] (n1.south) -- (n2.north);
\node[pmodelo, below=of n2] (n3) {{\sffamily\bfseries\small 004\hspace{0.5em}Núcleo narrativo}\hfill{\sffamily\tiny\color{white}\colorbox{unLavanda}{MODELO}}\\[1pt]{\tiny\ttfamily\color{unLavanda}A004NucleoNarrativo}\\[2pt]\ciclo{$\circlearrowright$ começa aqui: \cod{ComiteFinalLoop} (até 2 rodadas), \cod{ComiteCanaisLoop} (até 5 rodadas), \cod{NucleoNarrativoLoop} (até 2 rodadas)}\\[1pt]\campo{Faz}Seções 0, 1, 3 e 4 a partir do briefing, dos anexos, do mapa conceitual e do mecanismo escolhido: descoberta, conceitos vividos em tela, anatomia do aha, narrativa e as 6 etapas.\\[1pt]\campo{Objetivo}Sua função é uma só: dado um briefing, ideia ou one-pager bruto do time, gerar o núcleo narrativo da missão — as 4 seções que carregam toda a lógica criativa.\\[1pt]\campo{Entra}\cod{cobertura\_report}\,\arq{produto/00i-cobertura-conceitual.md}, \cod{comite\_canal\_a\_report}\,\arq{produto/00l-comite-canais-painel-a.md}, \cod{comite\_canal\_b\_report}\,\arq{produto/00m-comite-canais-painel-b.md}, \cod{comite\_final\_report}\,\arq{produto/03d-comite-final.md}, \cod{evidencia\_one\_pager\_report}\,\arq{produto/00k-contrato-evidencia-12a.md}, \cod{fidelidade\_nucleo\_report}, \cod{ideacao\_mecanismos}\,\arq{produto/00h-ideacao-mecanismos.md}, \cod{mapa\_conceitual}\,\arq{produto/00g-mapa-conceitual.md}, \cod{mecanica\_report}, +4\\[1pt]\campo{Sai}\cod{nucleo\_narrativo}\,\arq{produto/00a-nucleo-narrativo.md}\\[1pt]\campo{Ciclo}\cod{ComiteFinalLoop} $\supset$ \cod{ComiteCanaisLoop} $\supset$ \cod{NucleoNarrativoLoop}};
\draw[fluxo] (n2.south) -- (n3.north);
\node[pcodigo, below=of n3] (n4) {{\sffamily\bfseries\small 005\hspace{0.5em}Missões já geradas, pelo núcleo (determinístico)}\hfill{\sffamily\tiny\color{white}\colorbox{unNavy}{CÓDIGO}}\\[1pt]{\tiny\ttfamily\color{unLavanda}A005VizinhasCatalogo}\\[2pt]\campo{Faz}Refiltra o índice pelo ambiente, eixo e descoberta do núcleo recém-escrito. É esta lista que a verificação V5 do micro-gate julga.\\[1pt]\campo{Objetivo}Põe as missões já geradas na frente de quem está inventando uma nova.\\[1pt]\campo{Entra}\cod{nucleo\_narrativo}\,\arq{produto/00a-nucleo-narrativo.md}, \cod{briefing}\\[1pt]\campo{Sai}\cod{missoes\_vizinhas}\\[1pt]\campo{Ciclo}\cod{ComiteFinalLoop} $\supset$ \cod{ComiteCanaisLoop} $\supset$ \cod{NucleoNarrativoLoop}};
\draw[fluxo] (n3.south) -- (n4.north);
\node[pportao, below=of n4] (n5) {{\sffamily\bfseries\small 006\hspace{0.5em}Cobertura conceitual (determinístico)}\hfill{\sffamily\tiny\color{white}\colorbox{unAzul!80!unNavy}{PORTÃO}}\\[1pt]{\tiny\ttfamily\color{unLavanda}A006CoberturaConceitual}\\[2pt]\campo{Faz}Núcleo $\times$ mapa: o limiar vive no Decidir ou no aha, todo conceito núcleo tem destino, a anatomia do aha tem os seis passos e o trecho citado pelo mapa existe no material.\\[1pt]\campo{Objetivo}Confere o núcleo contra o mapa conceitual do material, a cada rodada do loop do núcleo.\\[1pt]\campo{Entra}\cod{nucleo\_narrativo}\,\arq{produto/00a-nucleo-narrativo.md}, \cod{mapa\_conceitual}\,\arq{produto/00g-mapa-conceitual.md}, \cod{material\_anexado}\,\arq{produto/00-material-anexado.md}\\[1pt]\campo{Sai}\cod{cobertura\_report}\,\arq{produto/00i-cobertura-conceitual.md}\\[1pt]\campo{Ciclo}\cod{ComiteFinalLoop} $\supset$ \cod{ComiteCanaisLoop} $\supset$ \cod{NucleoNarrativoLoop}};
\draw[fluxo] (n4.south) -- (n5.north);
\draw[loopfaixa] ([xshift=-0.25cm]n3.north west) -- ([xshift=-0.25cm]n5.south west);
\draw[loopfaixa] ([xshift=-0.55cm]n3.north west) -- ([xshift=-0.55cm]n5.south west);
\draw[loopfaixa] ([xshift=-0.85cm]n3.north west) -- ([xshift=-0.85cm]n5.south west);
\end{tikzpicture}
\end{center}
\newpage
\begin{center}
\begin{tikzpicture}[node distance=0.42cm]
\node[pmodelo] (n6) {{\sffamily\bfseries\small 007\hspace{0.5em}Micro-gate do núcleo}\hfill{\sffamily\tiny\color{white}\colorbox{unLavanda}{MODELO}}\\[1pt]{\tiny\ttfamily\color{unLavanda}A007MicroGate}\\[2pt]\campo{Faz}Spoiler, gap de informação, escolha significativa e job-to-be-done — binário, até 2 rodadas.\\[1pt]\campo{Objetivo}Você faz quatro verificações, todas binárias (passa/não passa).\\[1pt]\campo{Entra}\cod{missoes\_vizinhas}, \cod{nucleo\_narrativo}\,\arq{produto/00a-nucleo-narrativo.md}\\[1pt]\campo{Sai}\cod{micro\_gate\_report}\,\arq{produto/00b-micro-gate.md}\\[1pt]\campo{Ciclo}\cod{ComiteFinalLoop} $\supset$ \cod{ComiteCanaisLoop} $\supset$ \cod{NucleoNarrativoLoop}};
\node[pmodelo, below=of n6] (n7) {{\sffamily\bfseries\small 008\hspace{0.5em}Juiz de profundidade}\hfill{\sffamily\tiny\color{white}\colorbox{unLavanda}{MODELO}}\\[1pt]{\tiny\ttfamily\color{unLavanda}A008JuizProfundidade}\\[2pt]\campo{Faz}Nota de 1 a 4 em seis dimensões — conceito do material vivido, aposta comprometida, ápice do aha, explicar depois de viver, transferência e demanda à altura do ano —, com trecho citado e instrução de reescrita.\\[1pt]\campo{Objetivo}O micro-gate confere se a missão protege o gap de informação — se não há spoiler.\\[1pt]\campo{Entra}\cod{cobertura\_report}\,\arq{produto/00i-cobertura-conceitual.md}, \cod{mapa\_conceitual}\,\arq{produto/00g-mapa-conceitual.md}, \cod{material\_anexado}\,\arq{produto/00-material-anexado.md}, \cod{nucleo\_narrativo}\,\arq{produto/00a-nucleo-narrativo.md}\\[1pt]\campo{Sai}\cod{profundidade\_report}\,\arq{produto/00j-juiz-profundidade.md}\\[1pt]\campo{Ciclo}\cod{ComiteFinalLoop} $\supset$ \cod{ComiteCanaisLoop} $\supset$ \cod{NucleoNarrativoLoop}};
\draw[fluxo] (n6.south) -- (n7.north);
\node[pportao, below=of n7] (n8) {{\sffamily\bfseries\small 009\hspace{0.5em}Conta da virada (determinístico)}\hfill{\sffamily\tiny\color{white}\colorbox{unAzul!80!unNavy}{PORTÃO}}\\[1pt]{\tiny\ttfamily\color{unLavanda}A009ContaDaVirada}\\[2pt]\campo{Faz}Refaz o ranking da tabela 1.3 do núcleo: normaliza cada critério, aplica os pesos e confere que o vencedor antes e o depois são os declarados — e diferentes. Pega a virada que não fecha (o custo por vaga da exp09ef1b).\\[1pt]\campo{Objetivo}Refaz a conta da virada a partir da tabela 1.3 do núcleo (tools/mecanica.py).\\[1pt]\campo{Entra}\cod{nucleo\_narrativo}\,\arq{produto/00a-nucleo-narrativo.md}\\[1pt]\campo{Sai}\cod{mecanica\_report}, \cod{mecanica\_status}\\[1pt]\campo{Ciclo}\cod{ComiteFinalLoop} $\supset$ \cod{ComiteCanaisLoop} $\supset$ \cod{NucleoNarrativoLoop}};
\draw[fluxo] (n7.south) -- (n8.north);
\node[pportao, below=of n8] (n9) {{\sffamily\bfseries\small 010\hspace{0.5em}Fidelidade ao pedido do PO (núcleo)}\hfill{\sffamily\tiny\color{white}\colorbox{unAzul!80!unNavy}{PORTÃO}}\\[1pt]{\tiny\ttfamily\color{unLavanda}A010FidelidadeNucleo}\\[2pt]\campo{Faz}O núcleo contra todos os pedidos de refazer do PO (TRAVAS, AJUSTES, VALIDAÇÃO, frases entre aspas) e o briefing: o juiz cita, o código confere. Parcial: o núcleo não tem as outras seções, então item sem prova é  e só o pedido violado () devolve o núcleo ao autor; a prova é cobrada no 16c.\\[1pt]\campo{Objetivo}O que o PO pediu e o briefing fixou está nos documentos?\\[1pt]\campo{Entra}\cod{nucleo\_narrativo}\,\arq{produto/00a-nucleo-narrativo.md}, \cod{divisao\_do\_pedido}, \cod{pedidos\_do\_po}, \cod{pedido\_correcao}, \cod{etapa\_da\_correcao}, \cod{briefing}, \cod{edicoes\_pontuais}, \textit{+3 de controle}\\[1pt]\campo{Sai}\cod{fidelidade\_nucleo\_report}, \cod{fidelidade\_nucleo\_status}\\[1pt]\campo{Ciclo}\cod{ComiteFinalLoop} $\supset$ \cod{ComiteCanaisLoop} $\supset$ \cod{NucleoNarrativoLoop}};
\draw[fluxo] (n8.south) -- (n9.north);
\node[pportao, below=of n9] (n10) {{\sffamily\bfseries\small 011\hspace{0.5em}Veredito de profundidade (determinístico)}\hfill{\sffamily\tiny\color{white}\colorbox{unAzul!80!unNavy}{PORTÃO}}\\[1pt]{\tiny\ttfamily\color{unLavanda}A011ProfundidadeCheck}\\[2pt]\campo{Faz}Conta o veredito nas notas do juiz, não no que ele escreveu: P1, P2 ou P3 abaixo de 3 é . Funde a cobertura conceitual.\\[1pt]\campo{Objetivo}Veredito de profundidade contado em código: notas do juiz + cobertura determinística.\\[1pt]\campo{Entra}\cod{cobertura\_resultado}\\[1pt]\campo{Sai}\cod{profundidade\_status}, \cod{profundidade\_report}\,\arq{produto/00j-juiz-profundidade.md}, \textit{+1 de controle}\\[1pt]\campo{Ciclo}\cod{ComiteFinalLoop} $\supset$ \cod{ComiteCanaisLoop} $\supset$ \cod{NucleoNarrativoLoop}};
\draw[fluxo] (n9.south) -- (n10.north);
\node[pportao, below=of n10] (n11) {{\sffamily\bfseries\small 012\hspace{0.5em}Veredito do micro-gate (determinístico)}\hfill{\sffamily\tiny\color{white}\colorbox{unAzul!80!unNavy}{PORTÃO}}\\[1pt]{\tiny\ttfamily\color{unLavanda}A012MicroGateCheck}\\[2pt]\campo{Faz}Lê o envelope do micro-gate e registra se o núcleo passou em spoiler, gap, escolha e job.\\[1pt]\campo{Objetivo}Lê o veredito de um relatório e decide se o pipeline segue, repete ou para.\\[1pt]\campo{Entra}\cod{nucleo\_narrativo}\,\arq{produto/00a-nucleo-narrativo.md}, \cod{micro\_gate\_report}\,\arq{produto/00b-micro-gate.md}\\[1pt]\campo{Sai}\cod{micro\_gate\_status}, \textit{+1 de controle}\\[1pt]\campo{Ciclo}\cod{ComiteFinalLoop} $\supset$ \cod{ComiteCanaisLoop} $\supset$ \cod{NucleoNarrativoLoop}};
\draw[fluxo] (n10.south) -- (n11.north);
\draw[loopfaixa] ([xshift=-0.25cm]n6.north west) -- ([xshift=-0.25cm]n11.south west);
\draw[loopfaixa] ([xshift=-0.55cm]n6.north west) -- ([xshift=-0.55cm]n11.south west);
\draw[loopfaixa] ([xshift=-0.85cm]n6.north west) -- ([xshift=-0.85cm]n11.south west);
\end{tikzpicture}
\end{center}
\newpage
\begin{center}
\begin{tikzpicture}[node distance=0.42cm]
\node[psaida] (n12) {{\sffamily\bfseries\small 013\hspace{0.5em}Saída do loop do núcleo (determinístico)}\hfill{\sffamily\tiny\color{white}\colorbox{sinalAmbar}{SAÍDA DE LOOP}}\\[1pt]{\tiny\ttfamily\color{unLavanda}A013NucleoCheck}\\[2pt]\campo{Faz}Micro-gate e profundidade aprovados encerram o loop; qualquer um reprovado devolve o núcleo ao P0a com os dois relatórios.\\[1pt]\campo{Objetivo}Encerra um LoopAgent quando todos os gates da rodada passaram — ou quando o teto caiu.\\[1pt]\campo{Entra}\cod{micro\_gate\_status}, \cod{profundidade\_status}, \cod{fidelidade\_nucleo\_status}, \cod{mecanica\_status}\\[1pt]\campo{Sai}\textit{+1 de controle}\\[1pt]\campo{Ciclo}\cod{ComiteFinalLoop} $\supset$ \cod{ComiteCanaisLoop} $\supset$ \cod{NucleoNarrativoLoop}\\[1pt]\ciclo{$\circlearrowleft$ fim do ciclo: \cod{NucleoNarrativoLoop} --- reprovado, volta ao passo 004}};
\node[pmodelo, below=of n12] (n13) {{\sffamily\bfseries\small 014\hspace{0.5em}Seções de apoio e brief}\hfill{\sffamily\tiny\color{white}\colorbox{unLavanda}{MODELO}}\\[1pt]{\tiny\ttfamily\color{unLavanda}A014SecoesApoio}\\[2pt]\ciclo{$\circlearrowright$ começa aqui: \cod{SecoesApoioLoop} (até 2 rodadas)}\\[1pt]\campo{Faz}Completa o one-pager canônico (seções 2 e 5–18) e o brief de escopo sobre o núcleo aprovado.\\[1pt]\campo{Objetivo}Sua função: dado o núcleo narrativo validado (seções 0, 1, 3, 4 + decisões de direção), gerar as seções de apoio (2, 5–18) e o brief de escopo, completando o one-pager canônico da missão.\\[1pt]\campo{Entra}\cod{briefing}, \cod{comite\_canal\_a\_report}\,\arq{produto/00l-comite-canais-painel-a.md}, \cod{comite\_canal\_b\_report}\,\arq{produto/00m-comite-canais-painel-b.md}, \cod{comite\_final\_report}\,\arq{produto/03d-comite-final.md}, \cod{evidencia\_one\_pager\_report}\,\arq{produto/00k-contrato-evidencia-12a.md}, \cod{nucleo\_narrativo}\,\arq{produto/00a-nucleo-narrativo.md}, \cod{one\_pager\_po}, \cod{one\_pager\_po\_report}\\[1pt]\campo{Sai}\cod{one\_pager}\,\arq{produto/00-one-pager.md}\\[1pt]\campo{Ciclo}\cod{ComiteFinalLoop} $\supset$ \cod{ComiteCanaisLoop} $\supset$ \cod{SecoesApoioLoop}};
\draw[fluxo] (n12.south) -- (n13.north);
\node[pportao, below=of n13] (n14) {{\sffamily\bfseries\small 015\hspace{0.5em}One-pager do PO preservado (determinístico)}\hfill{\sffamily\tiny\color{white}\colorbox{unAzul!80!unNavy}{PORTÃO}}\\[1pt]{\tiny\ttfamily\color{unLavanda}A015OnePagerPoCheck}\\[2pt]\campo{Faz}Quando o PO anexou o one-pager: cada trecho dele (linha e célula de tabela) está no one-pager das seções de apoio, ou saiu porque o pedido mandou. Trecho perdido volta ao A014 com a lista. Sem anexo, não se aplica.\\[1pt]\campo{Objetivo}Gate determinístico de rastro: modo decide a verificação de tools/rastro.py.\\[1pt]\campo{Entra}\cod{one\_pager\_po}, \cod{briefing}, \cod{one\_pager}\,\arq{produto/00-one-pager.md}, \cod{brief\_escopo}\,\arq{produto/00f-brief-escopo.md}\\[1pt]\campo{Sai}\cod{one\_pager\_po\_report}, \cod{one\_pager\_po\_status}\\[1pt]\campo{Ciclo}\cod{ComiteFinalLoop} $\supset$ \cod{ComiteCanaisLoop} $\supset$ \cod{SecoesApoioLoop}};
\draw[fluxo] (n13.south) -- (n14.north);
\node[psaida, below=of n14] (n15) {{\sffamily\bfseries\small 016\hspace{0.5em}Saída das seções de apoio (determinístico)}\hfill{\sffamily\tiny\color{white}\colorbox{sinalAmbar}{SAÍDA DE LOOP}}\\[1pt]{\tiny\ttfamily\color{unLavanda}A016SecoesApoioRodada}\\[2pt]\campo{Faz}One-pager do PO preservado encerra o loop das seções de apoio; trecho perdido devolve ao A014.\\[1pt]\campo{Objetivo}Encerra um LoopAgent quando todos os gates da rodada passaram — ou quando o teto caiu.\\[1pt]\campo{Entra}\cod{one\_pager\_po\_status}, \cod{one\_pager}\,\arq{produto/00-one-pager.md}\\[1pt]\campo{Sai}\textit{+1 de controle}\\[1pt]\campo{Ciclo}\cod{ComiteFinalLoop} $\supset$ \cod{ComiteCanaisLoop} $\supset$ \cod{SecoesApoioLoop}\\[1pt]\ciclo{$\circlearrowleft$ fim do ciclo: \cod{SecoesApoioLoop} --- reprovado, volta ao passo 014}};
\draw[fluxo] (n14.south) -- (n15.north);
\node[pcodigo, below=of n15] (n16) {{\sffamily\bfseries\small 017\hspace{0.5em}One-pager existe (determinístico)}\hfill{\sffamily\tiny\color{white}\colorbox{unNavy}{CÓDIGO}}\\[1pt]{\tiny\ttfamily\color{unLavanda}A017OnePagerExiste}\\[2pt]\campo{Faz}Se o A014 terminou sem texto nas duas tentativas, a etapa para aqui com uma mensagem que diz quem ficou mudo — antes, o gate estrutural quebrava com «Context variable not found».\\[1pt]\campo{Objetivo}Para a etapa, com mensagem que se lê, quando um documento de que o resto depende não existe.\\[1pt]\campo{Entra}\cod{one\_pager}\,\arq{produto/00-one-pager.md}\\[1pt]\campo{Sai}\textit{nada do estado da sessão (configuração fixa do passo)}\\[1pt]\campo{Ciclo}\cod{ComiteFinalLoop} $\supset$ \cod{ComiteCanaisLoop}};
\draw[fluxo] (n15.south) -- (n16.north);
\node[pportao, below=of n16] (n17) {{\sffamily\bfseries\small 018\hspace{0.5em}Contrato de evidência do one-pager (determinístico)}\hfill{\sffamily\tiny\color{white}\colorbox{unAzul!80!unNavy}{PORTÃO}}\\[1pt]{\tiny\ttfamily\color{unLavanda}A018EvidenciaOnePager}\\[2pt]\campo{Faz}O \S{}12a onde ele nasce: toda habilidade da seção 2 com linha, Canal A e regra ternária por extenso, e Canal B com o nível de cada alternativa. Forma antes de voto: o que falha aqui volta ao autor sem gastar o comitê.\\[1pt]\campo{Objetivo}Gate determinístico de rastro: modo decide a verificação de tools/rastro.py.\\[1pt]\campo{Entra}\cod{one\_pager}\,\arq{produto/00-one-pager.md}\\[1pt]\campo{Sai}\cod{evidencia\_one\_pager\_report}\,\arq{produto/00k-contrato-evidencia-12a.md}, \cod{evidencia\_one\_pager\_status}\\[1pt]\campo{Ciclo}\cod{ComiteFinalLoop} $\supset$ \cod{ComiteCanaisLoop}};
\draw[fluxo] (n16.south) -- (n17.north);
\draw[loopfaixa] ([xshift=-0.25cm]n12.north west) -- ([xshift=-0.25cm]n17.south west);
\draw[loopfaixa] ([xshift=-0.55cm]n12.north west) -- ([xshift=-0.55cm]n17.south west);
\draw[loopfaixa] ([xshift=-0.85cm]n12.north west) -- ([xshift=-0.85cm]n12.south west);
\draw[loopfaixa] ([xshift=-0.85cm]n13.north west) -- ([xshift=-0.85cm]n15.south west);
\draw[loopvolta] ([xshift=-0.85cm]n15.south west) -- ++(-0.18,0) |- ([xshift=-0.85cm]n13.north west);
\end{tikzpicture}
\end{center}
\newpage
\begin{center}
\begin{tikzpicture}[node distance=0.42cm]
\node[pmodelo] (n18) {{\sffamily\bfseries\small 019a\hspace{0.5em}Comitê dos canais: A·1 Currículo}\hfill{\sffamily\tiny\color{white}\colorbox{unLavanda}{MODELO}}\\[1pt]{\tiny\ttfamily\color{unLavanda}A019aComiteCanalACurriculo}\\[2pt]\campo{Faz}Painel A: a ação do Canal A de cada habilidade é o verbo e o objeto do texto oficial, na decisão onde o limiar mora.\\[1pt]\campo{Objetivo}O painel A julga a mecânica da missão e a ação do Canal A de cada habilidade — a ação nativa do ambiente que o sistema observa e classifica em Atingiu, Em processo e Não atingiu.\\[1pt]\campo{Entra}\cod{briefing}, \cod{comite\_canal\_a\_report}\,\arq{produto/00l-comite-canais-painel-a.md}, \cod{evidencia\_one\_pager\_report}\,\arq{produto/00k-contrato-evidencia-12a.md}, \cod{mapa\_conceitual}\,\arq{produto/00g-mapa-conceitual.md}, \cod{nucleo\_narrativo}\,\arq{produto/00a-nucleo-narrativo.md}, \cod{one\_pager}\,\arq{produto/00-one-pager.md}\\[1pt]\campo{Sai}\cod{comite\_canal\_a\_curriculo}\\[1pt]\campo{Ciclo}\cod{ComiteFinalLoop} $\supset$ \cod{ComiteCanaisLoop}\\[1pt]\campo{Roda}em paralelo com os demais de \cod{A019ComiteCanais}};
\node[pmodelo, below=of n18] (n19) {{\sffamily\bfseries\small 019b\hspace{0.5em}Comitê dos canais: A·2 Medida}\hfill{\sffamily\tiny\color{white}\colorbox{unLavanda}{MODELO}}\\[1pt]{\tiny\ttfamily\color{unLavanda}A019bComiteCanalAMedida}\\[2pt]\campo{Faz}Painel A: a regra ternária é alcançável, sem buraco, sem Atingiu por sorte e calculável pelo que o evento registra.\\[1pt]\campo{Objetivo}O painel A julga a mecânica da missão e a ação do Canal A de cada habilidade.\\[1pt]\campo{Entra}\cod{comite\_canal\_a\_report}\,\arq{produto/00l-comite-canais-painel-a.md}, \cod{evidencia\_one\_pager\_report}\,\arq{produto/00k-contrato-evidencia-12a.md}, \cod{nucleo\_narrativo}\,\arq{produto/00a-nucleo-narrativo.md}, \cod{one\_pager}\,\arq{produto/00-one-pager.md}\\[1pt]\campo{Sai}\cod{comite\_canal\_a\_medida}\\[1pt]\campo{Ciclo}\cod{ComiteFinalLoop} $\supset$ \cod{ComiteCanaisLoop}\\[1pt]\campo{Roda}em paralelo com os demais de \cod{A019ComiteCanais}};
\draw[fluxo] (n18.south) -- (n19.north);
\node[pmodelo, below=of n19] (n20) {{\sffamily\bfseries\small 019c\hspace{0.5em}Comitê dos canais: A·3 Vivência e idade}\hfill{\sffamily\tiny\color{white}\colorbox{unLavanda}{MODELO}}\\[1pt]{\tiny\ttfamily\color{unLavanda}A019cComiteCanalAVivencia}\\[2pt]\campo{Faz}Painel A: a ação é verbo nativo, vivida na decisão e dentro dos tetos da faixa.\\[1pt]\campo{Objetivo}O painel A julga a mecânica da missão e a ação do Canal A de cada habilidade.\\[1pt]\campo{Entra}\cod{briefing}, \cod{comite\_canal\_a\_report}\,\arq{produto/00l-comite-canais-painel-a.md}, \cod{nucleo\_narrativo}\,\arq{produto/00a-nucleo-narrativo.md}, \cod{one\_pager}\,\arq{produto/00-one-pager.md}\\[1pt]\campo{Sai}\cod{comite\_canal\_a\_vivencia}\\[1pt]\campo{Ciclo}\cod{ComiteFinalLoop} $\supset$ \cod{ComiteCanaisLoop}\\[1pt]\campo{Roda}em paralelo com os demais de \cod{A019ComiteCanais}};
\draw[fluxo] (n19.south) -- (n20.north);
\node[pmodelo, below=of n20] (n21) {{\sffamily\bfseries\small 019d\hspace{0.5em}Comitê dos canais: B·1 Currículo}\hfill{\sffamily\tiny\color{white}\colorbox{unLavanda}{MODELO}}\\[1pt]{\tiny\ttfamily\color{unLavanda}A019dComiteCanalBCurriculo}\\[2pt]\campo{Faz}Painel B: a pergunta mede a habilidade da linha, e o Atingiu é domínio pleno e autônomo.\\[1pt]\campo{Objetivo}O painel B julga a pergunta do Canal B: uma pergunta de múltipla escolha, com três alternativas, cada uma mapeada a um nível (Atingiu · Em processo · Não atingiu), que aparece depois do Refletir e soma evidência ao Canal A da…\\[1pt]\campo{Entra}\cod{comite\_canal\_b\_report}\,\arq{produto/00m-comite-canais-painel-b.md}, \cod{evidencia\_one\_pager\_report}\,\arq{produto/00k-contrato-evidencia-12a.md}, \cod{nucleo\_narrativo}\,\arq{produto/00a-nucleo-narrativo.md}, \cod{one\_pager}\,\arq{produto/00-one-pager.md}\\[1pt]\campo{Sai}\cod{comite\_canal\_b\_curriculo}\\[1pt]\campo{Ciclo}\cod{ComiteFinalLoop} $\supset$ \cod{ComiteCanaisLoop}\\[1pt]\campo{Roda}em paralelo com os demais de \cod{A019ComiteCanais}};
\draw[fluxo] (n20.south) -- (n21.north);
\node[pmodelo, below=of n21] (n22) {{\sffamily\bfseries\small 019e\hspace{0.5em}Comitê dos canais: B·2 Medida}\hfill{\sffamily\tiny\color{white}\colorbox{unLavanda}{MODELO}}\\[1pt]{\tiny\ttfamily\color{unLavanda}A019eComiteCanalBMedida}\\[2pt]\campo{Faz}Painel B: as três alternativas discriminam os níveis e nenhuma se elimina pela linguagem.\\[1pt]\campo{Objetivo}O painel B julga a pergunta do Canal B das habilidades com canal A + B.\\[1pt]\campo{Entra}\cod{comite\_canal\_b\_report}\,\arq{produto/00m-comite-canais-painel-b.md}, \cod{evidencia\_one\_pager\_report}\,\arq{produto/00k-contrato-evidencia-12a.md}, \cod{nucleo\_narrativo}\,\arq{produto/00a-nucleo-narrativo.md}, \cod{one\_pager}\,\arq{produto/00-one-pager.md}\\[1pt]\campo{Sai}\cod{comite\_canal\_b\_medida}\\[1pt]\campo{Ciclo}\cod{ComiteFinalLoop} $\supset$ \cod{ComiteCanaisLoop}\\[1pt]\campo{Roda}em paralelo com os demais de \cod{A019ComiteCanais}};
\draw[fluxo] (n21.south) -- (n22.north);
\node[pmodelo, below=of n22] (n23) {{\sffamily\bfseries\small 019f\hspace{0.5em}Comitê dos canais: B·3 Experiência}\hfill{\sffamily\tiny\color{white}\colorbox{unLavanda}{MODELO}}\\[1pt]{\tiny\ttfamily\color{unLavanda}A019fComiteCanalBExperiencia}\\[2pt]\campo{Faz}Painel B: depois do Refletir sem repeti-lo, na língua do ambiente, sem vazar resultado e dentro do teto.\\[1pt]\campo{Objetivo}O painel B julga a pergunta do Canal B das habilidades com canal A + B.\\[1pt]\campo{Entra}\cod{comite\_canal\_b\_report}\,\arq{produto/00m-comite-canais-painel-b.md}, \cod{nucleo\_narrativo}\,\arq{produto/00a-nucleo-narrativo.md}, \cod{one\_pager}\,\arq{produto/00-one-pager.md}\\[1pt]\campo{Sai}\cod{comite\_canal\_b\_experiencia}\\[1pt]\campo{Ciclo}\cod{ComiteFinalLoop} $\supset$ \cod{ComiteCanaisLoop}\\[1pt]\campo{Roda}em paralelo com os demais de \cod{A019ComiteCanais}};
\draw[fluxo] (n22.south) -- (n23.north);
\draw[loopfaixa] ([xshift=-0.25cm]n18.north west) -- ([xshift=-0.25cm]n23.south west);
\draw[loopfaixa] ([xshift=-0.55cm]n18.north west) -- ([xshift=-0.55cm]n23.south west);
\end{tikzpicture}
\end{center}
\newpage
\begin{center}
\begin{tikzpicture}[node distance=0.42cm]
\node[pportao] (n24) {{\sffamily\bfseries\small 020\hspace{0.5em}Votos do painel A (determinístico)}\hfill{\sffamily\tiny\color{white}\colorbox{unAzul!80!unNavy}{PORTÃO}}\\[1pt]{\tiny\ttfamily\color{unLavanda}A020VotosCanalA}\\[2pt]\campo{Faz}Conta os três votos do painel A (mecânica e Canal A) com as citações conferidas: unanimidade aprova; sem ela, a missão volta ao núcleo (A004) e às seções de apoio (A014).\\[1pt]\campo{Objetivo}Conta os votos de um painel do comitê pedagógico (tools/comite.py).\\[1pt]\campo{Entra}\cod{nucleo\_narrativo}\,\arq{produto/00a-nucleo-narrativo.md}, \cod{one\_pager}\,\arq{produto/00-one-pager.md}, \cod{brief\_escopo}\,\arq{produto/00f-brief-escopo.md}, \cod{comite\_canal\_a\_curriculo}, \cod{comite\_canal\_a\_medida}, \cod{comite\_canal\_a\_vivencia}\\[1pt]\campo{Sai}\cod{comite\_canal\_a\_report}\,\arq{produto/00l-comite-canais-painel-a.md}, \cod{comite\_canal\_a\_status}, \textit{+1 de controle}\\[1pt]\campo{Ciclo}\cod{ComiteFinalLoop} $\supset$ \cod{ComiteCanaisLoop}};
\node[pportao, below=of n24] (n25) {{\sffamily\bfseries\small 021\hspace{0.5em}Votos do painel B (determinístico)}\hfill{\sffamily\tiny\color{white}\colorbox{unAzul!80!unNavy}{PORTÃO}}\\[1pt]{\tiny\ttfamily\color{unLavanda}A021VotosCanalB}\\[2pt]\campo{Faz}Conta os três votos do painel B (Canal B): unanimidade aprova; sem ela, a missão volta a quem a escreve. Sem Canal B na missão, os juízes são dispensados sem custo.\\[1pt]\campo{Objetivo}Conta os votos de um painel do comitê pedagógico (tools/comite.py).\\[1pt]\campo{Entra}\cod{nucleo\_narrativo}\,\arq{produto/00a-nucleo-narrativo.md}, \cod{one\_pager}\,\arq{produto/00-one-pager.md}, \cod{brief\_escopo}\,\arq{produto/00f-brief-escopo.md}, \cod{comite\_canal\_b\_curriculo}, \cod{comite\_canal\_b\_medida}, \cod{comite\_canal\_b\_experiencia}\\[1pt]\campo{Sai}\cod{comite\_canal\_b\_report}\,\arq{produto/00m-comite-canais-painel-b.md}, \cod{comite\_canal\_b\_status}, \textit{+1 de controle}\\[1pt]\campo{Ciclo}\cod{ComiteFinalLoop} $\supset$ \cod{ComiteCanaisLoop}};
\draw[fluxo] (n24.south) -- (n25.north);
\node[psaida, below=of n25] (n26) {{\sffamily\bfseries\small 022\hspace{0.5em}Saída do comitê dos canais (determinístico)}\hfill{\sffamily\tiny\color{white}\colorbox{sinalAmbar}{SAÍDA DE LOOP}}\\[1pt]{\tiny\ttfamily\color{unLavanda}A022ComiteCanaisRodada}\\[2pt]\campo{Faz}Os dois painéis unânimes e o \S{}12a em forma encerram o loop; senão, núcleo e seções de apoio reescrevem com os relatórios, até COMITE\_MAX\_RODADAS (padrão 5). Esgotado, segue  com os votos vencidos no parecer.\\[1pt]\campo{Objetivo}Encerra um LoopAgent quando todos os gates da rodada passaram — ou quando o teto caiu.\\[1pt]\campo{Entra}\cod{comite\_canal\_a\_status}, \cod{comite\_canal\_b\_status}, \cod{evidencia\_one\_pager\_status}, \cod{one\_pager}\,\arq{produto/00-one-pager.md}\\[1pt]\campo{Sai}\textit{+1 de controle}\\[1pt]\campo{Ciclo}\cod{ComiteFinalLoop} $\supset$ \cod{ComiteCanaisLoop}\\[1pt]\ciclo{$\circlearrowleft$ fim do ciclo: \cod{ComiteCanaisLoop} --- reprovado, volta ao passo 004}};
\draw[fluxo] (n25.south) -- (n26.north);
\node[pportao, below=of n26] (n27) {{\sffamily\bfseries\small 023\hspace{0.5em}Brief de escopo separado (determinístico)}\hfill{\sffamily\tiny\color{white}\colorbox{unAzul!80!unNavy}{PORTÃO}}\\[1pt]{\tiny\ttfamily\color{unLavanda}A023BriefEscopoSplit}\\[2pt]\ciclo{$\circlearrowright$ começa aqui: \cod{GateEstruturalLoop} (até 3 rodadas)}\\[1pt]\campo{Faz}Separa o brief de escopo do one-pager canônico para `brief\_escopo`, antes do revisor.\\[1pt]\campo{Objetivo}Copia seções de um relatório para chaves próprias do estado e registra um status informativo.\\[1pt]\campo{Entra}\cod{one\_pager}\,\arq{produto/00-one-pager.md}, \cod{notas\_de\_passagem}, \cod{pedidos\_do\_po}, \cod{spec\_missao\_gate0}\,\arq{produto/05a-spec-missao-gate0.md}, \cod{spec\_missao}\,\arq{produto/05-spec-missao.md}\\[1pt]\campo{Sai}\cod{brief\_escopo\_status}\\[1pt]\campo{Ciclo}\cod{ComiteFinalLoop} $\supset$ \cod{GateEstruturalLoop}};
\draw[fluxo] (n26.south) -- (n27.north);
\node[pportao, below=of n27] (n28) {{\sffamily\bfseries\small 024\hspace{0.5em}Contrato de plataforma e origem dos parâmetros (determinístico)}\hfill{\sffamily\tiny\color{white}\colorbox{unAzul!80!unNavy}{PORTÃO}}\\[1pt]{\tiny\ttfamily\color{unLavanda}A024ContratoProdutoCheck}\\[2pt]\campo{Faz}Lint do one-pager a cada rodada do gate: canal de sala por polling (nunca WebSocket), estado no servidor (sem localStorage/IndexedDB), SAEB só nos anos avaliados, número com origem ou marcador. O gate aplica; o revisor corrige.\\[1pt]\campo{Objetivo}Gate determinístico de rastro: modo decide a verificação de tools/rastro.py.\\[1pt]\campo{Entra}\cod{one\_pager}\,\arq{produto/00-one-pager.md}, \cod{nucleo\_narrativo}\,\arq{produto/00a-nucleo-narrativo.md}, \cod{briefing}, \cod{one\_pager\_po}, \cod{decisoes\_do\_po}\\[1pt]\campo{Sai}\cod{contrato\_produto\_report}\,\arq{produto/00e-contrato-produto.md}, \cod{contrato\_produto\_status}\\[1pt]\campo{Ciclo}\cod{ComiteFinalLoop} $\supset$ \cod{GateEstruturalLoop}};
\draw[fluxo] (n27.south) -- (n28.north);
\node[pportao, below=of n28] (n29) {{\sffamily\bfseries\small 025\hspace{0.5em}Juiz do contrato de produto}\hfill{\sffamily\tiny\color{white}\colorbox{unAzul!80!unNavy}{PORTÃO}}\\[1pt]{\tiny\ttfamily\color{unLavanda}A025JuizContratoProduto}\\[2pt]\campo{Faz}Juiz de achados sobre `Contrato de plataforma e origem dos parâmetros`: retira o achado de uso vedado quando o texto só cita para vedar — com citação literal conferida.\\[1pt]\campo{Objetivo}O regex acha, o LLM decide, o código confere a citação (tools/juizachados.py).\\[1pt]\campo{Entra}\cod{one\_pager}\,\arq{produto/00-one-pager.md}, \textit{+1 de controle}\\[1pt]\campo{Sai}\cod{contrato\_produto\_status}, \cod{contrato\_produto\_report}\,\arq{produto/00e-contrato-produto.md}\\[1pt]\campo{Ciclo}\cod{ComiteFinalLoop} $\supset$ \cod{GateEstruturalLoop}};
\draw[fluxo] (n28.south) -- (n29.north);
\draw[loopfaixa] ([xshift=-0.25cm]n24.north west) -- ([xshift=-0.25cm]n29.south west);
\draw[loopfaixa] ([xshift=-0.55cm]n24.north west) -- ([xshift=-0.55cm]n26.south west);
\draw[loopfaixa] ([xshift=-0.55cm]n27.north west) -- ([xshift=-0.55cm]n29.south west);
\end{tikzpicture}
\end{center}
\newpage
\begin{center}
\begin{tikzpicture}[node distance=0.42cm]
\node[pmodelo] (n30) {{\sffamily\bfseries\small 026\hspace{0.5em}Gate estrutural}\hfill{\sffamily\tiny\color{white}\colorbox{unLavanda}{MODELO}}\\[1pt]{\tiny\ttfamily\color{unLavanda}A026GateEstrutural}\\[2pt]\campo{Faz}critérios: estrutura, ambiente, logística, evidências, viabilidade e infraestrutura.\\[1pt]\campo{Objetivo}Responde duas perguntas: o documento da missão está completo?\\[1pt]\campo{Entra}\cod{briefing}, \cod{contrato\_produto\_report}\,\arq{produto/00e-contrato-produto.md}, \cod{one\_pager}\,\arq{produto/00-one-pager.md}, \cod{one\_pager\_po}, \textit{+1 de controle}\\[1pt]\campo{Sai}\cod{gate\_estrutural\_report}\,\arq{produto/01-gate-estrutural.md}\\[1pt]\campo{Ciclo}\cod{ComiteFinalLoop} $\supset$ \cod{GateEstruturalLoop}};
\node[pportao, below=of n30] (n31) {{\sffamily\bfseries\small 027\hspace{0.5em}Veredito do gate (determinístico)}\hfill{\sffamily\tiny\color{white}\colorbox{unAzul!80!unNavy}{PORTÃO}}\\[1pt]{\tiny\ttfamily\color{unLavanda}A027GateEstruturalCheck}\\[2pt]\campo{Faz}Lê o envelope JSON e decide repetir, seguir ou bloquear.\\[1pt]\campo{Objetivo}Lê o veredito de um relatório e decide se o pipeline segue, repete ou para.\\[1pt]\campo{Entra}\cod{one\_pager}\,\arq{produto/00-one-pager.md}, \cod{brief\_escopo}\,\arq{produto/00f-brief-escopo.md}, \cod{briefing}, \cod{gate\_estrutural\_report}\,\arq{produto/01-gate-estrutural.md}\\[1pt]\campo{Sai}\cod{gate\_estrutural\_status}, \textit{+1 de controle}\\[1pt]\campo{Ciclo}\cod{ComiteFinalLoop} $\supset$ \cod{GateEstruturalLoop}};
\draw[fluxo] (n30.south) -- (n31.north);
\node[pmodelo, below=of n31] (n32) {{\sffamily\bfseries\small 028\hspace{0.5em}Revisor do one-pager}\hfill{\sffamily\tiny\color{white}\colorbox{unLavanda}{MODELO}}\\[1pt]{\tiny\ttfamily\color{unLavanda}A028RevisorOnePager}\\[2pt]\campo{Faz}Aplica o complemento pronto-pra-colar do gate; não inventa decisão de produto.\\[1pt]\campo{Objetivo}O gate estrutural bloqueou a missão e escreveu, no bloco "Complemento — pronto pra colar", as correções prontas.\\[1pt]\campo{Entra}\cod{adequacao\_report}\,\arq{produto/03b-adequacao-pedagogica.md}, \cod{avaliacao\_pedagogica\_report}\,\arq{produto/02-avaliacao-pedagogica.md}, \cod{briefing}, \cod{contrato\_produto\_report}\,\arq{produto/00e-contrato-produto.md}, \cod{experiencia\_aluno\_report}\,\arq{produto/03-experiencia-aluno.md}, \cod{gate\_estrutural\_report}\,\arq{produto/01-gate-estrutural.md}, \cod{one\_pager\_po}\\[1pt]\campo{Sai}\cod{one\_pager}\,\arq{produto/00-one-pager.md}\\[1pt]\campo{Ciclo}\cod{ComiteFinalLoop} $\supset$ \cod{GateEstruturalLoop}\\[1pt]\ciclo{$\circlearrowleft$ fim do ciclo: \cod{GateEstruturalLoop} --- reprovado, volta ao passo 023}};
\draw[fluxo] (n31.south) -- (n32.north);
\node[pmodelo, below=of n32] (n33) {{\sffamily\bfseries\small 029a\hspace{0.5em}Avaliação pedagógica}\hfill{\sffamily\tiny\color{white}\colorbox{unLavanda}{MODELO}}\\[1pt]{\tiny\ttfamily\color{unLavanda}A029aAvaliacaoPedagogica}\\[2pt]\ciclo{$\circlearrowright$ começa aqui: \cod{AvaliacoesLoop} (até 2 rodadas)}\\[1pt]\campo{Faz}BNCC, professor, guardrails e affordance $\times$ guardrail (paralelo com P3).\\[1pt]\campo{Objetivo}Você responde: o conteúdo pedagógico é sólido, e o professor consegue conduzir a missão com o que está escrito?\\[1pt]\campo{Entra}\cod{briefing}, \cod{catalogo\_missoes}, \cod{gate\_estrutural\_report}\,\arq{produto/01-gate-estrutural.md}, \cod{material\_anexado}\,\arq{produto/00-material-anexado.md}, \cod{one\_pager}\,\arq{produto/00-one-pager.md}, \cod{one\_pager\_po}\\[1pt]\campo{Sai}\cod{avaliacao\_pedagogica\_report}\,\arq{produto/02-avaliacao-pedagogica.md}\\[1pt]\campo{Ciclo}\cod{ComiteFinalLoop} $\supset$ \cod{AvaliacoesLoop}\\[1pt]\campo{Roda}em paralelo com os demais de \cod{AvaliacoesParalelas}};
\draw[fluxo] (n32.south) -- (n33.north);
\node[pmodelo, below=of n33] (n34) {{\sffamily\bfseries\small 029b\hspace{0.5em}Experiência do estudante}\hfill{\sffamily\tiny\color{white}\colorbox{unLavanda}{MODELO}}\\[1pt]{\tiny\ttfamily\color{unLavanda}A029bExperienciaEstudante}\\[2pt]\campo{Faz}Carga cognitiva, faixa etária e agência (paralelo com P2).\\[1pt]\campo{Objetivo}Você responde: um estudante real dessa idade consegue viver essa experiência do jeito que foi desenhada, sem sobrecarga e sem ficar aquém do que é capaz?\\[1pt]\campo{Entra}\cod{briefing}, \cod{one\_pager}\,\arq{produto/00-one-pager.md}, \cod{one\_pager\_po}\\[1pt]\campo{Sai}\cod{experiencia\_aluno\_report}\,\arq{produto/03-experiencia-aluno.md}\\[1pt]\campo{Ciclo}\cod{ComiteFinalLoop} $\supset$ \cod{AvaliacoesLoop}\\[1pt]\campo{Roda}em paralelo com os demais de \cod{AvaliacoesParalelas}};
\draw[fluxo] (n33.south) -- (n34.north);
\node[pmodelo, below=of n34] (n35) {{\sffamily\bfseries\small 029c\hspace{0.5em}Adequação pedagógica}\hfill{\sffamily\tiny\color{white}\colorbox{unLavanda}{MODELO}}\\[1pt]{\tiny\ttfamily\color{unLavanda}A029cAdequacaoPedagogica}\\[2pt]\campo{Faz}Seis eixos com parecer e ajustes acionáveis: currículo, demanda cognitiva, carga acidental, viabilidade na escola pública, evidência de aprendizagem e equidade. Parecer entra como ressalva, não bloqueia.\\[1pt]\campo{Objetivo}Você avalia se uma atividade pedagógica de uma missão Unooria é adequada para a série a que se destina, considerando três coisas ao mesmo tempo: aderência à BNCC, compatibilidade com a capacidade cognitiva real dos estudantes e…\\[1pt]\campo{Entra}\cod{briefing}, \cod{gate\_estrutural\_report}\,\arq{produto/01-gate-estrutural.md}, \cod{one\_pager}\,\arq{produto/00-one-pager.md}, \cod{one\_pager\_po}\\[1pt]\campo{Sai}\cod{adequacao\_report}\,\arq{produto/03b-adequacao-pedagogica.md}\\[1pt]\campo{Ciclo}\cod{ComiteFinalLoop} $\supset$ \cod{AvaliacoesLoop}\\[1pt]\campo{Roda}em paralelo com os demais de \cod{AvaliacoesParalelas}};
\draw[fluxo] (n34.south) -- (n35.north);
\draw[loopfaixa] ([xshift=-0.25cm]n30.north west) -- ([xshift=-0.25cm]n35.south west);
\draw[loopfaixa] ([xshift=-0.55cm]n30.north west) -- ([xshift=-0.55cm]n32.south west);
\draw[loopfaixa] ([xshift=-0.55cm]n33.north west) -- ([xshift=-0.55cm]n35.south west);
\end{tikzpicture}
\end{center}
\newpage
\begin{center}
\begin{tikzpicture}[node distance=0.42cm]
\node[pportao] (n36) {{\sffamily\bfseries\small 030\hspace{0.5em}Veredito da avaliação pedagógica (determinístico)}\hfill{\sffamily\tiny\color{white}\colorbox{unAzul!80!unNavy}{PORTÃO}}\\[1pt]{\tiny\ttfamily\color{unLavanda}A030AvaliacaoCheck}\\[2pt]\campo{Faz}Lê o veredito da avaliação pedagógica (critérios 4.x, 5.1/5.2, 8.x).\\[1pt]\campo{Objetivo}Lê o veredito de um relatório e decide se o pipeline segue, repete ou para.\\[1pt]\campo{Entra}\cod{one\_pager}\,\arq{produto/00-one-pager.md}, \cod{brief\_escopo}\,\arq{produto/00f-brief-escopo.md}, \cod{material\_anexado}\,\arq{produto/00-material-anexado.md}, \cod{avaliacao\_pedagogica\_report}\,\arq{produto/02-avaliacao-pedagogica.md}\\[1pt]\campo{Sai}\cod{avaliacao\_status}\\[1pt]\campo{Ciclo}\cod{ComiteFinalLoop} $\supset$ \cod{AvaliacoesLoop}};
\node[pportao, below=of n36] (n37) {{\sffamily\bfseries\small 031\hspace{0.5em}Veredito da experiência do estudante (determinístico)}\hfill{\sffamily\tiny\color{white}\colorbox{unAzul!80!unNavy}{PORTÃO}}\\[1pt]{\tiny\ttfamily\color{unLavanda}A031ExperienciaCheck}\\[2pt]\campo{Faz}Lê o veredito da experiência do estudante (3.x, 6.x). Sinaliza, não trava: os tetos desta dimensão são [EXTRAPOLAÇÃO], não medida.\\[1pt]\campo{Objetivo}Lê o veredito de um relatório e decide se o pipeline segue, repete ou para.\\[1pt]\campo{Entra}\cod{one\_pager}\,\arq{produto/00-one-pager.md}, \cod{brief\_escopo}\,\arq{produto/00f-brief-escopo.md}, \cod{experiencia\_aluno\_report}\,\arq{produto/03-experiencia-aluno.md}\\[1pt]\campo{Sai}\cod{experiencia\_status}\\[1pt]\campo{Ciclo}\cod{ComiteFinalLoop} $\supset$ \cod{AvaliacoesLoop}};
\draw[fluxo] (n36.south) -- (n37.north);
\node[pportao, below=of n37] (n38) {{\sffamily\bfseries\small 032\hspace{0.5em}Veredito da adequação pedagógica (determinístico)}\hfill{\sffamily\tiny\color{white}\colorbox{unAzul!80!unNavy}{PORTÃO}}\\[1pt]{\tiny\ttfamily\color{unLavanda}A032AdequacaoCheck}\\[2pt]\campo{Faz}Lê o veredito do parecer de adequação pedagógica.\\[1pt]\campo{Objetivo}Lê o veredito de um relatório e decide se o pipeline segue, repete ou para.\\[1pt]\campo{Entra}\cod{one\_pager}\,\arq{produto/00-one-pager.md}, \cod{brief\_escopo}\,\arq{produto/00f-brief-escopo.md}, \cod{material\_anexado}\,\arq{produto/00-material-anexado.md}, \cod{adequacao\_report}\,\arq{produto/03b-adequacao-pedagogica.md}\\[1pt]\campo{Sai}\cod{adequacao\_status}\\[1pt]\campo{Ciclo}\cod{ComiteFinalLoop} $\supset$ \cod{AvaliacoesLoop}};
\draw[fluxo] (n37.south) -- (n38.north);
\node[psaida, below=of n38] (n39) {{\sffamily\bfseries\small 033\hspace{0.5em}Veredito do parecer (determinístico)}\hfill{\sffamily\tiny\color{white}\colorbox{sinalAmbar}{SAÍDA DE LOOP}}\\[1pt]{\tiny\ttfamily\color{unLavanda}A033AdequacaoCheck}\\[2pt]\campo{Faz}Parecer sem ajuste exigido encerra o ciclo das avaliações; com ajuste, devolve ao revisor do one-pager.\\[1pt]\campo{Objetivo}Encerra um LoopAgent quando todos os gates da rodada passaram — ou quando o teto caiu.\\[1pt]\campo{Entra}\cod{adequacao\_status}, \cod{avaliacao\_status}\\[1pt]\campo{Sai}\textit{+1 de controle}\\[1pt]\campo{Ciclo}\cod{ComiteFinalLoop} $\supset$ \cod{AvaliacoesLoop}};
\draw[fluxo] (n38.south) -- (n39.north);
\node[pmodelo, below=of n39] (n40) {{\sffamily\bfseries\small 034\hspace{0.5em}Revisor movido pelo parecer}\hfill{\sffamily\tiny\color{white}\colorbox{unLavanda}{MODELO}}\\[1pt]{\tiny\ttfamily\color{unLavanda}A034RevisorOnePagerParecer}\\[2pt]\campo{Faz}Aplica ao one-pager os ajustes exigidos pelo parecer de adequação (padrão: 2 rodadas; AVALIACOES\_MAX\_RODADAS=1 desliga).\\[1pt]\campo{Objetivo}O gate estrutural bloqueou a missão e escreveu, no bloco "Complemento — pronto pra colar", as correções prontas.\\[1pt]\campo{Entra}\cod{adequacao\_report}\,\arq{produto/03b-adequacao-pedagogica.md}, \cod{avaliacao\_pedagogica\_report}\,\arq{produto/02-avaliacao-pedagogica.md}, \cod{briefing}, \cod{contrato\_produto\_report}\,\arq{produto/00e-contrato-produto.md}, \cod{experiencia\_aluno\_report}\,\arq{produto/03-experiencia-aluno.md}, \cod{gate\_estrutural\_report}\,\arq{produto/01-gate-estrutural.md}, \cod{one\_pager\_po}\\[1pt]\campo{Sai}\cod{one\_pager}\,\arq{produto/00-one-pager.md}\\[1pt]\campo{Ciclo}\cod{ComiteFinalLoop} $\supset$ \cod{AvaliacoesLoop}\\[1pt]\ciclo{$\circlearrowleft$ fim do ciclo: \cod{AvaliacoesLoop} --- reprovado, volta ao passo 029a}};
\draw[fluxo] (n39.south) -- (n40.north);
\node[pportao, below=of n40] (n41) {{\sffamily\bfseries\small 035\hspace{0.5em}One-pager do PO preservado, depois dos pareceres (determinístico)}\hfill{\sffamily\tiny\color{white}\colorbox{unAzul!80!unNavy}{PORTÃO}}\\[1pt]{\tiny\ttfamily\color{unLavanda}A035OnePagerPoCheck}\\[2pt]\ciclo{$\circlearrowright$ começa aqui: \cod{OnePagerPoLoop} (até 2 rodadas)}\\[1pt]\campo{Faz}A conferência final, depois do gate estrutural, dos pareceres e dos revisores: trecho do one-pager do PO perdido sem pedido deixa produto  e vai ao restaurador.\\[1pt]\campo{Objetivo}Gate determinístico de rastro: modo decide a verificação de tools/rastro.py.\\[1pt]\campo{Entra}\cod{one\_pager\_po}, \cod{briefing}, \cod{one\_pager}\,\arq{produto/00-one-pager.md}, \cod{brief\_escopo}\,\arq{produto/00f-brief-escopo.md}\\[1pt]\campo{Sai}\cod{one\_pager\_po\_report}, \cod{one\_pager\_po\_status}\\[1pt]\campo{Ciclo}\cod{ComiteFinalLoop} $\supset$ \cod{OnePagerPoLoop}};
\draw[fluxo] (n40.south) -- (n41.north);
\node[psaida, below=of n41] (n42) {{\sffamily\bfseries\small 036\hspace{0.5em}Saída da conferência do one-pager do PO (determinístico)}\hfill{\sffamily\tiny\color{white}\colorbox{sinalAmbar}{SAÍDA DE LOOP}}\\[1pt]{\tiny\ttfamily\color{unLavanda}A036OnePagerPoRodada}\\[2pt]\campo{Faz}Preservado segue para o gate final; perdido vai ao restaurador uma vez.\\[1pt]\campo{Objetivo}Encerra um LoopAgent quando todos os gates da rodada passaram — ou quando o teto caiu.\\[1pt]\campo{Entra}\cod{one\_pager\_po\_status}, \cod{one\_pager}\,\arq{produto/00-one-pager.md}\\[1pt]\campo{Sai}\textit{+1 de controle}\\[1pt]\campo{Ciclo}\cod{ComiteFinalLoop} $\supset$ \cod{OnePagerPoLoop}};
\draw[fluxo] (n41.south) -- (n42.north);
\draw[loopfaixa] ([xshift=-0.25cm]n36.north west) -- ([xshift=-0.25cm]n42.south west);
\draw[loopfaixa] ([xshift=-0.55cm]n36.north west) -- ([xshift=-0.55cm]n40.south west);
\draw[loopfaixa] ([xshift=-0.55cm]n41.north west) -- ([xshift=-0.55cm]n42.south west);
\end{tikzpicture}
\end{center}
\newpage
\begin{center}
\begin{tikzpicture}[node distance=0.42cm]
\node[pmodelo] (n43) {{\sffamily\bfseries\small 037\hspace{0.5em}Restaurador do one-pager do PO}\hfill{\sffamily\tiny\color{white}\colorbox{unLavanda}{MODELO}}\\[1pt]{\tiny\ttfamily\color{unLavanda}A037RestaurarOnePagerPo}\\[2pt]\campo{Faz}Devolve ao one-pager, literalmente, os trechos do PO que a esteira reescreveu; o que um gate ou parecer pedia vira [PROPOSTA — validar com PO].\\[1pt]\campo{Objetivo}O PO anexou o one-pager dele, e uma conferência em código achou trechos que sumiram ou foram reescritos sem que o pedido do PO mandasse.\\[1pt]\campo{Entra}\cod{briefing}, \cod{comite\_canal\_a\_report}\,\arq{produto/00l-comite-canais-painel-a.md}, \cod{comite\_canal\_b\_report}\,\arq{produto/00m-comite-canais-painel-b.md}, \cod{comite\_final\_report}\,\arq{produto/03d-comite-final.md}, \cod{one\_pager\_po}, \cod{one\_pager\_po\_report}\\[1pt]\campo{Sai}\cod{one\_pager}\,\arq{produto/00-one-pager.md}\\[1pt]\campo{Ciclo}\cod{ComiteFinalLoop} $\supset$ \cod{OnePagerPoLoop}\\[1pt]\ciclo{$\circlearrowleft$ fim do ciclo: \cod{OnePagerPoLoop} --- reprovado, volta ao passo 035}};
\node[pportao, below=of n43] (n44) {{\sffamily\bfseries\small 038\hspace{0.5em}Contrato de evidência, antes do comitê final (determinístico)}\hfill{\sffamily\tiny\color{white}\colorbox{unAzul!80!unNavy}{PORTÃO}}\\[1pt]{\tiny\ttfamily\color{unLavanda}A038EvidenciaOnePagerFinal}\\[2pt]\campo{Faz}O \S{}12a de novo, sobre o one-pager depois do gate estrutural, dos pareceres e dos revisores.\\[1pt]\campo{Objetivo}Gate determinístico de rastro: modo decide a verificação de tools/rastro.py.\\[1pt]\campo{Entra}\cod{one\_pager}\,\arq{produto/00-one-pager.md}\\[1pt]\campo{Sai}\cod{evidencia\_one\_pager\_report}\,\arq{produto/00k-contrato-evidencia-12a.md}, \cod{evidencia\_one\_pager\_status}\\[1pt]\campo{Ciclo}\cod{ComiteFinalLoop}};
\draw[fluxo] (n43.south) -- (n44.north);
\node[pportao, below=of n44] (n45) {{\sffamily\bfseries\small 039\hspace{0.5em}Conta da virada no one-pager (determinístico)}\hfill{\sffamily\tiny\color{white}\colorbox{unAzul!80!unNavy}{PORTÃO}}\\[1pt]{\tiny\ttfamily\color{unLavanda}A039ContaDaViradaOnePager}\\[2pt]\campo{Faz}Refaz o ranking da tabela 1.3 sobre o one-pager de agora: os revisores podem ter mudado a mecânica depois do núcleo.\\[1pt]\campo{Objetivo}Refaz a conta da virada a partir da tabela 1.3 do núcleo (tools/mecanica.py).\\[1pt]\campo{Entra}\cod{one\_pager}\,\arq{produto/00-one-pager.md}\\[1pt]\campo{Sai}\cod{mecanica\_one\_pager\_report}\,\arq{produto/03c-conta-da-virada-one-pager.md}, \cod{mecanica\_one\_pager\_status}\\[1pt]\campo{Ciclo}\cod{ComiteFinalLoop}};
\draw[fluxo] (n44.south) -- (n45.north);
\node[pmodelo, below=of n45] (n46) {{\sffamily\bfseries\small 040a\hspace{0.5em}Comitê final: F·1 Pedagogia}\hfill{\sffamily\tiny\color{white}\colorbox{unLavanda}{MODELO}}\\[1pt]{\tiny\ttfamily\color{unLavanda}A040aComiteFinalPedagogia}\\[2pt]\campo{Faz}Comitê final: limiar vivido, aha com os seis passos e canais ainda coerentes depois das revisões.\\[1pt]\campo{Objetivo}O comitê final julga a missão completa — o one-pager depois do gate estrutural, das avaliações, dos revisores e da conferência do texto do PO —, antes do gate final e do parecer humano.\\[1pt]\campo{Entra}\cod{brief\_escopo}\,\arq{produto/00f-brief-escopo.md}, \cod{briefing}, \cod{comite\_final\_report}\,\arq{produto/03d-comite-final.md}, \cod{evidencia\_one\_pager\_report}\,\arq{produto/00k-contrato-evidencia-12a.md}, \cod{mapa\_conceitual}\,\arq{produto/00g-mapa-conceitual.md}, \cod{mecanica\_one\_pager\_report}\,\arq{produto/03c-conta-da-virada-one-pager.md}, \cod{one\_pager}\,\arq{produto/00-one-pager.md}\\[1pt]\campo{Sai}\cod{comite\_final\_pedagogia}\\[1pt]\campo{Ciclo}\cod{ComiteFinalLoop}\\[1pt]\campo{Roda}em paralelo com os demais de \cod{A040ComiteFinal}};
\draw[fluxo] (n45.south) -- (n46.north);
\node[pmodelo, below=of n46] (n47) {{\sffamily\bfseries\small 040b\hspace{0.5em}Comitê final: F·2 Qualidade de produto}\hfill{\sffamily\tiny\color{white}\colorbox{unLavanda}{MODELO}}\\[1pt]{\tiny\ttfamily\color{unLavanda}A040bComiteFinalQualidade}\\[2pt]\campo{Faz}Comitê final: uma missão só, tempos e tetos que fecham, professor claro e construível sem inventar.\\[1pt]\campo{Objetivo}O comitê final julga a missão completa, antes do gate final e do parecer humano.\\[1pt]\campo{Entra}\cod{brief\_escopo}\,\arq{produto/00f-brief-escopo.md}, \cod{briefing}, \cod{comite\_final\_report}\,\arq{produto/03d-comite-final.md}, \cod{evidencia\_one\_pager\_report}\,\arq{produto/00k-contrato-evidencia-12a.md}, \cod{mecanica\_one\_pager\_report}\,\arq{produto/03c-conta-da-virada-one-pager.md}, \cod{one\_pager}\,\arq{produto/00-one-pager.md}\\[1pt]\campo{Sai}\cod{comite\_final\_qualidade}\\[1pt]\campo{Ciclo}\cod{ComiteFinalLoop}\\[1pt]\campo{Roda}em paralelo com os demais de \cod{A040ComiteFinal}};
\draw[fluxo] (n46.south) -- (n47.north);
\node[pmodelo, below=of n47] (n48) {{\sffamily\bfseries\small 040c\hspace{0.5em}Comitê final: F·3 Imersão}\hfill{\sffamily\tiny\color{white}\colorbox{unLavanda}{MODELO}}\\[1pt]{\tiny\ttfamily\color{unLavanda}A040cComiteFinalImersao}\\[2pt]\campo{Faz}Comitê final: porta, troca e chegada diegéticas, verbos nativos e virada mostrada pelo ambiente.\\[1pt]\campo{Objetivo}O comitê final julga a missão completa, antes do gate final e do parecer humano.\\[1pt]\campo{Entra}\cod{brief\_escopo}\,\arq{produto/00f-brief-escopo.md}, \cod{briefing}, \cod{comite\_final\_report}\,\arq{produto/03d-comite-final.md}, \cod{one\_pager}\,\arq{produto/00-one-pager.md}\\[1pt]\campo{Sai}\cod{comite\_final\_imersao}\\[1pt]\campo{Ciclo}\cod{ComiteFinalLoop}\\[1pt]\campo{Roda}em paralelo com os demais de \cod{A040ComiteFinal}};
\draw[fluxo] (n47.south) -- (n48.north);
\draw[loopfaixa] ([xshift=-0.25cm]n43.north west) -- ([xshift=-0.25cm]n48.south west);
\draw[loopfaixa] ([xshift=-0.55cm]n43.north west) -- ([xshift=-0.55cm]n43.south west);
\end{tikzpicture}
\end{center}
\newpage
\begin{center}
\begin{tikzpicture}[node distance=0.42cm]
\node[pportao] (n49) {{\sffamily\bfseries\small 041\hspace{0.5em}Votos do comitê final (determinístico)}\hfill{\sffamily\tiny\color{white}\colorbox{unAzul!80!unNavy}{PORTÃO}}\\[1pt]{\tiny\ttfamily\color{unLavanda}A041VotosFinal}\\[2pt]\campo{Faz}Conta os votos de pedagogia, qualidade e imersão com as citações conferidas: unanimidade aprova; sem ela, a missão volta ao A004 e ao A014.\\[1pt]\campo{Objetivo}Conta os votos de um painel do comitê pedagógico (tools/comite.py).\\[1pt]\campo{Entra}\cod{nucleo\_narrativo}\,\arq{produto/00a-nucleo-narrativo.md}, \cod{one\_pager}\,\arq{produto/00-one-pager.md}, \cod{brief\_escopo}\,\arq{produto/00f-brief-escopo.md}, \cod{comite\_final\_pedagogia}, \cod{comite\_final\_qualidade}, \cod{comite\_final\_imersao}\\[1pt]\campo{Sai}\cod{comite\_final\_report}\,\arq{produto/03d-comite-final.md}, \cod{comite\_final\_status}, \textit{+1 de controle}\\[1pt]\campo{Ciclo}\cod{ComiteFinalLoop}};
\node[psaida, below=of n49] (n50) {{\sffamily\bfseries\small 042\hspace{0.5em}Saída do comitê final (determinístico)}\hfill{\sffamily\tiny\color{white}\colorbox{sinalAmbar}{SAÍDA DE LOOP}}\\[1pt]{\tiny\ttfamily\color{unLavanda}A042ComiteFinalRodada}\\[2pt]\campo{Faz}Os três de acordo, o \S{}12a em forma e a conta da virada fechando encerram o loop; senão, a cadeia de produto roda de novo a partir do núcleo, até COMITE\_FINAL\_MAX\_RODADAS (padrão 2). Esgotado, segue .\\[1pt]\campo{Objetivo}Encerra um LoopAgent quando todos os gates da rodada passaram — ou quando o teto caiu.\\[1pt]\campo{Entra}\cod{comite\_final\_status}, \cod{evidencia\_one\_pager\_status}, \cod{mecanica\_one\_pager\_status}, \cod{one\_pager}\,\arq{produto/00-one-pager.md}\\[1pt]\campo{Sai}\textit{+1 de controle}\\[1pt]\campo{Ciclo}\cod{ComiteFinalLoop}\\[1pt]\ciclo{$\circlearrowleft$ fim do ciclo: \cod{ComiteFinalLoop} --- reprovado, volta ao passo 004}};
\draw[fluxo] (n49.south) -- (n50.north);
\node[pmodelo, below=of n50] (n51) {{\sffamily\bfseries\small 043\hspace{0.5em}Gate final de produto}\hfill{\sffamily\tiny\color{white}\colorbox{unLavanda}{MODELO}}\\[1pt]{\tiny\ttfamily\color{unLavanda}A043GateFinalProduto}\\[2pt]\campo{Faz}Coerência entre missão e vereditos; libera ou bloqueia para design/dev.\\[1pt]\campo{Objetivo}Você responde uma única pergunta: a versão final da missão pode ser entregue para design/dev com segurança, sem contradizer a rubrica nem esconder pendências?\\[1pt]\campo{Entra}\cod{adequacao\_report}\,\arq{produto/03b-adequacao-pedagogica.md}, \cod{adequacao\_status}, \cod{avaliacao\_pedagogica\_report}\,\arq{produto/02-avaliacao-pedagogica.md}, \cod{briefing}, \cod{experiencia\_aluno\_report}\,\arq{produto/03-experiencia-aluno.md}, \cod{gate\_estrutural\_report}\,\arq{produto/01-gate-estrutural.md}, \cod{gate\_estrutural\_status}, \cod{one\_pager}\,\arq{produto/00-one-pager.md}, \cod{one\_pager\_po}, \textit{+1 de controle}\\[1pt]\campo{Sai}\cod{gate\_final\_report}\,\arq{produto/04-gate-final.md}};
\draw[fluxo] (n50.south) -- (n51.north);
\node[pportao, below=of n51] (n52) {{\sffamily\bfseries\small 044\hspace{0.5em}Veredito final (determinístico)}\hfill{\sffamily\tiny\color{white}\colorbox{unAzul!80!unNavy}{PORTÃO}}\\[1pt]{\tiny\ttfamily\color{unLavanda}A044GateFinalProdutoCheck}\\[2pt]\campo{Faz}Lê o envelope do gate final.\\[1pt]\campo{Objetivo}Lê o veredito de um relatório e decide se o pipeline segue, repete ou para.\\[1pt]\campo{Entra}\cod{one\_pager}\,\arq{produto/00-one-pager.md}, \cod{brief\_escopo}\,\arq{produto/00f-brief-escopo.md}, \cod{gate\_estrutural\_report}\,\arq{produto/01-gate-estrutural.md}, \cod{avaliacao\_pedagogica\_report}\,\arq{produto/02-avaliacao-pedagogica.md}, \cod{experiencia\_aluno\_report}\,\arq{produto/03-experiencia-aluno.md}, \cod{adequacao\_report}\,\arq{produto/03b-adequacao-pedagogica.md}, \cod{gate\_final\_report}\,\arq{produto/04-gate-final.md}\\[1pt]\campo{Sai}\cod{gate\_final\_status}};
\draw[fluxo] (n51.south) -- (n52.north);
\node[pcodigo, below=of n52] (n53) {{\sffamily\bfseries\small 044b\hspace{0.5em}Textos-chave com ID (determinístico)}\hfill{\sffamily\tiny\color{white}\colorbox{unNavy}{CÓDIGO}}\\[1pt]{\tiny\ttfamily\color{unLavanda}A044bTextosChave}\\[2pt]\campo{Faz}Os textos-chave do one-pager aprovado — aspas do \S{}4, pergunta e alternativas do Canal B (\S{}12a), opções de cada escolha com o nº de opções, roteiro do professor (\S{}9), podcast (\S{}11) — numa tabela com ID estável (`TC-07`, a opção `TC-12.2`), montada por código. O cartão do de acordo a mostra; design e dev a recebem no conteúdo aprovado, e o código copia cada texto até a copy (PR 7).\\[1pt]\campo{Objetivo}Determinístico: a tabela de textos-chave com ID, gravada no fim de produto (tools/textoschave.py, PR 7).\\[1pt]\campo{Entra}\cod{one\_pager}\,\arq{produto/00-one-pager.md}, \cod{briefing}, \cod{decisoes\_do\_po}\\[1pt]\campo{Sai}\cod{textos\_chave\_tabela}\,\arq{produto/00n-textos-chave.md}};
\draw[fluxo] (n52.south) -- (n53.north);
\draw[loopfaixa] ([xshift=-0.25cm]n49.north west) -- ([xshift=-0.25cm]n50.south west);
\end{tikzpicture}
\end{center}
\newpage
\begin{center}
\begin{tikzpicture}[node distance=0.42cm]
\node[pmodelo] (n54) {{\sffamily\bfseries\small 045\hspace{0.5em}Spec da missão}\hfill{\sffamily\tiny\color{white}\colorbox{unLavanda}{MODELO}}\\[1pt]{\tiny\ttfamily\color{unLavanda}A045SpecMissao}\\[2pt]\ciclo{$\circlearrowright$ começa aqui: \cod{SpecLoop} (até 2 rodadas)}\\[1pt]\campo{Faz}Os 12 blocos da spec — contrato contra o qual design e dev constroem.\\[1pt]\campo{Objetivo}Você traduz o one-pager aprovado na especificação de missão — o contrato contra o qual design e dev constroem.\\[1pt]\campo{Entra}\cod{adequacao\_report}\,\arq{produto/03b-adequacao-pedagogica.md}, \cod{adequacao\_status}, \cod{avaliacao\_pedagogica\_report}\,\arq{produto/02-avaliacao-pedagogica.md}, \cod{banca\_produto\_report}\,\arq{produto/05i-banca-produto.md}, \cod{bncc\_report}\,\arq{produto/05b-conformidade-bncc.md}, \cod{briefing}, \cod{conteudo\_aprovado\_spec\_report}\,\arq{produto/05g-conteudo-aprovado.md}, \cod{evidencia\_report}\,\arq{produto/05c-matriz-evidencia.md}, \cod{experiencia\_aluno\_report}\,\arq{produto/03-experiencia-aluno.md}, +11\\[1pt]\campo{Sai}\cod{spec\_missao}\,\arq{produto/05-spec-missao.md}\\[1pt]\campo{Ciclo}\cod{SpecLoop}};
\node[pmodelo, below=of n54] (n55) {{\sffamily\bfseries\small 045b\hspace{0.5em}Completude da spec}\hfill{\sffamily\tiny\color{white}\colorbox{unLavanda}{MODELO}}\\[1pt]{\tiny\ttfamily\color{unLavanda}A045bCompletudeSpec}\\[2pt]\ciclo{$\circlearrowright$ começa aqui: \cod{CompletudeLoop} (até 2 rodadas)}\\[1pt]\campo{Faz}O que o D0 fechava no design e é de produto (estado, handoff, evidência, dados e retenção, duração e tetos, painel docente, PQ e KPI, texto do professor, revezamento), em produto: decisões [CP-n] com base declarada no bloco 15 da spec, padrões da plataforma aplicados e listados no de acordo, e o que é do PO como [PARÂMETRO]/[PROPOSTA]/[ABERTO — dono: PO] na lista do de acordo. Nas rodadas seguintes do loop só roda se a spec perdeu o bloco 15; na correção pontual não roda.\\[1pt]\campo{Objetivo}Roda logo depois do especificador (A045), dentro do loop da spec, antes dos gates da spec, da banca de produto e do de acordo do PO.\\[1pt]\campo{Entra}\cod{catalogo\_missoes}, \cod{completude\_report}, \cod{gate\_estrutural\_report}\,\arq{produto/01-gate-estrutural.md}, \cod{one\_pager}\,\arq{produto/00-one-pager.md}, \cod{spec\_missao}\,\arq{produto/05-spec-missao.md}\\[1pt]\campo{Sai}\cod{completude\_spec}\,\arq{produto/05h-completude-spec.md}\\[1pt]\campo{Ciclo}\cod{SpecLoop} $\supset$ \cod{CompletudeLoop}};
\draw[fluxo] (n54.south) -- (n55.north);
\node[pcodigo, below=of n55] (n56) {{\sffamily\bfseries\small 045c\hspace{0.5em}Spec com o bloco 15 (determinístico)}\hfill{\sffamily\tiny\color{white}\colorbox{unNavy}{CÓDIGO}}\\[1pt]{\tiny\ttfamily\color{unLavanda}A045cCompletudeSpecSplit}\\[2pt]\campo{Faz}Monta em código a spec de produto — a do A045 com o bloco 15 da completude no lugar, o resto descartado —, separa o relatório de completude e registra `completude\_status`. [PARÂMETRO]/[PROPOSTA] fechado pela completude sem parecer datado do PO bloqueia.\\[1pt]\campo{Objetivo}Copia seções de um relatório para chaves próprias do estado e registra um status informativo.\\[1pt]\campo{Entra}\cod{completude\_spec}\,\arq{produto/05h-completude-spec.md}, \cod{completude\_report}, \cod{notas\_de\_passagem}, \cod{pedidos\_do\_po}, \cod{spec\_missao\_gate0}\,\arq{produto/05a-spec-missao-gate0.md}, \cod{spec\_missao}\,\arq{produto/05-spec-missao.md}\\[1pt]\campo{Sai}\cod{relatorio\_completude}\,\arq{produto/05i-relatorio-completude.md}, \cod{completude\_status}\\[1pt]\campo{Ciclo}\cod{SpecLoop} $\supset$ \cod{CompletudeLoop}};
\draw[fluxo] (n55.south) -- (n56.north);
\node[psaida, below=of n56] (n57) {{\sffamily\bfseries\small 045d\hspace{0.5em}Saída da completude (determinístico)}\hfill{\sffamily\tiny\color{white}\colorbox{sinalAmbar}{SAÍDA DE LOOP}}\\[1pt]{\tiny\ttfamily\color{unLavanda}A045dCompletudeRodada}\\[2pt]\campo{Faz}Completude conferida encerra o loop dela; reprovada pelo código (decidiu no lugar do PO, ou veio sem o bloco 15) volta só a ela, uma vez, com os achados — o autor da spec não reescreve por isso.\\[1pt]\campo{Objetivo}Encerra um LoopAgent quando todos os gates da rodada passaram — ou quando o teto caiu.\\[1pt]\campo{Entra}\cod{completude\_status}\\[1pt]\campo{Sai}\textit{+1 de controle}\\[1pt]\campo{Ciclo}\cod{SpecLoop} $\supset$ \cod{CompletudeLoop}\\[1pt]\ciclo{$\circlearrowleft$ fim do ciclo: \cod{CompletudeLoop} --- reprovado, volta ao passo 045b}};
\draw[fluxo] (n56.south) -- (n57.north);
\node[pportao, below=of n57] (n58) {{\sffamily\bfseries\small 046\hspace{0.5em}Conformidade BNCC Computação (determinístico)}\hfill{\sffamily\tiny\color{white}\colorbox{unAzul!80!unNavy}{PORTÃO}}\\[1pt]{\tiny\ttfamily\color{unLavanda}A046BnccCheck}\\[2pt]\campo{Faz}Bloco 11 $\times$ lista oficial do MEC: código existe, é do ano da missão e inclui letramento digital. Não conforme bloqueia design e dev.\\[1pt]\campo{Objetivo}Gate determinístico de conformidade com a BNCC Computação, aplicado à spec.\\[1pt]\campo{Entra}\cod{spec\_missao}\,\arq{produto/05-spec-missao.md}, \cod{material\_anexado}\,\arq{produto/00-material-anexado.md}\\[1pt]\campo{Sai}\cod{bncc\_report}\,\arq{produto/05b-conformidade-bncc.md}, \cod{bncc\_status}\\[1pt]\campo{Ciclo}\cod{SpecLoop}};
\draw[fluxo] (n57.south) -- (n58.north);
\draw[loopfaixa] ([xshift=-0.25cm]n54.north west) -- ([xshift=-0.25cm]n58.south west);
\draw[loopfaixa] ([xshift=-0.55cm]n55.north west) -- ([xshift=-0.55cm]n57.south west);
\draw[loopvolta] ([xshift=-0.55cm]n57.south west) -- ++(-0.18,0) |- ([xshift=-0.55cm]n55.north west);
\end{tikzpicture}
\end{center}
\newpage
\begin{center}
\begin{tikzpicture}[node distance=0.42cm]
\node[pportao] (n59) {{\sffamily\bfseries\small 047\hspace{0.5em}Ficha do manifesto (determinístico)}\hfill{\sffamily\tiny\color{white}\colorbox{unAzul!80!unNavy}{PORTÃO}}\\[1pt]{\tiny\ttfamily\color{unLavanda}A047ManifestoCheckProduto}\\[2pt]\campo{Faz}Bloco 0 da spec $\times$ listas fechadas da plataforma, conferido ainda em Produto: campo errado volta para quem escreveu a spec, em vez de custar o design inteiro.\\[1pt]\campo{Objetivo}Gate determinístico de rastro: modo decide a verificação de tools/rastro.py.\\[1pt]\campo{Entra}\cod{briefing}, \cod{one\_pager\_po}\\[1pt]\campo{Sai}\cod{manifesto\_report}\,\arq{dev/00-ficha-manifesto.md}, \cod{manifesto\_status}\\[1pt]\campo{Ciclo}\cod{SpecLoop}};
\node[pportao, below=of n59] (n60) {{\sffamily\bfseries\small 048\hspace{0.5em}Matriz de evidência pedagógica (determinístico)}\hfill{\sffamily\tiny\color{white}\colorbox{unAzul!80!unNavy}{PORTÃO}}\\[1pt]{\tiny\ttfamily\color{unLavanda}A048EvidenciaCheck}\\[2pt]\campo{Faz}Bloco 9 $\times$ bloco 11: grau `scored` com dois sinais nomeados, toda habilidade com momento em tela e medida ligada, gabarito declarado no servidor.\\[1pt]\campo{Objetivo}Gate determinístico de rastro: modo decide a verificação de tools/rastro.py.\\[1pt]\campo{Entra}\textit{nada do estado da sessão (configuração fixa do passo)}\\[1pt]\campo{Sai}\cod{evidencia\_report}\,\arq{produto/05c-matriz-evidencia.md}, \cod{evidencia\_status}\\[1pt]\campo{Ciclo}\cod{SpecLoop}};
\draw[fluxo] (n59.south) -- (n60.north);
\node[pportao, below=of n60] (n61) {{\sffamily\bfseries\small 049\hspace{0.5em}Juiz da matriz de evidência}\hfill{\sffamily\tiny\color{white}\colorbox{unAzul!80!unNavy}{PORTÃO}}\\[1pt]{\tiny\ttfamily\color{unLavanda}A049JuizEvidencia}\\[2pt]\campo{Faz}Juiz de achados sobre `Matriz de evidência pedagógica`: retira o achado de uso vedado quando o texto só cita para vedar — com citação literal conferida.\\[1pt]\campo{Objetivo}O regex acha, o LLM decide, o código confere a citação (tools/juizachados.py).\\[1pt]\campo{Entra}\cod{spec\_missao}\,\arq{produto/05-spec-missao.md}, \textit{+1 de controle}\\[1pt]\campo{Sai}\cod{evidencia\_status}, \cod{evidencia\_report}\,\arq{produto/05c-matriz-evidencia.md}\\[1pt]\campo{Ciclo}\cod{SpecLoop}};
\draw[fluxo] (n60.south) -- (n61.north);
\node[pportao, below=of n61] (n62) {{\sffamily\bfseries\small 050\hspace{0.5em}Perguntas de produto, sub-bloco 9f (determinístico)}\hfill{\sffamily\tiny\color{white}\colorbox{unAzul!80!unNavy}{PORTÃO}}\\[1pt]{\tiny\ttfamily\color{unLavanda}A050PerguntasCheck}\\[2pt]\campo{Faz}As cinco perguntas de produto: quantidade fechada, decisão com limiar, corte do vocabulário, e a ação que responde cada uma declarada na tabela de eventos do bloco 9 com payload que mede.\\[1pt]\campo{Objetivo}Gate determinístico de rastro: modo decide a verificação de tools/rastro.py.\\[1pt]\campo{Entra}\textit{nada do estado da sessão (configuração fixa do passo)}\\[1pt]\campo{Sai}\cod{perguntas\_report}\,\arq{produto/05d-perguntas-de-produto.md}, \cod{perguntas\_status}\\[1pt]\campo{Ciclo}\cod{SpecLoop}};
\draw[fluxo] (n61.south) -- (n62.north);
\node[pmodelo, below=of n62] (n63) {{\sffamily\bfseries\small 051\hspace{0.5em}Juiz das perguntas de produto}\hfill{\sffamily\tiny\color{white}\colorbox{unLavanda}{MODELO}}\\[1pt]{\tiny\ttfamily\color{unLavanda}A051JuizPerguntas}\\[2pt]\campo{Faz}O que não se confere por regex: a pergunta muda uma decisão nossa, é sobre o produto e não sobre a criança, a ação nomeada a responde, e ela não é uma transversal repintada.\\[1pt]\campo{Objetivo}Um gate determinístico já rodou antes de você e conferiu o que é fato: são cinco perguntas, cada uma tem decisão com limiar, corte do vocabulário, e uma ação verbo objeto:slug que existe na tabela de eventos do bloco 9 e carrega…\\[1pt]\campo{Entra}\cod{perguntas\_report}\,\arq{produto/05d-perguntas-de-produto.md}, \cod{perguntas\_status}, \cod{spec\_missao}\,\arq{produto/05-spec-missao.md}\\[1pt]\campo{Sai}\cod{juiz\_perguntas\_report}\,\arq{produto/05e-juiz-perguntas-de-produto.md}\\[1pt]\campo{Ciclo}\cod{SpecLoop}};
\draw[fluxo] (n62.south) -- (n63.north);
\node[pportao, below=of n63] (n64) {{\sffamily\bfseries\small 052\hspace{0.5em}Veredito do juiz das perguntas (determinístico)}\hfill{\sffamily\tiny\color{white}\colorbox{unAzul!80!unNavy}{PORTÃO}}\\[1pt]{\tiny\ttfamily\color{unLavanda}A052JuizPerguntasCheck}\\[2pt]\campo{Faz}Lê o envelope do juiz:  em pergunta declarada pelo PO bloqueia;  só em proposta segue como ressalva para o PO decidir.\\[1pt]\campo{Objetivo}Lê o veredito de um relatório e decide se o pipeline segue, repete ou para.\\[1pt]\campo{Entra}\cod{spec\_missao}\,\arq{produto/05-spec-missao.md}, \cod{perguntas\_po}, \cod{juiz\_perguntas\_report}\,\arq{produto/05e-juiz-perguntas-de-produto.md}\\[1pt]\campo{Sai}\cod{juiz\_perguntas\_status}, \textit{+1 de controle}\\[1pt]\campo{Ciclo}\cod{SpecLoop}};
\draw[fluxo] (n63.south) -- (n64.north);
\draw[loopfaixa] ([xshift=-0.25cm]n59.north west) -- ([xshift=-0.25cm]n64.south west);
\end{tikzpicture}
\end{center}
\newpage
\begin{center}
\begin{tikzpicture}[node distance=0.42cm]
\node[pportao] (n65) {{\sffamily\bfseries\small 053\hspace{0.5em}Rastro de lacunas da spec (determinístico)}\hfill{\sffamily\tiny\color{white}\colorbox{unAzul!80!unNavy}{PORTÃO}}\\[1pt]{\tiny\ttfamily\color{unLavanda}A053RastroSpec}\\[2pt]\campo{Faz}Marcadores de lacuna na spec, por dono: o de produto ou de insumo volta para quem escreveu a spec; a decisão do PO fica  no semáforo do parecer; o de design e o de dev ficam para a etapa delas.\\[1pt]\campo{Objetivo}Gate determinístico de rastro: modo decide a verificação de tools/rastro.py.\\[1pt]\campo{Entra}\cod{spec\_missao}\,\arq{produto/05-spec-missao.md}\\[1pt]\campo{Sai}\cod{rastro\_spec\_report}\,\arq{produto/05f-rastro-lacunas-spec.md}, \cod{rastro\_spec\_status}\\[1pt]\campo{Ciclo}\cod{SpecLoop}};
\node[pportao, below=of n65] (n66) {{\sffamily\bfseries\small 054\hspace{0.5em}Spec $\times$ one-pager aprovado (determinístico)}\hfill{\sffamily\tiny\color{white}\colorbox{unAzul!80!unNavy}{PORTÃO}}\\[1pt]{\tiny\ttfamily\color{unLavanda}A054ConteudoAprovadoSpec}\\[2pt]\campo{Faz}A spec traduz o one-pager que o PO aprovou: mesmas habilidades do contrato de evidência (\S{}12a), mesmo tempo total em tela (\S{}10) e o podcast (\S{}11) declarado. Ressalva de parecer que o PO não incorporou não muda a spec.\\[1pt]\campo{Objetivo}Gate determinístico de rastro: modo decide a verificação de tools/rastro.py.\\[1pt]\campo{Entra}\cod{spec\_missao}\,\arq{produto/05-spec-missao.md}, \cod{one\_pager}\,\arq{produto/00-one-pager.md}, \cod{prd\_final}\,\arq{PRD.md}, \cod{one\_pager\_resumido}\,\arq{produto/00d-one-pager-resumido.md}, \cod{one\_pager\_pedagogico}\,\arq{one-pager-pedagogico.md}\\[1pt]\campo{Sai}\cod{conteudo\_aprovado\_spec\_report}\,\arq{produto/05g-conteudo-aprovado.md}, \cod{conteudo\_aprovado\_spec\_status}\\[1pt]\campo{Ciclo}\cod{SpecLoop}};
\draw[fluxo] (n65.south) -- (n66.north);
\node[pcodigo, below=of n66] (n67) {{\sffamily\bfseries\small 055\hspace{0.5em}Conteúdo aprovado para a banca de produto (determinístico)}\hfill{\sffamily\tiny\color{white}\colorbox{unNavy}{CÓDIGO}}\\[1pt]{\tiny\ttfamily\color{unLavanda}A055ConteudoAprovado}\\[2pt]\campo{Faz}Grava as seções literais do one-pager para a banca do fim de produto.\\[1pt]\campo{Objetivo}Grava em conteudoaprovado as seções literais do one-pager aprovado (tools/conteudoaprovado.py).\\[1pt]\campo{Entra}\cod{conteudo\_aprovado}, \cod{one\_pager}\,\arq{produto/00-one-pager.md}, \cod{briefing}\\[1pt]\campo{Sai}\textit{nada do estado da sessão (configuração fixa do passo)}\\[1pt]\campo{Ciclo}\cod{SpecLoop}};
\draw[fluxo] (n66.south) -- (n67.north);
\node[pcodigo, below=of n67] (n68) {{\sffamily\bfseries\small 056\hspace{0.5em}Escopo da banca de produto (determinístico)}\hfill{\sffamily\tiny\color{white}\colorbox{unNavy}{CÓDIGO}}\\[1pt]{\tiny\ttfamily\color{unLavanda}A056EscopoBanca}\\[2pt]\campo{Faz}Grava o escopo da banca na posição «produto».\\[1pt]\campo{Objetivo}Grava em escopobanca onde a banca está rodando e o que ela pode fazer ali (tools/banca.ESCOPOS).\\[1pt]\campo{Entra}\textit{nada do estado da sessão (configuração fixa do passo)}\\[1pt]\campo{Sai}\cod{escopo\_banca}\\[1pt]\campo{Ciclo}\cod{SpecLoop}};
\draw[fluxo] (n67.south) -- (n68.north);
\node[pmodelo, below=of n68] (n69) {{\sffamily\bfseries\small 057a\hspace{0.5em}Banca no fim de produto: produto e narrativa}\hfill{\sffamily\tiny\color{white}\colorbox{unLavanda}{MODELO}}\\[1pt]{\tiny\ttfamily\color{unLavanda}A057aBancaProdutoNarrativa}\\[2pt]\campo{Faz}Os critérios PROD da banca que a spec decide, antes do parecer do PO. O achado volta ao autor da spec na próxima rodada.\\[1pt]\campo{Objetivo}A mesma banca roda em quatro posições, com os mesmos critérios — fim de produto, fim de design, entrada do dev e dentro do dev.\\[1pt]\campo{Entra}\cod{anexo\_design\_spec}\,\arq{design/00-anexo-design-spec.md}, \cod{banco\_conteudo}\,\arq{dev/01e-banco-conteudo.md}, \cod{brief\_assets}\,\arq{dev/01b-brief-assets.md}, \cod{brief\_oficio}\,\arq{design/01-briefing-oficio.md}, \cod{camada\_caa}\,\arq{design/04b-camada-caa.md}, \cod{conteudo\_aprovado}, \cod{copy\_strings\_revisada}\,\arq{design/06-copy-strings-revisada.md}, \cod{escopo\_banca}, \cod{fluxo\_arestas}\,\arq{design/03-fluxo-arestas.md}, +6\\[1pt]\campo{Sai}\cod{banca\_produto\_narrativa}\\[1pt]\campo{Ciclo}\cod{SpecLoop}\\[1pt]\campo{Roda}em paralelo com os demais de \cod{A057Banca}};
\draw[fluxo] (n68.south) -- (n69.north);
\node[pmodelo, below=of n69] (n70) {{\sffamily\bfseries\small 057b\hspace{0.5em}Banca no fim de produto: pedagogia e BNCC}\hfill{\sffamily\tiny\color{white}\colorbox{unLavanda}{MODELO}}\\[1pt]{\tiny\ttfamily\color{unLavanda}A057bBancaPedagogiaBncc}\\[2pt]\campo{Faz}Os critérios PED da banca sobre a spec, antes do parecer do PO.\\[1pt]\campo{Objetivo}A mesma banca roda em quatro posições, com os mesmos critérios — fim de produto, fim de design, entrada do dev e dentro do dev.\\[1pt]\campo{Entra}\cod{anexo\_design\_spec}\,\arq{design/00-anexo-design-spec.md}, \cod{banco\_conteudo}\,\arq{dev/01e-banco-conteudo.md}, \cod{brief\_assets}\,\arq{dev/01b-brief-assets.md}, \cod{brief\_oficio}\,\arq{design/01-briefing-oficio.md}, \cod{camada\_caa}\,\arq{design/04b-camada-caa.md}, \cod{conteudo\_aprovado}, \cod{copy\_strings\_revisada}\,\arq{design/06-copy-strings-revisada.md}, \cod{escopo\_banca}, \cod{fluxo\_arestas}\,\arq{design/03-fluxo-arestas.md}, +6\\[1pt]\campo{Sai}\cod{banca\_produto\_pedagogia}\\[1pt]\campo{Ciclo}\cod{SpecLoop}\\[1pt]\campo{Roda}em paralelo com os demais de \cod{A057Banca}};
\draw[fluxo] (n69.south) -- (n70.north);
\draw[loopfaixa] ([xshift=-0.25cm]n65.north west) -- ([xshift=-0.25cm]n70.south west);
\end{tikzpicture}
\end{center}
\newpage
\begin{center}
\begin{tikzpicture}[node distance=0.42cm]
\node[pportao] (n71) {{\sffamily\bfseries\small 058\hspace{0.5em}Banca de produto (determinístico)}\hfill{\sffamily\tiny\color{white}\colorbox{unAzul!80!unNavy}{PORTÃO}}\\[1pt]{\tiny\ttfamily\color{unLavanda}A058BancaConsolidacao}\\[2pt]\campo{Faz}Veredito da banca no fim de produto: o pior das dimensões, com os bloqueantes lidos dos arquivos de critério.  devolve a spec ao autor e, esgotadas as rodadas, recusa o de acordo.\\[1pt]\campo{Objetivo}Determinístico: o veredito da banca é a pior das cinco dimensões, contada em código.\\[1pt]\campo{Entra}\cod{banca\_produto\_narrativa}, \cod{banca\_produto\_pedagogia}, \cod{spec\_missao}\,\arq{produto/05-spec-missao.md}, \cod{one\_pager}\,\arq{produto/00-one-pager.md}\\[1pt]\campo{Sai}\cod{banca\_produto\_report}\,\arq{produto/05i-banca-produto.md}, \cod{banca\_produto\_status}\\[1pt]\campo{Ciclo}\cod{SpecLoop}};
\node[psaida, below=of n71] (n72) {{\sffamily\bfseries\small 059\hspace{0.5em}Veredito da spec (determinístico)}\hfill{\sffamily\tiny\color{white}\colorbox{sinalAmbar}{SAÍDA DE LOOP}}\\[1pt]{\tiny\ttfamily\color{unLavanda}A059SpecCheck}\\[2pt]\campo{Faz}Currículo e manifesto conformes encerram o loop da spec; reprovados devolvem ao autor com os dois relatórios.\\[1pt]\campo{Objetivo}Encerra um LoopAgent quando todos os gates da rodada passaram — ou quando o teto caiu.\\[1pt]\campo{Entra}\cod{bncc\_status}, \cod{manifesto\_status}, \cod{evidencia\_status}, \cod{perguntas\_status}, \cod{juiz\_perguntas\_status}, \cod{rastro\_spec\_status}, \cod{conteudo\_aprovado\_spec\_status}, \cod{banca\_produto\_status}, \cod{spec\_missao}\,\arq{produto/05-spec-missao.md}\\[1pt]\campo{Sai}\textit{+1 de controle}\\[1pt]\campo{Ciclo}\cod{SpecLoop}\\[1pt]\ciclo{$\circlearrowleft$ fim do ciclo: \cod{SpecLoop} --- reprovado, volta ao passo 045}};
\draw[fluxo] (n71.south) -- (n72.north);
\node[pportao, below=of n72] (n73) {{\sffamily\bfseries\small 060\hspace{0.5em}Fidelidade ao pedido do PO e ao briefing}\hfill{\sffamily\tiny\color{white}\colorbox{unAzul!80!unNavy}{PORTÃO}}\\[1pt]{\tiny\ttfamily\color{unLavanda}A060FidelidadeProduto}\\[2pt]\campo{Faz}Núcleo, one-pager e spec contra todos os pedidos do PO e o briefing. Pedido sem prova é  e recusa o de acordo de produto; conflito com regra da esteira vira decisão do PO.\\[1pt]\campo{Objetivo}O que o PO pediu e o briefing fixou está nos documentos?\\[1pt]\campo{Entra}\cod{nucleo\_narrativo}\,\arq{produto/00a-nucleo-narrativo.md}, \cod{one\_pager}\,\arq{produto/00-one-pager.md}, \cod{spec\_missao}\,\arq{produto/05-spec-missao.md}, \cod{divisao\_do\_pedido}, \cod{pedidos\_do\_po}, \cod{pedido\_correcao}, \cod{etapa\_da\_correcao}, \cod{briefing}, \cod{edicoes\_pontuais}, \textit{+3 de controle}\\[1pt]\campo{Sai}\cod{fidelidade\_report}, \cod{fidelidade\_status}};
\draw[fluxo] (n72.south) -- (n73.north);
\node[pmodelo, below=of n73] (n74) {{\sffamily\bfseries\small 061\hspace{0.5em}One-pager resumido}\hfill{\sffamily\tiny\color{white}\colorbox{unLavanda}{MODELO}}\\[1pt]{\tiny\ttfamily\color{unLavanda}A061OnePagerResumido}\\[2pt]\ciclo{$\circlearrowright$ começa aqui: \cod{OnePagerResumidoLoop} (até 2 rodadas)}\\[1pt]\campo{Faz}O one-pager no formato do time de produto para avaliação pedagógica: ficha, ideia, componente curricular, socioemocionais, 6 etapas, professor, contrato de evidência ternário e takeaway. Sai em .md, .docx e .html.\\[1pt]\campo{Objetivo}O one-pager canônico de 18 seções já foi escrito, auditado pelo gate estrutural, avaliado pelos três pareceres e traduzido em spec com o currículo conferido contra a lista oficial da BNCC.\\[1pt]\campo{Entra}\cod{adequacao\_report}\,\arq{produto/03b-adequacao-pedagogica.md}, \cod{avaliacao\_pedagogica\_report}\,\arq{produto/02-avaliacao-pedagogica.md}, \cod{bncc\_report}\,\arq{produto/05b-conformidade-bncc.md}, \cod{briefing}, \cod{catalogo\_missoes}, \cod{conteudo\_aprovado\_resumido\_report}\,\arq{produto/05h-conteudo-aprovado-resumido.md}, \cod{experiencia\_aluno\_report}\,\arq{produto/03-experiencia-aluno.md}, \cod{gate\_final\_report}\,\arq{produto/04-gate-final.md}, \cod{mapa\_conceitual}\,\arq{produto/00g-mapa-conceitual.md}, +3\\[1pt]\campo{Sai}\cod{one\_pager\_resumido}\,\arq{produto/00d-one-pager-resumido.md}\\[1pt]\campo{Ciclo}\cod{OnePagerResumidoLoop}};
\draw[fluxo] (n73.south) -- (n74.north);
\node[pportao, below=of n74] (n75) {{\sffamily\bfseries\small 062\hspace{0.5em}Resumido $\times$ one-pager aprovado (determinístico)}\hfill{\sffamily\tiny\color{white}\colorbox{unAzul!80!unNavy}{PORTÃO}}\\[1pt]{\tiny\ttfamily\color{unLavanda}A062ConteudoAprovadoResumido}\\[2pt]\campo{Faz}As habilidades do bloco 11, o tempo em tela do \S{}10 e o Canal B com o texto do \S{}12a no resumido que o time lê para aprovar. Divergência volta ao autor do resumido.\\[1pt]\campo{Objetivo}Gate determinístico de rastro: modo decide a verificação de tools/rastro.py.\\[1pt]\campo{Entra}\cod{spec\_missao}\,\arq{produto/05-spec-missao.md}, \cod{one\_pager}\,\arq{produto/00-one-pager.md}, \cod{prd\_final}\,\arq{PRD.md}, \cod{one\_pager\_resumido}\,\arq{produto/00d-one-pager-resumido.md}, \cod{one\_pager\_pedagogico}\,\arq{one-pager-pedagogico.md}\\[1pt]\campo{Sai}\cod{conteudo\_aprovado\_resumido\_report}\,\arq{produto/05h-conteudo-aprovado-resumido.md}, \cod{conteudo\_aprovado\_resumido\_status}\\[1pt]\campo{Ciclo}\cod{OnePagerResumidoLoop}};
\draw[fluxo] (n74.south) -- (n75.north);
\draw[loopfaixa] ([xshift=-0.25cm]n71.north west) -- ([xshift=-0.25cm]n72.south west);
\draw[loopfaixa] ([xshift=-0.25cm]n74.north west) -- ([xshift=-0.25cm]n75.south west);
\end{tikzpicture}
\end{center}
\newpage
\begin{center}
\begin{tikzpicture}[node distance=0.42cm]
\node[psaida] (n76) {{\sffamily\bfseries\small 063\hspace{0.5em}Veredito do resumido (determinístico)}\hfill{\sffamily\tiny\color{white}\colorbox{sinalAmbar}{SAÍDA DE LOOP}}\\[1pt]{\tiny\ttfamily\color{unLavanda}A063ResumidoRodada}\\[2pt]\campo{Faz}Resumido fiel encerra o loop; divergente volta ao autor com o relatório, até o teto de rodadas.\\[1pt]\campo{Objetivo}Encerra um LoopAgent quando todos os gates da rodada passaram — ou quando o teto caiu.\\[1pt]\campo{Entra}\cod{conteudo\_aprovado\_resumido\_status}, \cod{one\_pager\_resumido}\,\arq{produto/00d-one-pager-resumido.md}\\[1pt]\campo{Sai}\textit{+1 de controle}\\[1pt]\campo{Ciclo}\cod{OnePagerResumidoLoop}\\[1pt]\ciclo{$\circlearrowleft$ fim do ciclo: \cod{OnePagerResumidoLoop} --- reprovado, volta ao passo 061}};
\draw[loopfaixa] ([xshift=-0.25cm]n76.north west) -- ([xshift=-0.25cm]n76.south west);
\end{tikzpicture}
\end{center}
\newpage
\markboth{Design}{}
{\sffamily\bfseries\Large\color{unNavy}Design}\\[0.3em]
\begin{center}
\begin{tikzpicture}[node distance=0.42cm]
\node[pcodigo] (n77) {{\sffamily\bfseries\small 064\hspace{0.5em}Conteúdo aprovado (determinístico)}\hfill{\sffamily\tiny\color{white}\colorbox{unNavy}{CÓDIGO}}\\[1pt]{\tiny\ttfamily\color{unLavanda}A064ConteudoAprovadoDesign}\\[2pt]\campo{Faz}As seções do one-pager aprovado — o núcleo (descoberta, narrativa e as 6 etapas com o texto de tela), currículo, falhas, professor, tempos e cenários, podcast, contrato de evidência, dispositivos, hipóteses —, literais, para os agentes de design que não leem o one-pager.\\[1pt]\campo{Objetivo}Grava em conteudoaprovado as seções literais do one-pager aprovado (tools/conteudoaprovado.py).\\[1pt]\campo{Entra}\cod{one\_pager}\,\arq{produto/00-one-pager.md}, \cod{briefing}\\[1pt]\campo{Sai}\cod{conteudo\_aprovado}};
\node[pmodelo, below=of n77] (n78) {{\sffamily\bfseries\small 065\hspace{0.5em}Fechamento de lacunas de UI/UX (D0-UX)}\hfill{\sffamily\tiny\color{white}\colorbox{unLavanda}{MODELO}}\\[1pt]{\tiny\ttfamily\color{unLavanda}A065FechaLacunas}\\[2pt]\campo{Faz}Mandato só de UI/UX desde o PR 6: varredura do que a spec de produto deixa da tela de cada estado, gesto, scaffolding, acessibilidade de componente e dispositivo, decisões [D0-n] com base declarada, bloco 13 do anexo de design (com o mapa célula $\rightarrow$ D0) e relatório de preenchimento. O que é de produto vira `[ABERTO — dono: PO]`.\\[1pt]\campo{Objetivo}Roda depois do de acordo de produto e antes do briefing de ofício (D1).\\[1pt]\campo{Entra}\cod{anexo\_design\_spec}\,\arq{design/00-anexo-design-spec.md}, \cod{catalogo\_missoes}, \cod{gate\_estrutural\_report}\,\arq{produto/01-gate-estrutural.md}, \cod{one\_pager}\,\arq{produto/00-one-pager.md}, \cod{spec\_missao}\,\arq{produto/05-spec-missao.md}, \textit{+1 de controle}\\[1pt]\campo{Sai}\cod{fechamento\_lacunas}\,\arq{design/00-fechamento-lacunas.md}};
\draw[fluxo] (n77.south) -- (n78.north);
\node[pcodigo, below=of n78] (n79) {{\sffamily\bfseries\small 066\hspace{0.5em}Relatório de lacunas + anexo de design (determinístico)}\hfill{\sffamily\tiny\color{white}\colorbox{unNavy}{CÓDIGO}}\\[1pt]{\tiny\ttfamily\color{unLavanda}A066FechaLacunasSplit}\\[2pt]\campo{Faz}Monta em código o anexo de design (`anexo\_design\_spec`, `design/00-anexo-design-spec.md`) — bloco 13 do D0 + bloco 14 que já houver —, sem tocar a spec de produto; separa o relatório e registra `lacunas\_status`. Mudança no texto de produto vira `[ABERTO — dono: PO]`; decisão de produto (classe fora do D0-UX, tópico de produto, [PARÂMETRO]/[PROPOSTA] fechado sem parecer datado) bloqueia.\\[1pt]\campo{Objetivo}Copia seções de um relatório para chaves próprias do estado e registra um status informativo.\\[1pt]\campo{Entra}\cod{fechamento\_lacunas}\,\arq{design/00-fechamento-lacunas.md}, \cod{notas\_de\_passagem}, \cod{pedidos\_do\_po}, \cod{spec\_missao\_gate0}\,\arq{produto/05a-spec-missao-gate0.md}, \cod{spec\_missao}\,\arq{produto/05-spec-missao.md}\\[1pt]\campo{Sai}\cod{relatorio\_lacunas}\,\arq{design/00a-relatorio-lacunas.md}, \cod{lacunas\_status}};
\draw[fluxo] (n78.south) -- (n79.north);
\node[pmodelo, below=of n79] (n80) {{\sffamily\bfseries\small 067\hspace{0.5em}Ambiente simulado}\hfill{\sffamily\tiny\color{white}\colorbox{unLavanda}{MODELO}}\\[1pt]{\tiny\ttfamily\color{unLavanda}A067AmbienteSimulado}\\[2pt]\campo{Faz}Condicional: decide o ambiente (app clonado existente, app novo ou mini-jogo) quando a spec cala; dormente emite só a ficha do bloco 14 que D1, D2, D4, D5 e o QA consomem.\\[1pt]\campo{Objetivo}Agente condicional: roda depois do D0 (fechamento de lacunas) e antes do D1 (briefing de ofício), e só acorda quando a spec não nomeia o ambiente simulado da missão — ou nomeia um que não existe como seção do Design System…\\[1pt]\campo{Entra}\cod{anexo\_design\_spec}\,\arq{design/00-anexo-design-spec.md}, \cod{catalogo\_missoes}, \cod{gate\_estrutural\_report}\,\arq{produto/01-gate-estrutural.md}, \cod{one\_pager}\,\arq{produto/00-one-pager.md}, \cod{spec\_missao}\,\arq{produto/05-spec-missao.md}, \textit{+1 de controle}\\[1pt]\campo{Sai}\cod{ambiente\_simulado}\,\arq{design/00b-ambiente-simulado.md}};
\draw[fluxo] (n79.south) -- (n80.north);
\node[pportao, below=of n80] (n81) {{\sffamily\bfseries\small 068\hspace{0.5em}Anexo com bloco 14 (determinístico)}\hfill{\sffamily\tiny\color{white}\colorbox{unAzul!80!unNavy}{PORTÃO}}\\[1pt]{\tiny\ttfamily\color{unLavanda}A068AmbienteSimuladoSplit}\\[2pt]\campo{Faz}Grava o bloco 14 do D9 no anexo de design, sem tocar a spec de produto, e registra `ambiente\_status`.\\[1pt]\campo{Objetivo}Copia seções de um relatório para chaves próprias do estado e registra um status informativo.\\[1pt]\campo{Entra}\cod{ambiente\_simulado}\,\arq{design/00b-ambiente-simulado.md}, \cod{notas\_de\_passagem}, \cod{pedidos\_do\_po}, \cod{spec\_missao\_gate0}\,\arq{produto/05a-spec-missao-gate0.md}, \cod{spec\_missao}\,\arq{produto/05-spec-missao.md}\\[1pt]\campo{Sai}\cod{ambiente\_status}};
\draw[fluxo] (n80.south) -- (n81.north);
\end{tikzpicture}
\end{center}
\newpage
\begin{center}
\begin{tikzpicture}[node distance=0.42cm]
\node[pportao] (n82) {{\sffamily\bfseries\small 069\hspace{0.5em}Spec consolidada $\times$ one-pager aprovado (determinístico)}\hfill{\sffamily\tiny\color{white}\colorbox{unAzul!80!unNavy}{PORTÃO}}\\[1pt]{\tiny\ttfamily\color{unLavanda}A069SpecConsolidada}\\[2pt]\campo{Faz}A spec de produto mais o anexo de design (D0 e D9) contra o one-pager aprovado: habilidades, tempo em tela, podcast, Canal B, debate e parâmetros a definir. Só o que o anexo piorou contra a spec de produto; o conserto vai ao D0 pelo QA.\\[1pt]\campo{Objetivo}Gate determinístico de rastro: modo decide a verificação de tools/rastro.py.\\[1pt]\campo{Entra}\cod{spec\_missao\_gate0}\,\arq{produto/05a-spec-missao-gate0.md}, \cod{spec\_missao}\,\arq{produto/05-spec-missao.md}, \cod{one\_pager}\,\arq{produto/00-one-pager.md}\\[1pt]\campo{Sai}\cod{conteudo\_aprovado\_spec\_design\_report}\,\arq{design/00c-spec-consolidada-x-one-pager.md}, \cod{conteudo\_aprovado\_spec\_design\_status}};
\node[pmodelo, below=of n82] (n83) {{\sffamily\bfseries\small 070\hspace{0.5em}Briefing de ofício}\hfill{\sffamily\tiny\color{white}\colorbox{unLavanda}{MODELO}}\\[1pt]{\tiny\ttfamily\color{unLavanda}A070BriefOficio}\\[2pt]\campo{Faz}Aposta, forma do arco, pico, beat memorável, mapa de beats, não-negociáveis.\\[1pt]\campo{Objetivo}Roda depois do fechamento de lacunas (D0) e do ambiente simulado (D9): quando você começa, a spec já está consolidada e o ambiente já está decidido.\\[1pt]\campo{Entra}\cod{ambiente\_simulado}\,\arq{design/00b-ambiente-simulado.md}, \cod{anexo\_design\_spec}\,\arq{design/00-anexo-design-spec.md}, \cod{catalogo\_missoes}, \cod{gate\_estrutural\_report}\,\arq{produto/01-gate-estrutural.md}, \cod{one\_pager}\,\arq{produto/00-one-pager.md}, \cod{spec\_missao}\,\arq{produto/05-spec-missao.md}, \textit{+1 de controle}\\[1pt]\campo{Sai}\cod{brief\_oficio}\,\arq{design/01-briefing-oficio.md}};
\draw[fluxo] (n82.south) -- (n83.north);
\node[pmodelo, below=of n83] (n84) {{\sffamily\bfseries\small 071\hspace{0.5em}Jornada (Mermaid)}\hfill{\sffamily\tiny\color{white}\colorbox{unLavanda}{MODELO}}\\[1pt]{\tiny\ttfamily\color{unLavanda}A071Jornada}\\[2pt]\campo{Faz}Raias $\times$ 6 etapas, cenários de dispositivo, inventário de telas e direção por tela.\\[1pt]\campo{Objetivo}Anexo de design da spec (decisões de UI/UX do D0 e ficha do ambiente do D9, em design/00-anexo-design-spec.md — junto com a spec de produto, é o contrato; vazio numa missão anterior a 24/10/2026, que traz as duas coisas dentro…\\[1pt]\campo{Entra}\cod{ambiente\_simulado}\,\arq{design/00b-ambiente-simulado.md}, \cod{anexo\_design\_spec}\,\arq{design/00-anexo-design-spec.md}, \cod{brief\_oficio}\,\arq{design/01-briefing-oficio.md}, \cod{one\_pager}\,\arq{produto/00-one-pager.md}, \cod{spec\_missao}\,\arq{produto/05-spec-missao.md}, \textit{+1 de controle}\\[1pt]\campo{Sai}\cod{jornada\_telas}\,\arq{design/02-jornada-telas.md}};
\draw[fluxo] (n83.south) -- (n84.north);
\node[pmodelo, below=of n84] (n85) {{\sffamily\bfseries\small 072\hspace{0.5em}Fluxo de telas (Mermaid)}\hfill{\sffamily\tiny\color{white}\colorbox{unLavanda}{MODELO}}\\[1pt]{\tiny\ttfamily\color{unLavanda}A072Fluxo}\\[2pt]\campo{Faz}Grafo de telas com condições, gates e os 5 testes de fluxo.\\[1pt]\campo{Objetivo}Jornada e fluxo não são o mesmo artefato.\\[1pt]\campo{Entra}\cod{anexo\_design\_spec}\,\arq{design/00-anexo-design-spec.md}, \cod{brief\_oficio}\,\arq{design/01-briefing-oficio.md}, \cod{conteudo\_aprovado}, \cod{jornada\_telas}\,\arq{design/02-jornada-telas.md}, \cod{spec\_missao}\,\arq{produto/05-spec-missao.md}, \textit{+1 de controle}\\[1pt]\campo{Sai}\cod{fluxo\_arestas}\,\arq{design/03-fluxo-arestas.md}};
\draw[fluxo] (n84.south) -- (n85.north);
\node[pmodelo, below=of n85] (n86) {{\sffamily\bfseries\small 073\hspace{0.5em}Copy PT-BR (passada 1)}\hfill{\sffamily\tiny\color{white}\colorbox{unLavanda}{MODELO}}\\[1pt]{\tiny\ttfamily\color{unLavanda}A073CopyPtBr}\\[2pt]\campo{Faz}Sheet de strings por tela, calibrado à faixa etária.\\[1pt]\campo{Objetivo}Você escreve o que o estudante e o professor leem na tela.\\[1pt]\campo{Entra}\cod{ambiente\_simulado}\,\arq{design/00b-ambiente-simulado.md}, \cod{anexo\_design\_spec}\,\arq{design/00-anexo-design-spec.md}, \cod{brief\_oficio}\,\arq{design/01-briefing-oficio.md}, \cod{fluxo\_arestas}\,\arq{design/03-fluxo-arestas.md}, \cod{jornada\_telas}\,\arq{design/02-jornada-telas.md}, \cod{one\_pager}\,\arq{produto/00-one-pager.md}, \cod{spec\_missao}\,\arq{produto/05-spec-missao.md}, \cod{textos\_chave\_tabela}\,\arq{produto/00n-textos-chave.md}\\[1pt]\campo{Sai}\cod{copy\_strings}\,\arq{design/04-copy-strings.md}};
\draw[fluxo] (n85.south) -- (n86.north);
\node[pmodelo, below=of n86] (n87) {{\sffamily\bfseries\small 074\hspace{0.5em}Camada CAA}\hfill{\sffamily\tiny\color{white}\colorbox{unLavanda}{MODELO}}\\[1pt]{\tiny\ttfamily\color{unLavanda}A074Caa}\\[2pt]\campo{Faz}Símbolos, modos de resposta alternativos e pranchas; só roda quando a faixa ou a spec pedem.\\[1pt]\campo{Objetivo}Roda entre a copy passada 1 e os wireframes, e só quando a missão justifica.\\[1pt]\campo{Entra}\cod{anexo\_design\_spec}\,\arq{design/00-anexo-design-spec.md}, \cod{brief\_oficio}\,\arq{design/01-briefing-oficio.md}, \cod{conteudo\_aprovado}, \cod{copy\_strings}\,\arq{design/04-copy-strings.md}, \cod{jornada\_telas}\,\arq{design/02-jornada-telas.md}, \cod{spec\_missao}\,\arq{produto/05-spec-missao.md}, \textit{+1 de controle}\\[1pt]\campo{Sai}\cod{camada\_caa}\,\arq{design/04b-camada-caa.md}};
\draw[fluxo] (n86.south) -- (n87.north);
\end{tikzpicture}
\end{center}
\newpage
\begin{center}
\begin{tikzpicture}[node distance=0.42cm]
\node[pmodelo] (n88) {{\sffamily\bfseries\small 075\hspace{0.5em}Wireframes lo-fi (Stitch)}\hfill{\sffamily\tiny\color{white}\colorbox{unLavanda}{MODELO}}\\[1pt]{\tiny\ttfamily\color{unLavanda}A075Telas}\\[2pt]\campo{Faz}Telas em escala de cinza sobre a seção do DS do ambiente.\\[1pt]\campo{Objetivo}Você constrói todas as telas do inventário — estudante e professor — em baixa fidelidade, escala de cinza, sobre a seção do design system do ambiente simulado.\\[1pt]\campo{Entra}\cod{ambiente\_simulado}\,\arq{design/00b-ambiente-simulado.md}, \cod{anexo\_design\_spec}\,\arq{design/00-anexo-design-spec.md}, \cod{brief\_oficio}\,\arq{design/01-briefing-oficio.md}, \cod{camada\_caa}\,\arq{design/04b-camada-caa.md}, \cod{conteudo\_aprovado}, \cod{copy\_strings}\,\arq{design/04-copy-strings.md}, \cod{fluxo\_arestas}\,\arq{design/03-fluxo-arestas.md}, \cod{jornada\_telas}\,\arq{design/02-jornada-telas.md}, \cod{spec\_missao}\,\arq{produto/05-spec-missao.md}, +1, \textit{+1 de controle}\\[1pt]\campo{Sai}\cod{wireframes\_cobertura}\,\arq{design/05-wireframes-cobertura.md}};
\node[pcodigo, below=of n88] (n89) {{\sffamily\bfseries\small 076\hspace{0.5em}Telas no Stitch (determinístico)}\hfill{\sffamily\tiny\color{white}\colorbox{unNavy}{CÓDIGO}}\\[1pt]{\tiny\ttfamily\color{unLavanda}A076StitchGeracao}\\[2pt]\campo{Faz}Gera no Stitch, fora do modelo, toda (tela, estado) da tabela do agente de telas e do inventário da jornada, em paralelo; o screenId é o que o Stitch devolve.\\[1pt]\campo{Objetivo}Determinístico: gera no Stitch toda (tela, estado) que o agente de telas especificou (tools/stitchgeracao.py).\\[1pt]\campo{Entra}\cod{wireframes\_cobertura}\,\arq{design/05-wireframes-cobertura.md}, \cod{stitch\_project\_id}, \cod{stitch\_design\_system}, \cod{spec\_missao}\,\arq{produto/05-spec-missao.md}, \cod{jornada\_telas}\,\arq{design/02-jornada-telas.md}, \cod{stitch\_links}\,\arq{design/10-stitch-telas.md}\\[1pt]\campo{Sai}\cod{stitch\_geracao}};
\draw[fluxo] (n88.south) -- (n89.north);
\node[pmodelo, below=of n89] (n90) {{\sffamily\bfseries\small 077\hspace{0.5em}Copy PT-BR (passada 2)}\hfill{\sffamily\tiny\color{white}\colorbox{unLavanda}{MODELO}}\\[1pt]{\tiny\ttfamily\color{unLavanda}A077CopyRevisao}\\[2pt]\campo{Faz}Revisão em contexto de tela.\\[1pt]\campo{Objetivo}Na passada 1 você escreveu as strings a partir do inventário.\\[1pt]\campo{Entra}\cod{anexo\_design\_spec}\,\arq{design/00-anexo-design-spec.md}, \cod{brief\_oficio}\,\arq{design/01-briefing-oficio.md}, \cod{camada\_caa}\,\arq{design/04b-camada-caa.md}, \cod{conteudo\_aprovado}, \cod{copy\_strings}\,\arq{design/04-copy-strings.md}, \cod{jornada\_telas}\,\arq{design/02-jornada-telas.md}, \cod{spec\_missao}\,\arq{produto/05-spec-missao.md}, \cod{wireframes\_cobertura}\,\arq{design/05-wireframes-cobertura.md}, \textit{+1 de controle}\\[1pt]\campo{Sai}\cod{copy\_strings\_revisada}\,\arq{design/06-copy-strings-revisada.md}};
\draw[fluxo] (n89.south) -- (n90.north);
\node[pcodigo, below=of n90] (n91) {{\sffamily\bfseries\small 078\hspace{0.5em}Copy preservada (determinístico)}\hfill{\sffamily\tiny\color{white}\colorbox{unNavy}{CÓDIGO}}\\[1pt]{\tiny\ttfamily\color{unLavanda}A078CopyPreservada}\\[2pt]\campo{Faz}Reanexa à copy revisada, sem mudança, as seções da passada 1 que a revisão não trouxe (prompts e schemas, banco curado, roteiro do podcast).\\[1pt]\campo{Objetivo}Determinístico: o que a passada 1 da copy escreveu e a revisão não reemitiu volta, sem mudança.\\[1pt]\campo{Entra}\cod{copy\_strings}\,\arq{design/04-copy-strings.md}\\[1pt]\campo{Sai}\cod{copy\_strings\_revisada}\,\arq{design/06-copy-strings-revisada.md}};
\draw[fluxo] (n90.south) -- (n91.north);
\node[pportao, below=of n91] (n92) {{\sffamily\bfseries\small 078b\hspace{0.5em}Textos-chave na copy (determinístico)}\hfill{\sffamily\tiny\color{white}\colorbox{unAzul!80!unNavy}{PORTÃO}}\\[1pt]{\tiny\ttfamily\color{unLavanda}A078bTextosChave}\\[2pt]\campo{Faz}Copia por código os textos-chave do one-pager aprovado (aspas do \S{}4, Canal B, opções de cada escolha, podcast, roteiro do professor) para a linha da copy que os traz, pelo ID ou pelo texto. A redação que o design tinha escrito vira [PROPOSTA — validar com PO], para o PO pedir no Refazer; opção a menos ou a mais bloqueia e volta ao dono da copy (PR 7).\\[1pt]\campo{Objetivo}Determinístico: os textos-chave copiados por código para a copy, e o gate que confere (PR 7, D1).\\[1pt]\campo{Entra}\cod{copy\_strings\_revisada}\,\arq{design/06-copy-strings-revisada.md}, \cod{spec\_missao}\,\arq{produto/05-spec-missao.md}, \cod{one\_pager}\,\arq{produto/00-one-pager.md}, \cod{edicoes\_pontuais}\\[1pt]\campo{Sai}\cod{textos\_chave\_report}\,\arq{design/06d-textos-chave.md}, \cod{textos\_chave\_status}};
\draw[fluxo] (n91.south) -- (n92.north);
\node[pmodelo, below=of n92] (n93) {{\sffamily\bfseries\small 079\hspace{0.5em}Protótipo navegável (HTML)}\hfill{\sffamily\tiny\color{white}\colorbox{unLavanda}{MODELO}}\\[1pt]{\tiny\ttfamily\color{unLavanda}A079Prototipo}\\[2pt]\campo{Faz}Arquivo HTML único e sem dependências com as telas navegáveis, o painel da etapa vindo do one-pager e as cadências do pico. Salvo como `design/08-prototipo.html`.\\[1pt]\campo{Objetivo}Roda depois da copy revisada (D6) e antes do QA de entregáveis (D7), que vai auditar o que você construiu.\\[1pt]\campo{Entra}\cod{ambiente\_simulado}\,\arq{design/00b-ambiente-simulado.md}, \cod{anexo\_design\_spec}\,\arq{design/00-anexo-design-spec.md}, \cod{brief\_oficio}\,\arq{design/01-briefing-oficio.md}, \cod{camada\_caa}\,\arq{design/04b-camada-caa.md}, \cod{copy\_strings\_revisada}\,\arq{design/06-copy-strings-revisada.md}, \cod{fluxo\_arestas}\,\arq{design/03-fluxo-arestas.md}, \cod{jornada\_telas}\,\arq{design/02-jornada-telas.md}, \cod{one\_pager}\,\arq{produto/00-one-pager.md}, \cod{spec\_missao}\,\arq{produto/05-spec-missao.md}, +1, \textit{+1 de controle}\\[1pt]\campo{Sai}\cod{prototipo\_html}};
\draw[fluxo] (n92.south) -- (n93.north);
\end{tikzpicture}
\end{center}
\newpage
\begin{center}
\begin{tikzpicture}[node distance=0.42cm]
\node[pcodigo] (n94) {{\sffamily\bfseries\small 080\hspace{0.5em}Artefatos e links (determinístico)}\hfill{\sffamily\tiny\color{white}\colorbox{unNavy}{CÓDIGO}}\\[1pt]{\tiny\ttfamily\color{unLavanda}A080ArtefatosDesign}\\[2pt]\campo{Faz}Separa o HTML do protótipo do relatório, colhe os links das telas no Stitch e monta a folha de diagramas. Stitch com menos telas que o inventário da jornada bloqueia (`stitch\_status`, dono `telas`).\\[1pt]\campo{Objetivo}Determinístico: separa o protótipo do relatório, colhe os links do Stitch e monta os diagramas.\\[1pt]\campo{Entra}\cod{prototipo\_html}, \cod{prototipo\_arquivo}\,\arq{design/08-prototipo.html}, \cod{prototipo\_relatorio}\,\arq{design/08-prototipo.md}, \cod{diagramas\_html}\,\arq{design/09-diagramas.html}, \cod{spec\_missao}\,\arq{produto/05-spec-missao.md}, \cod{jornada\_telas}\,\arq{design/02-jornada-telas.md}, \cod{fluxo\_arestas}\,\arq{design/03-fluxo-arestas.md}, \cod{stitch\_project\_id}, \cod{stitch\_geracao}, +1\\[1pt]\campo{Sai}\cod{stitch\_links}\,\arq{design/10-stitch-telas.md}, \cod{stitch\_status}};
\node[pportao, below=of n94] (n95) {{\sffamily\bfseries\small 081\hspace{0.5em}Lint do protótipo (determinístico)}\hfill{\sffamily\tiny\color{white}\colorbox{unAzul!80!unNavy}{PORTÃO}}\\[1pt]{\tiny\ttfamily\color{unLavanda}A081PrototipoLint}\\[2pt]\campo{Faz}Confere o arquivo do protótipo HTML: o que o QA de entregáveis não vê por ler só o texto. Reprovado devolve ao protótipo no retrabalho de design.\\[1pt]\campo{Objetivo}Determinístico: confere o HTML do protótipo contra o inventário de telas da jornada.\\[1pt]\campo{Entra}\cod{prototipo\_arquivo}\,\arq{design/08-prototipo.html}, \cod{jornada\_telas}\,\arq{design/02-jornada-telas.md}, \cod{wireframes\_cobertura}\,\arq{design/05-wireframes-cobertura.md}, \cod{spec\_missao}\,\arq{produto/05-spec-missao.md}, \cod{one\_pager}\,\arq{produto/00-one-pager.md}\\[1pt]\campo{Sai}\cod{prototipo\_lint\_report}\,\arq{design/06b-lint-prototipo.md}, \cod{prototipo\_lint\_status}};
\draw[fluxo] (n94.south) -- (n95.north);
\node[pportao, below=of n95] (n96) {{\sffamily\bfseries\small 082\hspace{0.5em}Coerência entre documentos (determinístico)}\hfill{\sffamily\tiny\color{white}\colorbox{unAzul!80!unNavy}{PORTÃO}}\\[1pt]{\tiny\ttfamily\color{unLavanda}A082Consistencia}\\[2pt]\campo{Faz}Jornada, fluxo, copy revisada e protótipo contra o glossário canônico da spec: variante proibida de entidade bloqueia; papel ou idade de personagem, número com unidade, rótulo de botão e nome de arquivo divergentes ficam . Pegou os quatro nomes do órgão da exp09ef1b.\\[1pt]\campo{Objetivo}Gate determinístico de rastro: modo decide a verificação de tools/rastro.py.\\[1pt]\campo{Entra}\cod{copy\_strings\_revisada}\,\arq{design/06-copy-strings-revisada.md}, \cod{one\_pager}\,\arq{produto/00-one-pager.md}, \cod{camada\_caa}\,\arq{design/04b-camada-caa.md}, \cod{prototipo\_arquivo}\,\arq{design/08-prototipo.html}\\[1pt]\campo{Sai}\cod{consistencia\_report}\,\arq{design/06c-coerencia-documentos.md}, \cod{consistencia\_status}};
\draw[fluxo] (n95.south) -- (n96.north);
\node[pportao, below=of n96] (n97) {{\sffamily\bfseries\small 083\hspace{0.5em}Cobertura de telas (determinístico)}\hfill{\sffamily\tiny\color{white}\colorbox{unAzul!80!unNavy}{PORTÃO}}\\[1pt]{\tiny\ttfamily\color{unLavanda}A083CoberturaTelas}\\[2pt]\campo{Faz}Toda tela do inventário da jornada no fluxo, na copy revisada, na cobertura dos wireframes e nos quadros do protótipo, conferida por ID; falta vira ticket com o dono (`cobertura\_telas\_status`).\\[1pt]\campo{Objetivo}Determinístico: toda tela do inventário da jornada em todo artefato que a desenha, por ID.\\[1pt]\campo{Entra}\cod{spec\_missao}\,\arq{produto/05-spec-missao.md}, \cod{one\_pager}\,\arq{produto/00-one-pager.md}\\[1pt]\campo{Sai}\cod{cobertura\_telas\_report}, \cod{cobertura\_telas\_status}};
\draw[fluxo] (n96.south) -- (n97.north);
\node[pportao, below=of n97] (n98) {{\sffamily\bfseries\small 083b\hspace{0.5em}Fluxo e telas $\times$ bloco 2 (determinístico)}\hfill{\sffamily\tiny\color{white}\colorbox{unAzul!80!unNavy}{PORTÃO}}\\[1pt]{\tiny\ttfamily\color{unLavanda}A083bArestasContrato}\\[2pt]\campo{Faz}Arestas do fluxo e regras de habilitação das telas contra a mecânica da spec de produto: guarda que só deixa seguir com a resposta certa, ou que troca o «e» de um gatilho da spec por «ou», vira [PROPOSTA — validar com PO]; estado novo de decomposição é livre (D7). Bloqueia com ESTEIRA\_ESTRITA; sem ela, é decisão humana no cartão de design (PR 7).\\[1pt]\campo{Objetivo}Determinístico: arestas do fluxo e regras de habilitação das telas $\times$ bloco 2 da spec (tools/arestas.py, PR 7).\\[1pt]\campo{Entra}\cod{fluxo\_arestas}\,\arq{design/03-fluxo-arestas.md}, \cod{spec\_missao}\,\arq{produto/05-spec-missao.md}, \cod{one\_pager}\,\arq{produto/00-one-pager.md}\\[1pt]\campo{Sai}\cod{arestas\_report}\,\arq{design/03b-fluxo-x-spec.md}, \cod{arestas\_status}};
\draw[fluxo] (n97.south) -- (n98.north);
\end{tikzpicture}
\end{center}
\newpage
\begin{center}
\begin{tikzpicture}[node distance=0.42cm]
\node[pcodigo] (n99) {{\sffamily\bfseries\small 084\hspace{0.5em}Espelho no Miro (determinístico)}\hfill{\sffamily\tiny\color{white}\colorbox{unNavy}{CÓDIGO}}\\[1pt]{\tiny\ttfamily\color{unLavanda}A084EspelhoMiro}\\[2pt]\campo{Faz}Publica no Miro, como board navegável, os mesmos diagramas Mermaid da jornada e do fluxo. O Mermaid continua sendo a fonte; Miro indisponível não interrompe a esteira.\\[1pt]\campo{Objetivo}Determinístico: espelha no Miro o Mermaid que a jornada e o fluxo já escreveram.\\[1pt]\campo{Entra}\cod{jornada\_telas}\,\arq{design/02-jornada-telas.md}, \cod{fluxo\_arestas}\,\arq{design/03-fluxo-arestas.md}, \cod{miro\_board}, \cod{spec\_missao}\,\arq{produto/05-spec-missao.md}\\[1pt]\campo{Sai}\cod{miro\_links}\,\arq{design/11-miro-boards.md}};
\node[pcodigo, below=of n99] (n100) {{\sffamily\bfseries\small 085\hspace{0.5em}Escopo da banca de design (determinístico)}\hfill{\sffamily\tiny\color{white}\colorbox{unNavy}{CÓDIGO}}\\[1pt]{\tiny\ttfamily\color{unLavanda}A085EscopoBanca}\\[2pt]\ciclo{$\circlearrowright$ começa aqui: \cod{QaDesignLoop} (até 2 rodadas)}\\[1pt]\campo{Faz}Grava o escopo da banca na posição «design».\\[1pt]\campo{Objetivo}Grava em escopobanca onde a banca está rodando e o que ela pode fazer ali (tools/banca.ESCOPOS).\\[1pt]\campo{Entra}\textit{nada do estado da sessão (configuração fixa do passo)}\\[1pt]\campo{Sai}\cod{escopo\_banca}\\[1pt]\campo{Ciclo}\cod{QaDesignLoop}};
\draw[fluxo] (n99.south) -- (n100.north);
\node[pmodelo, below=of n100] (n101) {{\sffamily\bfseries\small 086a\hspace{0.5em}Banca no fim de design: produto realizado em tela}\hfill{\sffamily\tiny\color{white}\colorbox{unLavanda}{MODELO}}\\[1pt]{\tiny\ttfamily\color{unLavanda}A086aBancaProdutoNarrativa}\\[2pt]\campo{Faz}Os critérios PROD que se realizam em tela — fluxo $\times$ contrato de estados, momento de conclusão, condução do debate, modos de erro —, antes do parecer de design.\\[1pt]\campo{Objetivo}A mesma banca roda em quatro posições, com os mesmos critérios — fim de produto, fim de design, entrada do dev e dentro do dev.\\[1pt]\campo{Entra}\cod{anexo\_design\_spec}\,\arq{design/00-anexo-design-spec.md}, \cod{banco\_conteudo}\,\arq{dev/01e-banco-conteudo.md}, \cod{brief\_assets}\,\arq{dev/01b-brief-assets.md}, \cod{brief\_oficio}\,\arq{design/01-briefing-oficio.md}, \cod{camada\_caa}\,\arq{design/04b-camada-caa.md}, \cod{conteudo\_aprovado}, \cod{copy\_strings\_revisada}\,\arq{design/06-copy-strings-revisada.md}, \cod{escopo\_banca}, \cod{fluxo\_arestas}\,\arq{design/03-fluxo-arestas.md}, +6\\[1pt]\campo{Sai}\cod{banca\_design\_narrativa}\\[1pt]\campo{Ciclo}\cod{QaDesignLoop}\\[1pt]\campo{Roda}em paralelo com os demais de \cod{A086Banca}};
\draw[fluxo] (n100.south) -- (n101.north);
\node[pmodelo, below=of n101] (n102) {{\sffamily\bfseries\small 086b\hspace{0.5em}Banca no fim de design: design e imersão}\hfill{\sffamily\tiny\color{white}\colorbox{unLavanda}{MODELO}}\\[1pt]{\tiny\ttfamily\color{unLavanda}A086bBancaDesignImersao}\\[2pt]\campo{Faz}Os critérios DES da banca sobre os documentos de design, antes do parecer de design.\\[1pt]\campo{Objetivo}A mesma banca roda em quatro posições, com os mesmos critérios — fim de produto, fim de design, entrada do dev e dentro do dev.\\[1pt]\campo{Entra}\cod{anexo\_design\_spec}\,\arq{design/00-anexo-design-spec.md}, \cod{banco\_conteudo}\,\arq{dev/01e-banco-conteudo.md}, \cod{brief\_assets}\,\arq{dev/01b-brief-assets.md}, \cod{brief\_oficio}\,\arq{design/01-briefing-oficio.md}, \cod{camada\_caa}\,\arq{design/04b-camada-caa.md}, \cod{conteudo\_aprovado}, \cod{copy\_strings\_revisada}\,\arq{design/06-copy-strings-revisada.md}, \cod{escopo\_banca}, \cod{fluxo\_arestas}\,\arq{design/03-fluxo-arestas.md}, +6\\[1pt]\campo{Sai}\cod{banca\_design\_imersao}\\[1pt]\campo{Ciclo}\cod{QaDesignLoop}\\[1pt]\campo{Roda}em paralelo com os demais de \cod{A086Banca}};
\draw[fluxo] (n101.south) -- (n102.north);
\node[pportao, below=of n102] (n103) {{\sffamily\bfseries\small 087\hspace{0.5em}Banca de design (determinístico)}\hfill{\sffamily\tiny\color{white}\colorbox{unAzul!80!unNavy}{PORTÃO}}\\[1pt]{\tiny\ttfamily\color{unLavanda}A087BancaConsolidacao}\\[2pt]\campo{Faz}Veredito da banca no fim de design. O bloqueio vira ticket pelo A089 e o dono refaz nesta etapa; esgotadas as rodadas, recusa o de acordo.\\[1pt]\campo{Objetivo}Determinístico: o veredito da banca é a pior das cinco dimensões, contada em código.\\[1pt]\campo{Entra}\cod{banca\_design\_narrativa}, \cod{banca\_design\_imersao}, \cod{spec\_missao}\,\arq{produto/05-spec-missao.md}, \cod{one\_pager}\,\arq{produto/00-one-pager.md}\\[1pt]\campo{Sai}\cod{banca\_design\_report}\,\arq{design/07c-banca-design.md}, \cod{banca\_design\_status}\\[1pt]\campo{Ciclo}\cod{QaDesignLoop}};
\draw[fluxo] (n102.south) -- (n103.north);
\node[pmodelo, below=of n103] (n104) {{\sffamily\bfseries\small 088\hspace{0.5em}QA de entregáveis}\hfill{\sffamily\tiny\color{white}\colorbox{unLavanda}{MODELO}}\\[1pt]{\tiny\ttfamily\color{unLavanda}A088QaEntregaveis}\\[2pt]\campo{Faz}Fidelidade à spec, diff de strings, currículo, faixa, fluxo, professor, briefing, ambiente (D8) e CAA (D8b).\\[1pt]\campo{Objetivo}Seu escopo é a esteira de design, não a de produto.\\[1pt]\campo{Entra}\cod{ambiente\_simulado}\,\arq{design/00b-ambiente-simulado.md}, \cod{anexo\_design\_spec}\,\arq{design/00-anexo-design-spec.md}, \cod{brief\_oficio}\,\arq{design/01-briefing-oficio.md}, \cod{camada\_caa}\,\arq{design/04b-camada-caa.md}, \cod{consistencia\_report}\,\arq{design/06c-coerencia-documentos.md}, \cod{conteudo\_aprovado}, \cod{copy\_strings\_revisada}\,\arq{design/06-copy-strings-revisada.md}, \cod{fluxo\_arestas}\,\arq{design/03-fluxo-arestas.md}, \cod{jornada\_telas}\,\arq{design/02-jornada-telas.md}, +7\\[1pt]\campo{Sai}\cod{qa\_entregaveis\_report}\,\arq{design/07-qa-entregaveis.md}\\[1pt]\campo{Ciclo}\cod{QaDesignLoop}};
\draw[fluxo] (n103.south) -- (n104.north);
\draw[loopfaixa] ([xshift=-0.25cm]n100.north west) -- ([xshift=-0.25cm]n104.south west);
\end{tikzpicture}
\end{center}
\newpage
\begin{center}
\begin{tikzpicture}[node distance=0.42cm]
\node[pportao] (n105) {{\sffamily\bfseries\small 089\hspace{0.5em}Afirmações do QA (determinístico)}\hfill{\sffamily\tiny\color{white}\colorbox{unAzul!80!unNavy}{PORTÃO}}\\[1pt]{\tiny\ttfamily\color{unLavanda}A089QaAfirmacoes}\\[2pt]\campo{Faz}Contagens de telas ("todas as N telas", "N de M") e nomes de arquivo que o QA escreveu, contra o inventário da jornada, o Stitch, o protótipo e `pipeline/artifacts.py`. Divergência reescreve o veredito do QA como bloqueado; Stitch com menos telas que o inventário vira ticket para `telas`. Pegou as "16 telas" de um Stitch com 6 na exp09ef1b.\\[1pt]\campo{Objetivo}Confere em código as contagens de telas e os nomes de arquivo que o QA de design escreveu.\\[1pt]\campo{Entra}\cod{qa\_entregaveis\_report}\,\arq{design/07-qa-entregaveis.md}, \cod{stitch\_links}\,\arq{design/10-stitch-telas.md}, \cod{jornada\_telas}\,\arq{design/02-jornada-telas.md}, \cod{prototipo\_arquivo}\,\arq{design/08-prototipo.html}, \cod{spec\_missao}\,\arq{produto/05-spec-missao.md}\\[1pt]\campo{Sai}\cod{qa\_afirmacoes\_report}\,\arq{design/07b-qa-afirmacoes.md}, \cod{qa\_afirmacoes\_status}\\[1pt]\campo{Ciclo}\cod{QaDesignLoop}};
\node[pportao, below=of n105] (n106) {{\sffamily\bfseries\small 090\hspace{0.5em}Veredito de design (determinístico)}\hfill{\sffamily\tiny\color{white}\colorbox{unAzul!80!unNavy}{PORTÃO}}\\[1pt]{\tiny\ttfamily\color{unLavanda}A090QaEntregaveisCheck}\\[2pt]\campo{Faz}Lê o envelope do QA de entregáveis.\\[1pt]\campo{Objetivo}Lê o veredito de um relatório e decide se o pipeline segue, repete ou para.\\[1pt]\campo{Entra}\cod{spec\_missao}\,\arq{produto/05-spec-missao.md}, \cod{anexo\_design\_spec}\,\arq{design/00-anexo-design-spec.md}, \cod{jornada\_telas}\,\arq{design/02-jornada-telas.md}, \cod{fluxo\_arestas}\,\arq{design/03-fluxo-arestas.md}, \cod{copy\_strings\_revisada}\,\arq{design/06-copy-strings-revisada.md}, \cod{camada\_caa}\,\arq{design/04b-camada-caa.md}, \cod{wireframes\_cobertura}\,\arq{design/05-wireframes-cobertura.md}, \cod{prototipo\_arquivo}\,\arq{design/08-prototipo.html}, \cod{prototipo\_relatorio}\,\arq{design/08-prototipo.md}, +4\\[1pt]\campo{Sai}\cod{design\_status}, \textit{+1 de controle}\\[1pt]\campo{Ciclo}\cod{QaDesignLoop}};
\draw[fluxo] (n105.south) -- (n106.north);
\node[pcodigo, below=of n106] (n107) {{\sffamily\bfseries\small 091\hspace{0.5em}Retrabalho de design (determinístico)}\hfill{\sffamily\tiny\color{white}\colorbox{unNavy}{CÓDIGO}}\\[1pt]{\tiny\ttfamily\color{unLavanda}A091DespachoDesign}\\[2pt]\campo{Faz}Reexecuta só os agentes com ticket do QA, na ordem de precedência do prompt, e refaz o protótipo depois.\\[1pt]\campo{Objetivo}Re-executa apenas os agentes endereçados pelos tickets de um gate.\\[1pt]\campo{Entra}\cod{qa\_entregaveis\_report}\,\arq{design/07-qa-entregaveis.md}\\[1pt]\campo{Sai}\textit{+1 de controle}\\[1pt]\campo{Ciclo}\cod{QaDesignLoop}\\[1pt]\ciclo{$\circlearrowleft$ fim do ciclo: \cod{QaDesignLoop} --- reprovado, volta ao passo 085}};
\draw[fluxo] (n106.south) -- (n107.north);
\node[pportao, below=of n107] (n108) {{\sffamily\bfseries\small 092\hspace{0.5em}Rastro de lacunas do design (determinístico)}\hfill{\sffamily\tiny\color{white}\colorbox{unAzul!80!unNavy}{PORTÃO}}\\[1pt]{\tiny\ttfamily\color{unLavanda}A092RastroDesign}\\[2pt]\campo{Faz}Marcadores de lacuna na spec de produto, no anexo de design e nos documentos de design: o de design bloqueia a etapa; o que o D0 fechou (mapa célula $\rightarrow$ D0 do bloco 13) e a completude de produto (mapa do bloco 15) não conta; o de produto já aprovado e o de dev ficam ; [D0-n] citado e ausente do bloco 13 do anexo bloqueia.\\[1pt]\campo{Objetivo}Gate determinístico de rastro: modo decide a verificação de tools/rastro.py.\\[1pt]\campo{Entra}\cod{spec\_missao}\,\arq{produto/05-spec-missao.md}, \cod{anexo\_design\_spec}\,\arq{design/00-anexo-design-spec.md}, \cod{jornada\_telas}\,\arq{design/02-jornada-telas.md}, \cod{fluxo\_arestas}\,\arq{design/03-fluxo-arestas.md}, \cod{copy\_strings\_revisada}\,\arq{design/06-copy-strings-revisada.md}, \cod{wireframes\_cobertura}\,\arq{design/05-wireframes-cobertura.md}, \cod{stitch\_links}\,\arq{design/10-stitch-telas.md}, \cod{camada\_caa}\,\arq{design/04b-camada-caa.md}, \cod{ambiente\_simulado}\,\arq{design/00b-ambiente-simulado.md}\\[1pt]\campo{Sai}\cod{rastro\_design\_report}\,\arq{design/12-rastro-lacunas-design.md}, \cod{rastro\_design\_status}};
\draw[fluxo] (n107.south) -- (n108.north);
\node[pportao, below=of n108] (n109) {{\sffamily\bfseries\small 093\hspace{0.5em}Matriz de rastreabilidade (determinístico)}\hfill{\sffamily\tiny\color{white}\colorbox{unAzul!80!unNavy}{PORTÃO}}\\[1pt]{\tiny\ttfamily\color{unLavanda}A093RastreabilidadeDesign}\\[2pt]\campo{Faz}Cada requisito do produto (habilidades, textos literais, tempo, telas, estados, seção do DS) procurado em cada artefato de design; o que se perdeu sai com o dono.\\[1pt]\campo{Objetivo}Determinístico: a matriz requisito do produto $\times$ artefato (tools/rastreabilidade.py).\\[1pt]\campo{Entra}\cod{spec\_missao}\,\arq{produto/05-spec-missao.md}, \cod{one\_pager}\,\arq{produto/00-one-pager.md}\\[1pt]\campo{Sai}\cod{rastreabilidade\_report}\,\arq{matriz-de-rastreabilidade.md}, \cod{rastreabilidade\_status}};
\draw[fluxo] (n108.south) -- (n109.north);
\draw[loopfaixa] ([xshift=-0.25cm]n105.north west) -- ([xshift=-0.25cm]n107.south west);
\end{tikzpicture}
\end{center}
\newpage
\begin{center}
\begin{tikzpicture}[node distance=0.42cm]
\node[pportao] (n110) {{\sffamily\bfseries\small 094\hspace{0.5em}Fidelidade ao refazer de design}\hfill{\sffamily\tiny\color{white}\colorbox{unAzul!80!unNavy}{PORTÃO}}\\[1pt]{\tiny\ttfamily\color{unLavanda}A094FidelidadeDesign}\\[2pt]\campo{Faz}Cada pedido de refazer feito ao design (TRAVAS, AJUSTES, VALIDAÇÃO, frases entre aspas) conferido no anexo de design, no brief, na jornada, no fluxo, na copy, na CAA e nas telas: o juiz cita, o código confere. Sem pedido de design, não roda.  recusa o de acordo.\\[1pt]\campo{Objetivo}O que o PO pediu e o briefing fixou está nos documentos?\\[1pt]\campo{Entra}\cod{anexo\_design\_spec}\,\arq{design/00-anexo-design-spec.md}, \cod{brief\_oficio}\,\arq{design/01-briefing-oficio.md}, \cod{jornada\_telas}\,\arq{design/02-jornada-telas.md}, \cod{fluxo\_arestas}\,\arq{design/03-fluxo-arestas.md}, \cod{copy\_strings\_revisada}\,\arq{design/06-copy-strings-revisada.md}, \cod{camada\_caa}\,\arq{design/04b-camada-caa.md}, \cod{wireframes\_cobertura}\,\arq{design/05-wireframes-cobertura.md}, \cod{divisao\_do\_pedido}, \cod{pedidos\_do\_po}, +4, \textit{+3 de controle}\\[1pt]\campo{Sai}\cod{fidelidade\_design\_report}\,\arq{design/13-fidelidade-refazer-design.md}, \cod{fidelidade\_design\_status}};
\end{tikzpicture}
\end{center}
\newpage
\markboth{Dev}{}
{\sffamily\bfseries\Large\color{unNavy}Dev}\\[0.3em]
\begin{center}
\begin{tikzpicture}[node distance=0.42cm]
\node[pcodigo] (n111) {{\sffamily\bfseries\small 095\hspace{0.5em}Conteúdo aprovado (determinístico)}\hfill{\sffamily\tiny\color{white}\colorbox{unNavy}{CÓDIGO}}\\[1pt]{\tiny\ttfamily\color{unLavanda}A095ConteudoAprovadoDev}\\[2pt]\campo{Faz}As seções do one-pager aprovado, literais — as mesmas do design: núcleo (\S{}1, \S{}3, \S{}4), currículo, falhas, professor, tempos, podcast, evidência, dispositivos e hipóteses —, para o banco de conteúdo, a banca e os agentes de design que ela reaciona.\\[1pt]\campo{Objetivo}Grava em conteudoaprovado as seções literais do one-pager aprovado (tools/conteudoaprovado.py).\\[1pt]\campo{Entra}\cod{one\_pager}\,\arq{produto/00-one-pager.md}, \cod{briefing}\\[1pt]\campo{Sai}\cod{conteudo\_aprovado}};
\node[pcodigo, below=of n111] (n112) {{\sffamily\bfseries\small 096\hspace{0.5em}Escopo da banca de entrada do dev (determinístico)}\hfill{\sffamily\tiny\color{white}\colorbox{unNavy}{CÓDIGO}}\\[1pt]{\tiny\ttfamily\color{unLavanda}A096EscopoBanca}\\[2pt]\campo{Faz}Grava o escopo da banca na posição «entrada».\\[1pt]\campo{Objetivo}Grava em escopobanca onde a banca está rodando e o que ela pode fazer ali (tools/banca.ESCOPOS).\\[1pt]\campo{Entra}\textit{nada do estado da sessão (configuração fixa do passo)}\\[1pt]\campo{Sai}\cod{escopo\_banca}};
\draw[fluxo] (n111.south) -- (n112.north);
\node[pmodelo, below=of n112] (n113) {{\sffamily\bfseries\small 097a\hspace{0.5em}Banca de entrada: produto e narrativa}\hfill{\sffamily\tiny\color{white}\colorbox{unLavanda}{MODELO}}\\[1pt]{\tiny\ttfamily\color{unLavanda}A097aBancaProdutoNarrativa}\\[2pt]\campo{Faz}Segunda verificação de produto, antes de qualquer agente de dev, com a visão cruzada de produto e design. Não reescreve nada.\\[1pt]\campo{Objetivo}A mesma banca roda em quatro posições, com os mesmos critérios — fim de produto, fim de design, entrada do dev e dentro do dev.\\[1pt]\campo{Entra}\cod{anexo\_design\_spec}\,\arq{design/00-anexo-design-spec.md}, \cod{banco\_conteudo}\,\arq{dev/01e-banco-conteudo.md}, \cod{brief\_assets}\,\arq{dev/01b-brief-assets.md}, \cod{brief\_oficio}\,\arq{design/01-briefing-oficio.md}, \cod{camada\_caa}\,\arq{design/04b-camada-caa.md}, \cod{conteudo\_aprovado}, \cod{copy\_strings\_revisada}\,\arq{design/06-copy-strings-revisada.md}, \cod{escopo\_banca}, \cod{fluxo\_arestas}\,\arq{design/03-fluxo-arestas.md}, +6\\[1pt]\campo{Sai}\cod{banca\_entrada\_narrativa}\\[1pt]\campo{Roda}em paralelo com os demais de \cod{A097Banca}};
\draw[fluxo] (n112.south) -- (n113.north);
\node[pmodelo, below=of n113] (n114) {{\sffamily\bfseries\small 097b\hspace{0.5em}Banca de entrada: design e imersão}\hfill{\sffamily\tiny\color{white}\colorbox{unLavanda}{MODELO}}\\[1pt]{\tiny\ttfamily\color{unLavanda}A097bBancaDesignImersao}\\[2pt]\campo{Faz}Segunda verificação de design, com a visão cruzada: cada estado da spec com tela, a conclusão em tela, o professor com as telas dele, os textos aprovados na copy.\\[1pt]\campo{Objetivo}A mesma banca roda em quatro posições, com os mesmos critérios — fim de produto, fim de design, entrada do dev e dentro do dev.\\[1pt]\campo{Entra}\cod{anexo\_design\_spec}\,\arq{design/00-anexo-design-spec.md}, \cod{banco\_conteudo}\,\arq{dev/01e-banco-conteudo.md}, \cod{brief\_assets}\,\arq{dev/01b-brief-assets.md}, \cod{brief\_oficio}\,\arq{design/01-briefing-oficio.md}, \cod{camada\_caa}\,\arq{design/04b-camada-caa.md}, \cod{conteudo\_aprovado}, \cod{copy\_strings\_revisada}\,\arq{design/06-copy-strings-revisada.md}, \cod{escopo\_banca}, \cod{fluxo\_arestas}\,\arq{design/03-fluxo-arestas.md}, +6\\[1pt]\campo{Sai}\cod{banca\_entrada\_imersao}\\[1pt]\campo{Roda}em paralelo com os demais de \cod{A097Banca}};
\draw[fluxo] (n113.south) -- (n114.north);
\node[pmodelo, below=of n114] (n115) {{\sffamily\bfseries\small 097c\hspace{0.5em}Banca de entrada: pedagogia e BNCC}\hfill{\sffamily\tiny\color{white}\colorbox{unLavanda}{MODELO}}\\[1pt]{\tiny\ttfamily\color{unLavanda}A097cBancaPedagogiaBncc}\\[2pt]\campo{Faz}Segunda verificação dos critérios PED sobre a spec e o design aprovados.\\[1pt]\campo{Objetivo}A mesma banca roda em quatro posições, com os mesmos critérios — fim de produto, fim de design, entrada do dev e dentro do dev.\\[1pt]\campo{Entra}\cod{anexo\_design\_spec}\,\arq{design/00-anexo-design-spec.md}, \cod{banco\_conteudo}\,\arq{dev/01e-banco-conteudo.md}, \cod{brief\_assets}\,\arq{dev/01b-brief-assets.md}, \cod{brief\_oficio}\,\arq{design/01-briefing-oficio.md}, \cod{camada\_caa}\,\arq{design/04b-camada-caa.md}, \cod{conteudo\_aprovado}, \cod{copy\_strings\_revisada}\,\arq{design/06-copy-strings-revisada.md}, \cod{escopo\_banca}, \cod{fluxo\_arestas}\,\arq{design/03-fluxo-arestas.md}, +6\\[1pt]\campo{Sai}\cod{banca\_entrada\_pedagogia}\\[1pt]\campo{Roda}em paralelo com os demais de \cod{A097Banca}};
\draw[fluxo] (n114.south) -- (n115.north);
\end{tikzpicture}
\end{center}
\newpage
\begin{center}
\begin{tikzpicture}[node distance=0.42cm]
\node[pportao] (n116) {{\sffamily\bfseries\small 098\hspace{0.5em}Banca de entrada do dev (determinístico)}\hfill{\sffamily\tiny\color{white}\colorbox{unAzul!80!unNavy}{PORTÃO}}\\[1pt]{\tiny\ttfamily\color{unLavanda}A098BancaConsolidacao}\\[2pt]\campo{Faz}O pior das três dimensões.  para o dev antes de qualquer agente de dev (BANCA\_ENTRADA\_ESTRITA) e lista os bloqueios por dono; o que passou na etapa e reprova aqui sai numa seção própria.\\[1pt]\campo{Objetivo}Determinístico: o veredito da banca é a pior das cinco dimensões, contada em código.\\[1pt]\campo{Entra}\cod{banca\_entrada\_narrativa}, \cod{banca\_entrada\_imersao}, \cod{banca\_entrada\_pedagogia}, \cod{spec\_missao}\,\arq{produto/05-spec-missao.md}, \cod{one\_pager}\,\arq{produto/00-one-pager.md}\\[1pt]\campo{Sai}\cod{banca\_entrada\_report}\,\arq{dev/00b-banca-entrada.md}, \cod{banca\_entrada\_status}};
\node[pportao, below=of n116] (n117) {{\sffamily\bfseries\small 099\hspace{0.5em}Ficha do manifesto (determinístico)}\hfill{\sffamily\tiny\color{white}\colorbox{unAzul!80!unNavy}{PORTÃO}}\\[1pt]{\tiny\ttfamily\color{unLavanda}A099ManifestoCheck}\\[2pt]\campo{Faz}Bloco 0 da spec $\times$ listas fechadas da plataforma: produto, séries, unidade, público, viewport, permissões, capabilities, seção do DS, liberação por turma. Não conforme impede o PRD.\\[1pt]\campo{Objetivo}Gate determinístico de rastro: modo decide a verificação de tools/rastro.py.\\[1pt]\campo{Entra}\cod{briefing}, \cod{one\_pager\_po}\\[1pt]\campo{Sai}\cod{manifesto\_report}\,\arq{dev/00-ficha-manifesto.md}, \cod{manifesto\_status}};
\draw[fluxo] (n116.south) -- (n117.north);
\node[pmodelo, below=of n117] (n118) {{\sffamily\bfseries\small 100a\hspace{0.5em}Handoff para dev}\hfill{\sffamily\tiny\color{white}\colorbox{unLavanda}{MODELO}}\\[1pt]{\tiny\ttfamily\color{unLavanda}A100aHandoffDev}\\[2pt]\ciclo{$\circlearrowright$ começa aqui: \cod{HandoffLoop} (até 2 rodadas)}\\[1pt]\campo{Faz}Modelo de insumo preenchido, stories com , eventos, escopo V1/V2, hipóteses (paralelo com V1b).\\[1pt]\campo{Objetivo}Você fecha os critérios 9.1, 9.2 e 9.3 da rubrica.\\[1pt]\campo{Entra}\cod{ambiente\_simulado}\,\arq{design/00b-ambiente-simulado.md}, \cod{anexo\_design\_spec}\,\arq{design/00-anexo-design-spec.md}, \cod{brief\_oficio}\,\arq{design/01-briefing-oficio.md}, \cod{camada\_caa}\,\arq{design/04b-camada-caa.md}, \cod{copy\_strings\_revisada}\,\arq{design/06-copy-strings-revisada.md}, \cod{design\_status}, \cod{fluxo\_arestas}\,\arq{design/03-fluxo-arestas.md}, \cod{jornada\_telas}\,\arq{design/02-jornada-telas.md}, \cod{miro\_links}\,\arq{design/11-miro-boards.md}, +11\\[1pt]\campo{Sai}\cod{handoff\_dev}\,\arq{dev/01-handoff-dev.md}\\[1pt]\campo{Ciclo}\cod{HandoffLoop}\\[1pt]\campo{Roda}em paralelo com os demais de \cod{HandoffParalelo}};
\draw[fluxo] (n117.south) -- (n118.north);
\node[pmodelo, below=of n118] (n119) {{\sffamily\bfseries\small 100b\hspace{0.5em}Brief de assets}\hfill{\sffamily\tiny\color{white}\colorbox{unLavanda}{MODELO}}\\[1pt]{\tiny\ttfamily\color{unLavanda}A100bBriefAssets}\\[2pt]\campo{Faz}Inventário de assets, spec por tipo, critérios de reprovação, ordem de produção.\\[1pt]\campo{Objetivo}Seu leitor é curadoria e produção de conteúdo, não dev.\\[1pt]\campo{Entra}\cod{anexo\_design\_spec}\,\arq{design/00-anexo-design-spec.md}, \cod{camada\_caa}\,\arq{design/04b-camada-caa.md}, \cod{copy\_strings\_revisada}\,\arq{design/06-copy-strings-revisada.md}, \cod{jornada\_telas}\,\arq{design/02-jornada-telas.md}, \cod{one\_pager}\,\arq{produto/00-one-pager.md}, \cod{spec\_missao}\,\arq{produto/05-spec-missao.md}, \cod{wireframes\_cobertura}\,\arq{design/05-wireframes-cobertura.md}, \textit{+1 de controle}\\[1pt]\campo{Sai}\cod{brief\_assets}\,\arq{dev/01b-brief-assets.md}\\[1pt]\campo{Ciclo}\cod{HandoffLoop}\\[1pt]\campo{Roda}em paralelo com os demais de \cod{HandoffParalelo}};
\draw[fluxo] (n118.south) -- (n119.north);
\node[pportao, below=of n119] (n120) {{\sffamily\bfseries\small 101\hspace{0.5em}Contrato de telemetria e arquitetura (determinístico)}\hfill{\sffamily\tiny\color{white}\colorbox{unAzul!80!unNavy}{PORTÃO}}\\[1pt]{\tiny\ttfamily\color{unLavanda}A101TelemetriaCheck}\\[2pt]\campo{Faz}Evento = verbo · objeto · resultado · contexto (sem `dados`), medida com tabela de leitura, sem WebSocket, seis papéis, quota em `ia\_quotas`.\\[1pt]\campo{Objetivo}Gate determinístico de rastro: modo decide a verificação de tools/rastro.py.\\[1pt]\campo{Entra}\cod{handoff\_dev}\,\arq{dev/01-handoff-dev.md}, \textit{+1 de controle}\\[1pt]\campo{Sai}\cod{telemetria\_report}\,\arq{dev/01c-telemetria-check.md}, \cod{telemetria\_status}\\[1pt]\campo{Ciclo}\cod{HandoffLoop}};
\draw[fluxo] (n119.south) -- (n120.north);
\node[pportao, below=of n120] (n121) {{\sffamily\bfseries\small 102\hspace{0.5em}Juiz da telemetria}\hfill{\sffamily\tiny\color{white}\colorbox{unAzul!80!unNavy}{PORTÃO}}\\[1pt]{\tiny\ttfamily\color{unLavanda}A102JuizTelemetria}\\[2pt]\campo{Faz}Juiz de achados sobre `Contrato de telemetria e arquitetura`: retira o achado de uso vedado quando o texto só cita para vedar — com citação literal conferida.\\[1pt]\campo{Objetivo}O regex acha, o LLM decide, o código confere a citação (tools/juizachados.py).\\[1pt]\campo{Entra}\cod{handoff\_dev}\,\arq{dev/01-handoff-dev.md}, \textit{+1 de controle}\\[1pt]\campo{Sai}\cod{telemetria\_status}, \cod{telemetria\_report}\,\arq{dev/01c-telemetria-check.md}\\[1pt]\campo{Ciclo}\cod{HandoffLoop}};
\draw[fluxo] (n120.south) -- (n121.north);
\draw[loopfaixa] ([xshift=-0.25cm]n118.north west) -- ([xshift=-0.25cm]n121.south west);
\end{tikzpicture}
\end{center}
\newpage
\begin{center}
\begin{tikzpicture}[node distance=0.42cm]
\node[pportao] (n122) {{\sffamily\bfseries\small 101b\hspace{0.5em}Handoff monta, não decide (determinístico)}\hfill{\sffamily\tiny\color{white}\colorbox{unAzul!80!unNavy}{PORTÃO}}\\[1pt]{\tiny\ttfamily\color{unLavanda}A101bMontagemHandoff}\\[2pt]\campo{Faz}O handoff contra o contrato de produto (one-pager, spec de produto, lista branca do bloco 13, decisões do PO): número de regra sem origem, hipótese reescrita do \S{}17, meta fora do 9f, estado inventado, texto-chave reescrito e decisão de design citada como produto voltam ao V1.\\[1pt]\campo{Objetivo}Gate determinístico de rastro: modo decide a verificação de tools/rastro.py.\\[1pt]\campo{Entra}\textit{nada do estado da sessão (configuração fixa do passo)}\\[1pt]\campo{Sai}\cod{montagem\_handoff\_report}\,\arq{dev/01g-montagem-handoff.md}, \cod{montagem\_handoff\_status}\\[1pt]\campo{Ciclo}\cod{HandoffLoop}};
\node[psaida, below=of n122] (n123) {{\sffamily\bfseries\small 103\hspace{0.5em}Veredito do handoff (determinístico)}\hfill{\sffamily\tiny\color{white}\colorbox{sinalAmbar}{SAÍDA DE LOOP}}\\[1pt]{\tiny\ttfamily\color{unLavanda}A103HandoffCheck}\\[2pt]\campo{Faz}Contrato de telemetria e montagem conformes encerram o loop; reprovado devolve o handoff ao V1 com os relatórios.\\[1pt]\campo{Objetivo}Encerra um LoopAgent quando todos os gates da rodada passaram — ou quando o teto caiu.\\[1pt]\campo{Entra}\cod{telemetria\_status}, \cod{montagem\_handoff\_status}, \cod{handoff\_dev}\,\arq{dev/01-handoff-dev.md}, \cod{brief\_assets}\,\arq{dev/01b-brief-assets.md}\\[1pt]\campo{Sai}\textit{+1 de controle}\\[1pt]\campo{Ciclo}\cod{HandoffLoop}\\[1pt]\ciclo{$\circlearrowleft$ fim do ciclo: \cod{HandoffLoop} --- reprovado, volta ao passo 100a}};
\draw[fluxo] (n122.south) -- (n123.north);
\node[pmodelo, below=of n123] (n124) {{\sffamily\bfseries\small 104\hspace{0.5em}Banco de conteúdo}\hfill{\sffamily\tiny\color{white}\colorbox{unLavanda}{MODELO}}\\[1pt]{\tiny\ttfamily\color{unLavanda}A104BancoConteudo}\\[2pt]\ciclo{$\circlearrowright$ começa aqui: \cod{BancoLoop} (até 2 rodadas)}\\[1pt]\campo{Faz}O texto integral de cada item de conteúdo do brief (documentos fictícios, atas, tabelas cujas linhas não entregam a resposta, transcrições, roteiros) e a árvore de falas da IA simulada por estado. Status `escrito` — homologar é da curadoria. Salvo em `dev/01e-banco-conteudo.md`.\\[1pt]\campo{Objetivo}O brief de assets descreve o que a missão exibe.\\[1pt]\campo{Entra}\cod{anexo\_design\_spec}\,\arq{design/00-anexo-design-spec.md}, \cod{brief\_assets}\,\arq{dev/01b-brief-assets.md}, \cod{briefing}, \cod{camada\_caa}\,\arq{design/04b-camada-caa.md}, \cod{conteudo\_aprovado}, \cod{conteudo\_report}\,\arq{dev/01f-banco-conteudo-check.md}, \cod{copy\_strings\_revisada}\,\arq{design/06-copy-strings-revisada.md}, \cod{spec\_missao}\,\arq{produto/05-spec-missao.md}, \textit{+1 de controle}\\[1pt]\campo{Sai}\cod{banco\_conteudo}\,\arq{dev/01e-banco-conteudo.md}\\[1pt]\campo{Ciclo}\cod{BancoLoop}};
\draw[fluxo] (n123.south) -- (n124.north);
\node[pportao, below=of n124] (n125) {{\sffamily\bfseries\small 105\hspace{0.5em}Banco de conteúdo conferido (determinístico)}\hfill{\sffamily\tiny\color{white}\colorbox{unAzul!80!unNavy}{PORTÃO}}\\[1pt]{\tiny\ttfamily\color{unLavanda}A105ConteudoCheck}\\[2pt]\campo{Faz}ID de texto/documento/tabela/roteiro do brief sem bloco no banco bloqueia; Status do brief fora de `a produzir | rascunho | escrito` ("Homologado", "Pronto para corte") bloqueia; estado de fala da IA sem fala fica .\\[1pt]\campo{Objetivo}Gate determinístico de rastro: modo decide a verificação de tools/rastro.py.\\[1pt]\campo{Entra}\cod{banco\_conteudo}\,\arq{dev/01e-banco-conteudo.md}, \cod{brief\_assets}\,\arq{dev/01b-brief-assets.md}, \cod{one\_pager}\,\arq{produto/00-one-pager.md}\\[1pt]\campo{Sai}\cod{conteudo\_report}\,\arq{dev/01f-banco-conteudo-check.md}, \cod{conteudo\_status}\\[1pt]\campo{Ciclo}\cod{BancoLoop}};
\draw[fluxo] (n124.south) -- (n125.north);
\node[psaida, below=of n125] (n126) {{\sffamily\bfseries\small 106\hspace{0.5em}Saída do loop do banco de conteúdo (determinístico)}\hfill{\sffamily\tiny\color{white}\colorbox{sinalAmbar}{SAÍDA DE LOOP}}\\[1pt]{\tiny\ttfamily\color{unLavanda}A106BancoRodada}\\[2pt]\campo{Faz}Banco conferido encerra o loop; com bloqueio, o banco reescreve com o relatório.\\[1pt]\campo{Objetivo}Encerra um LoopAgent quando todos os gates da rodada passaram — ou quando o teto caiu.\\[1pt]\campo{Entra}\cod{conteudo\_status}\\[1pt]\campo{Sai}\textit{+1 de controle}\\[1pt]\campo{Ciclo}\cod{BancoLoop}\\[1pt]\ciclo{$\circlearrowleft$ fim do ciclo: \cod{BancoLoop} --- reprovado, volta ao passo 104}};
\draw[fluxo] (n125.south) -- (n126.north);
\node[pportao, below=of n126] (n127) {{\sffamily\bfseries\small 107\hspace{0.5em}Rastro de lacunas (determinístico)}\hfill{\sffamily\tiny\color{white}\colorbox{unAzul!80!unNavy}{PORTÃO}}\\[1pt]{\tiny\ttfamily\color{unLavanda}A107LacunaCheck}\\[2pt]\campo{Faz}Marcadores de lacuna em spec, design, handoff e assets, por dono: o que faltou de input, produto, design ou dev impede o PRD; plataforma vira pendência humana ().\\[1pt]\campo{Objetivo}Gate determinístico de rastro: modo decide a verificação de tools/rastro.py.\\[1pt]\campo{Entra}\cod{spec\_missao}\,\arq{produto/05-spec-missao.md}, \cod{anexo\_design\_spec}\,\arq{design/00-anexo-design-spec.md}, \cod{jornada\_telas}\,\arq{design/02-jornada-telas.md}, \cod{fluxo\_arestas}\,\arq{design/03-fluxo-arestas.md}, \cod{copy\_strings\_revisada}\,\arq{design/06-copy-strings-revisada.md}, \cod{wireframes\_cobertura}\,\arq{design/05-wireframes-cobertura.md}, \cod{stitch\_links}\,\arq{design/10-stitch-telas.md}, \cod{camada\_caa}\,\arq{design/04b-camada-caa.md}, \cod{ambiente\_simulado}\,\arq{design/00b-ambiente-simulado.md}, +3\\[1pt]\campo{Sai}\cod{rastro\_lacunas\_report}\,\arq{dev/01d-rastro-lacunas.md}, \cod{rastro\_lacunas\_status}};
\draw[fluxo] (n126.south) -- (n127.north);
\draw[loopfaixa] ([xshift=-0.25cm]n122.north west) -- ([xshift=-0.25cm]n123.south west);
\draw[loopfaixa] ([xshift=-0.25cm]n124.north west) -- ([xshift=-0.25cm]n126.south west);
\draw[loopvolta] ([xshift=-0.25cm]n126.south west) -- ++(-0.18,0) |- ([xshift=-0.25cm]n124.north west);
\end{tikzpicture}
\end{center}
\newpage
\begin{center}
\begin{tikzpicture}[node distance=0.42cm]
\node[pcodigo] (n128) {{\sffamily\bfseries\small 108\hspace{0.5em}Escopo da banca de dev (determinístico)}\hfill{\sffamily\tiny\color{white}\colorbox{unNavy}{CÓDIGO}}\\[1pt]{\tiny\ttfamily\color{unLavanda}A108EscopoBancaDev}\\[2pt]\ciclo{$\circlearrowright$ começa aqui: \cod{BancaLoop} (até 2 rodadas)}\\[1pt]\campo{Faz}Grava o escopo da banca dentro do dev: dados e técnica, conserto só de documento de dev.\\[1pt]\campo{Objetivo}Grava em escopobanca onde a banca está rodando e o que ela pode fazer ali (tools/banca.ESCOPOS).\\[1pt]\campo{Entra}\textit{nada do estado da sessão (configuração fixa do passo)}\\[1pt]\campo{Sai}\cod{escopo\_banca}\\[1pt]\campo{Ciclo}\cod{BancaLoop}};
\node[pmodelo, below=of n128] (n129) {{\sffamily\bfseries\small 109a\hspace{0.5em}Banca: Dados, Telemetria \& LGPD}\hfill{\sffamily\tiny\color{white}\colorbox{unLavanda}{MODELO}}\\[1pt]{\tiny\ttfamily\color{unLavanda}A109aBancaDadosLgpd}\\[2pt]\campo{Faz}DAD-01…15.\\[1pt]\campo{Objetivo}A mesma banca roda em quatro posições, com os mesmos critérios — fim de produto, fim de design, entrada do dev e dentro do dev.\\[1pt]\campo{Entra}\cod{anexo\_design\_spec}\,\arq{design/00-anexo-design-spec.md}, \cod{banco\_conteudo}\,\arq{dev/01e-banco-conteudo.md}, \cod{brief\_assets}\,\arq{dev/01b-brief-assets.md}, \cod{brief\_oficio}\,\arq{design/01-briefing-oficio.md}, \cod{camada\_caa}\,\arq{design/04b-camada-caa.md}, \cod{conteudo\_aprovado}, \cod{copy\_strings\_revisada}\,\arq{design/06-copy-strings-revisada.md}, \cod{escopo\_banca}, \cod{fluxo\_arestas}\,\arq{design/03-fluxo-arestas.md}, +6\\[1pt]\campo{Sai}\cod{validacao\_dados\_lgpd}\,\arq{dev/02d-banca-dados-lgpd.md}\\[1pt]\campo{Ciclo}\cod{BancaLoop}\\[1pt]\campo{Roda}em paralelo com os demais de \cod{BancaValidacaoDev}};
\draw[fluxo] (n128.south) -- (n129.north);
\node[pmodelo, below=of n129] (n130) {{\sffamily\bfseries\small 109b\hspace{0.5em}Banca: Técnica \& Arquitetura}\hfill{\sffamily\tiny\color{white}\colorbox{unLavanda}{MODELO}}\\[1pt]{\tiny\ttfamily\color{unLavanda}A109bBancaTecnicaArquitetura}\\[2pt]\campo{Faz}TEC-01…13.\\[1pt]\campo{Objetivo}A mesma banca roda em quatro posições, com os mesmos critérios — fim de produto, fim de design, entrada do dev e dentro do dev.\\[1pt]\campo{Entra}\cod{anexo\_design\_spec}\,\arq{design/00-anexo-design-spec.md}, \cod{banco\_conteudo}\,\arq{dev/01e-banco-conteudo.md}, \cod{brief\_assets}\,\arq{dev/01b-brief-assets.md}, \cod{brief\_oficio}\,\arq{design/01-briefing-oficio.md}, \cod{camada\_caa}\,\arq{design/04b-camada-caa.md}, \cod{conteudo\_aprovado}, \cod{copy\_strings\_revisada}\,\arq{design/06-copy-strings-revisada.md}, \cod{escopo\_banca}, \cod{fluxo\_arestas}\,\arq{design/03-fluxo-arestas.md}, +6\\[1pt]\campo{Sai}\cod{validacao\_tecnica\_arquitetura}\,\arq{dev/02e-banca-tecnica-arquitetura.md}\\[1pt]\campo{Ciclo}\cod{BancaLoop}\\[1pt]\campo{Roda}em paralelo com os demais de \cod{BancaValidacaoDev}};
\draw[fluxo] (n129.south) -- (n130.north);
\node[pportao, below=of n130] (n131) {{\sffamily\bfseries\small 110\hspace{0.5em}Consolidação da banca}\hfill{\sffamily\tiny\color{white}\colorbox{unAzul!80!unNavy}{PORTÃO}}\\[1pt]{\tiny\ttfamily\color{unLavanda}A110ConsolidacaoBanca}\\[2pt]\campo{Faz}Determinístico: veredito = a pior das 5 dimensões; bloqueantes lidos de `references/criterios-dev/`.\\[1pt]\campo{Objetivo}Determinístico: o veredito da banca é a pior das cinco dimensões, contada em código.\\[1pt]\campo{Entra}\cod{validacao\_dados\_lgpd}\,\arq{dev/02d-banca-dados-lgpd.md}, \cod{validacao\_tecnica\_arquitetura}\,\arq{dev/02e-banca-tecnica-arquitetura.md}, \cod{spec\_missao}\,\arq{produto/05-spec-missao.md}, \cod{one\_pager}\,\arq{produto/00-one-pager.md}\\[1pt]\campo{Sai}\cod{validacao\_dev\_report}\,\arq{dev/03-banca-consolidada.md}, \cod{dev\_status}\\[1pt]\campo{Ciclo}\cod{BancaLoop}};
\draw[fluxo] (n130.south) -- (n131.north);
\node[pportao, below=of n131] (n132) {{\sffamily\bfseries\small 111\hspace{0.5em}Veredito de dev (determinístico)}\hfill{\sffamily\tiny\color{white}\colorbox{unAzul!80!unNavy}{PORTÃO}}\\[1pt]{\tiny\ttfamily\color{unLavanda}A111ValidacaoDevCheck}\\[2pt]\campo{Faz}Lê o envelope da banca.\\[1pt]\campo{Objetivo}Lê o veredito de um relatório e decide se o pipeline segue, repete ou para.\\[1pt]\campo{Entra}\cod{validacao\_dev\_report}\,\arq{dev/03-banca-consolidada.md}\\[1pt]\campo{Sai}\cod{dev\_status}, \textit{+1 de controle}\\[1pt]\campo{Ciclo}\cod{BancaLoop}};
\draw[fluxo] (n131.south) -- (n132.north);
\node[pcodigo, below=of n132] (n133) {{\sffamily\bfseries\small 112\hspace{0.5em}Retrabalho da banca (determinístico)}\hfill{\sffamily\tiny\color{white}\colorbox{unNavy}{CÓDIGO}}\\[1pt]{\tiny\ttfamily\color{unLavanda}A112DespachoBanca}\\[2pt]\campo{Faz}Reexecuta o dono de cada achado da banca — inclusive quando o dono está na etapa de design. Achado de produto vira pendência nomeada.\\[1pt]\campo{Objetivo}Re-executa apenas os agentes endereçados pelos tickets de um gate.\\[1pt]\campo{Entra}\cod{validacao\_dev\_report}\,\arq{dev/03-banca-consolidada.md}, \textit{+1 de controle}\\[1pt]\campo{Sai}\textit{+1 de controle}\\[1pt]\campo{Ciclo}\cod{BancaLoop}\\[1pt]\ciclo{$\circlearrowleft$ fim do ciclo: \cod{BancaLoop} --- reprovado, volta ao passo 108}};
\draw[fluxo] (n132.south) -- (n133.north);
\draw[loopfaixa] ([xshift=-0.25cm]n128.north west) -- ([xshift=-0.25cm]n133.south west);
\draw[loopvolta] ([xshift=-0.25cm]n133.south west) -- ++(-0.18,0) |- ([xshift=-0.25cm]n128.north west);
\end{tikzpicture}
\end{center}
\newpage
\begin{center}
\begin{tikzpicture}[node distance=0.42cm]
\node[pmodelo] (n134) {{\sffamily\bfseries\small 113\hspace{0.5em}PRD padrão ouro}\hfill{\sffamily\tiny\color{white}\colorbox{unLavanda}{MODELO}}\\[1pt]{\tiny\ttfamily\color{unLavanda}A113PrdWriter}\\[2pt]\ciclo{$\circlearrowright$ começa aqui: \cod{PrdLoop} (até 2 rodadas)}\\[1pt]\campo{Faz}PRD completo com contexto arquitetural da plataforma — entregável do pipeline.\\[1pt]\campo{Objetivo}Você escreve o PRD padrão ouro da missão: o documento que um agente de codificação vai receber para implementar a missão no monorepo plataforma-unooria sem adivinhar nada — e sem decidir nada de produto, pedagogia ou design: o…\\[1pt]\campo{Entra}\cod{anexo\_design\_spec}\,\arq{design/00-anexo-design-spec.md}, \cod{banco\_conteudo}\,\arq{dev/01e-banco-conteudo.md}, \cod{brief\_assets}\,\arq{dev/01b-brief-assets.md}, \cod{brief\_oficio}\,\arq{design/01-briefing-oficio.md}, \cod{camada\_caa}\,\arq{design/04b-camada-caa.md}, \cod{coerencia\_prd\_report}\,\arq{dev/04d-coerencia-prd.md}, \cod{conteudo\_aprovado\_prd\_report}\,\arq{dev/04c-conteudo-aprovado-prd.md}, \cod{copy\_strings\_revisada}\,\arq{design/06-copy-strings-revisada.md}, \cod{fluxo\_arestas}\,\arq{design/03-fluxo-arestas.md}, +17\\[1pt]\campo{Sai}\cod{prd\_final}\,\arq{PRD.md}\\[1pt]\campo{Ciclo}\cod{PrdLoop}};
\node[pportao, below=of n134] (n135) {{\sffamily\bfseries\small 114\hspace{0.5em}Pendências do PRD (determinístico)}\hfill{\sffamily\tiny\color{white}\colorbox{unAzul!80!unNavy}{PORTÃO}}\\[1pt]{\tiny\ttfamily\color{unLavanda}A114PrdCheck}\\[2pt]\campo{Faz}Marcadores que o redator deixou no PRD, mais a consistência interna do documento (prefixo de tabela único, requisito P0 com teste e tarefa, evento com tabela leitora, caminho único). lacuna de produto/design/dev marca o PRD como não pronto, plataforma deixa .\\[1pt]\campo{Objetivo}Gate determinístico de rastro: modo decide a verificação de tools/rastro.py.\\[1pt]\campo{Entra}\cod{prd\_final}\,\arq{PRD.md}, \cod{banco\_conteudo}\,\arq{dev/01e-banco-conteudo.md}, \cod{one\_pager}\,\arq{produto/00-one-pager.md}, \cod{handoff\_dev}\,\arq{dev/01-handoff-dev.md}\\[1pt]\campo{Sai}\cod{prd\_check\_report}\,\arq{dev/04-pendencias-do-prd.md}, \cod{prd\_status}\\[1pt]\campo{Ciclo}\cod{PrdLoop}};
\draw[fluxo] (n134.south) -- (n135.north);
\node[pportao, below=of n135] (n136) {{\sffamily\bfseries\small 115\hspace{0.5em}Coerência KPI $\leftrightarrow$ evento $\leftrightarrow$ tela (determinístico)}\hfill{\sffamily\tiny\color{white}\colorbox{unAzul!80!unNavy}{PORTÃO}}\\[1pt]{\tiny\ttfamily\color{unLavanda}A115KpisCheck}\\[2pt]\campo{Faz}Todo KPI com pergunta, decisão, corte e evento que dispara na tela que ele nomeia; os sinais universais das transversais de dono missão na tabela de eventos, com a IA ao estudante decidida pela spec.\\[1pt]\campo{Objetivo}Gate determinístico de rastro: modo decide a verificação de tools/rastro.py.\\[1pt]\campo{Entra}\cod{prd\_final}\,\arq{PRD.md}, \cod{one\_pager}\,\arq{produto/00-one-pager.md}\\[1pt]\campo{Sai}\cod{kpis\_report}\,\arq{dev/04b-kpis-check.md}, \cod{kpis\_status}\\[1pt]\campo{Ciclo}\cod{PrdLoop}};
\draw[fluxo] (n135.south) -- (n136.north);
\node[pportao, below=of n136] (n137) {{\sffamily\bfseries\small 116\hspace{0.5em}PRD $\times$ one-pager aprovado (determinístico)}\hfill{\sffamily\tiny\color{white}\colorbox{unAzul!80!unNavy}{PORTÃO}}\\[1pt]{\tiny\ttfamily\color{unLavanda}A116ConteudoAprovadoPrd}\\[2pt]\campo{Faz}As habilidades do bloco 11 na matriz da seção 14, um KPI por habilidade, o tempo em tela do one-pager, o podcast com requisito funcional e só caminhos de artefato que existem na sessão. Divergência volta ao redator.\\[1pt]\campo{Objetivo}Gate determinístico de rastro: modo decide a verificação de tools/rastro.py.\\[1pt]\campo{Entra}\cod{spec\_missao}\,\arq{produto/05-spec-missao.md}, \cod{one\_pager}\,\arq{produto/00-one-pager.md}, \cod{prd\_final}\,\arq{PRD.md}, \cod{one\_pager\_resumido}\,\arq{produto/00d-one-pager-resumido.md}, \cod{one\_pager\_pedagogico}\,\arq{one-pager-pedagogico.md}\\[1pt]\campo{Sai}\cod{conteudo\_aprovado\_prd\_report}\,\arq{dev/04c-conteudo-aprovado-prd.md}, \cod{conteudo\_aprovado\_prd\_status}\\[1pt]\campo{Ciclo}\cod{PrdLoop}};
\draw[fluxo] (n136.south) -- (n137.north);
\node[pportao, below=of n137] (n138) {{\sffamily\bfseries\small 116b\hspace{0.5em}PRD monta, não decide (determinístico)}\hfill{\sffamily\tiny\color{white}\colorbox{unAzul!80!unNavy}{PORTÃO}}\\[1pt]{\tiny\ttfamily\color{unLavanda}A116bMontagemPrd}\\[2pt]\campo{Faz}O corpo do PRD contra o contrato de produto: número de regra com origem no contrato, hipóteses do \S{}17 e metas do 9f com os números deles, tetos do \S{}15 presentes, estados da \S{}6 iguais aos do contrato, nenhuma decisão de design citada como produto. Divergência do próprio contrato só passa registrada na seção 19 com dono produto.\\[1pt]\campo{Objetivo}Gate determinístico de rastro: modo decide a verificação de tools/rastro.py.\\[1pt]\campo{Entra}\textit{nada do estado da sessão (configuração fixa do passo)}\\[1pt]\campo{Sai}\cod{montagem\_prd\_report}\,\arq{dev/04e-montagem-prd.md}, \cod{montagem\_prd\_status}\\[1pt]\campo{Ciclo}\cod{PrdLoop}};
\draw[fluxo] (n137.south) -- (n138.north);
\draw[loopfaixa] ([xshift=-0.25cm]n134.north west) -- ([xshift=-0.25cm]n138.south west);
\end{tikzpicture}
\end{center}
\newpage
\begin{center}
\begin{tikzpicture}[node distance=0.42cm]
\node[pportao] (n139) {{\sffamily\bfseries\small 117\hspace{0.5em}Coerência do PRD com o one-pager aprovado}\hfill{\sffamily\tiny\color{white}\colorbox{unAzul!80!unNavy}{PORTÃO}}\\[1pt]{\tiny\ttfamily\color{unLavanda}A117CoerenciaPrd}\\[2pt]\campo{Faz}Um juiz aponta onde o corpo do PRD afirma o contrário do conteúdo aprovado — mecânica, número, regra, texto —, citando os dois lados; o código confere as citações. Contradição comprovada volta ao redator.\\[1pt]\campo{Objetivo}O documento contradiz o conteúdo aprovado?\\[1pt]\campo{Entra}\cod{prd\_final}\,\arq{PRD.md}, \cod{conteudo\_aprovado}, \cod{contradicoes\_documento}, \cod{contradicoes\_referencia}, \cod{spec\_missao}\,\arq{produto/05-spec-missao.md}, \cod{one\_pager}\,\arq{produto/00-one-pager.md}\\[1pt]\campo{Sai}\cod{coerencia\_prd\_report}\,\arq{dev/04d-coerencia-prd.md}, \cod{coerencia\_prd\_status}\\[1pt]\campo{Ciclo}\cod{PrdLoop}};
\node[psaida, below=of n139] (n140) {{\sffamily\bfseries\small 118\hspace{0.5em}Veredito do PRD (determinístico)}\hfill{\sffamily\tiny\color{white}\colorbox{sinalAmbar}{SAÍDA DE LOOP}}\\[1pt]{\tiny\ttfamily\color{unLavanda}A118PrdRodada}\\[2pt]\campo{Faz}PRD consistente encerra o loop; reprovado devolve ao redator com o relatório do gate, até o teto de rodadas.\\[1pt]\campo{Objetivo}Encerra um LoopAgent quando todos os gates da rodada passaram — ou quando o teto caiu.\\[1pt]\campo{Entra}\cod{prd\_status}, \cod{kpis\_status}, \cod{conteudo\_aprovado\_prd\_status}, \cod{montagem\_prd\_status}, \cod{coerencia\_prd\_status}, \cod{prd\_final}\,\arq{PRD.md}\\[1pt]\campo{Sai}\textit{+1 de controle}\\[1pt]\campo{Ciclo}\cod{PrdLoop}\\[1pt]\ciclo{$\circlearrowleft$ fim do ciclo: \cod{PrdLoop} --- reprovado, volta ao passo 113}};
\draw[fluxo] (n139.south) -- (n140.north);
\node[pcodigo, below=of n140] (n141) {{\sffamily\bfseries\small 119\hspace{0.5em}Anexo do conteúdo aprovado no PRD (determinístico)}\hfill{\sffamily\tiny\color{white}\colorbox{unNavy}{CÓDIGO}}\\[1pt]{\tiny\ttfamily\color{unLavanda}A119AnexoPrd}\\[2pt]\campo{Faz}Cola no fim do PRD, literais, as seções do one-pager aprovado que o gabarito não tinha onde pôr: currículo, papel do professor, tempos e cenários de aula, podcast com roteiro, contrato de evidência e hipóteses. E põe na seção 19, por código, as pendências de governança da liberação de produto — liberado sem a conta produto, exceção do admin com gate vermelho — como estão agora, também na correção pontual.\\[1pt]\campo{Objetivo}Cola no fim de chave as seções do one-pager aprovado, literais (tools/conteudoaprovado.py).\\[1pt]\campo{Entra}\cod{one\_pager}\,\arq{produto/00-one-pager.md}\\[1pt]\campo{Sai}\cod{prd\_final}\,\arq{PRD.md}};
\draw[fluxo] (n140.south) -- (n141.north);
\node[pmodelo, below=of n141] (n142) {{\sffamily\bfseries\small 120\hspace{0.5em}One-pager pedagógico (validação humana)}\hfill{\sffamily\tiny\color{white}\colorbox{unLavanda}{MODELO}}\\[1pt]{\tiny\ttfamily\color{unLavanda}A120OnePagerPedagogico}\\[2pt]\ciclo{$\circlearrowright$ começa aqui: \cod{OnePagerPedagogicoLoop} (até 2 rodadas)}\\[1pt]\campo{Faz}A missão inteira em linguagem de sala de aula, com o papel do professor, os links do protótipo e as perguntas de aprovação. Salvo em Markdown e `.docx`.\\[1pt]\campo{Objetivo}O PRD que veio antes é para um agente de codificação: cheio de nome de arquivo, tabela e contrato.\\[1pt]\campo{Entra}\cod{ambiente\_simulado}\,\arq{design/00b-ambiente-simulado.md}, \cod{anexo\_design\_spec}\,\arq{design/00-anexo-design-spec.md}, \cod{avaliacao\_pedagogica\_report}\,\arq{produto/02-avaliacao-pedagogica.md}, \cod{bncc\_report}\,\arq{produto/05b-conformidade-bncc.md}, \cod{brief\_oficio}\,\arq{design/01-briefing-oficio.md}, \cod{camada\_caa}\,\arq{design/04b-camada-caa.md}, \cod{coerencia\_pedagogico\_report}, \cod{conteudo\_aprovado\_pedagogico\_report}\,\arq{dev/05-conteudo-aprovado-pedagogico.md}, \cod{copy\_strings\_revisada}\,\arq{design/06-copy-strings-revisada.md}, +15\\[1pt]\campo{Sai}\cod{one\_pager\_pedagogico}\,\arq{one-pager-pedagogico.md}\\[1pt]\campo{Ciclo}\cod{OnePagerPedagogicoLoop}};
\draw[fluxo] (n141.south) -- (n142.north);
\node[pportao, below=of n142] (n143) {{\sffamily\bfseries\small 121\hspace{0.5em}Pedagógico $\times$ one-pager aprovado (determinístico)}\hfill{\sffamily\tiny\color{white}\colorbox{unAzul!80!unNavy}{PORTÃO}}\\[1pt]{\tiny\ttfamily\color{unLavanda}A121ConteudoAprovadoPedagogico}\\[2pt]\campo{Faz}As habilidades do bloco 11, o tempo em tela, o podcast e os caminhos de artefato do pedagógico contra o one-pager aprovado. Divergência volta ao redator.\\[1pt]\campo{Objetivo}Gate determinístico de rastro: modo decide a verificação de tools/rastro.py.\\[1pt]\campo{Entra}\cod{spec\_missao}\,\arq{produto/05-spec-missao.md}, \cod{one\_pager}\,\arq{produto/00-one-pager.md}, \cod{prd\_final}\,\arq{PRD.md}, \cod{one\_pager\_resumido}\,\arq{produto/00d-one-pager-resumido.md}, \cod{one\_pager\_pedagogico}\,\arq{one-pager-pedagogico.md}\\[1pt]\campo{Sai}\cod{conteudo\_aprovado\_pedagogico\_report}\,\arq{dev/05-conteudo-aprovado-pedagogico.md}, \cod{conteudo\_aprovado\_pedagogico\_status}\\[1pt]\campo{Ciclo}\cod{OnePagerPedagogicoLoop}};
\draw[fluxo] (n142.south) -- (n143.north);
\draw[loopfaixa] ([xshift=-0.25cm]n139.north west) -- ([xshift=-0.25cm]n140.south west);
\draw[loopfaixa] ([xshift=-0.25cm]n142.north west) -- ([xshift=-0.25cm]n143.south west);
\end{tikzpicture}
\end{center}
\newpage
\begin{center}
\begin{tikzpicture}[node distance=0.42cm]
\node[pportao] (n144) {{\sffamily\bfseries\small 122\hspace{0.5em}Coerência do pedagógico com o one-pager aprovado}\hfill{\sffamily\tiny\color{white}\colorbox{unAzul!80!unNavy}{PORTÃO}}\\[1pt]{\tiny\ttfamily\color{unLavanda}A122CoerenciaPedagogico}\\[2pt]\campo{Faz}O mesmo juiz do PRD, sobre o documento que o time pedagógico aprova: contradição com o conteúdo aprovado, com as duas citações conferidas, volta ao redator.\\[1pt]\campo{Objetivo}O documento contradiz o conteúdo aprovado?\\[1pt]\campo{Entra}\cod{one\_pager\_pedagogico}\,\arq{one-pager-pedagogico.md}, \cod{conteudo\_aprovado}, \cod{contradicoes\_documento}, \cod{contradicoes\_referencia}, \cod{spec\_missao}\,\arq{produto/05-spec-missao.md}, \cod{one\_pager}\,\arq{produto/00-one-pager.md}\\[1pt]\campo{Sai}\cod{coerencia\_pedagogico\_report}, \cod{coerencia\_pedagogico\_status}\\[1pt]\campo{Ciclo}\cod{OnePagerPedagogicoLoop}};
\node[psaida, below=of n144] (n145) {{\sffamily\bfseries\small 123\hspace{0.5em}Veredito do pedagógico (determinístico)}\hfill{\sffamily\tiny\color{white}\colorbox{sinalAmbar}{SAÍDA DE LOOP}}\\[1pt]{\tiny\ttfamily\color{unLavanda}A123PedagogicoRodada}\\[2pt]\campo{Faz}Pedagógico fiel encerra o loop; divergente volta ao redator com o relatório, até o teto de rodadas.\\[1pt]\campo{Objetivo}Encerra um LoopAgent quando todos os gates da rodada passaram — ou quando o teto caiu.\\[1pt]\campo{Entra}\cod{conteudo\_aprovado\_pedagogico\_status}, \cod{coerencia\_pedagogico\_status}, \cod{one\_pager\_pedagogico}\,\arq{one-pager-pedagogico.md}\\[1pt]\campo{Sai}\textit{+1 de controle}\\[1pt]\campo{Ciclo}\cod{OnePagerPedagogicoLoop}\\[1pt]\ciclo{$\circlearrowleft$ fim do ciclo: \cod{OnePagerPedagogicoLoop} --- reprovado, volta ao passo 120}};
\draw[fluxo] (n144.south) -- (n145.north);
\node[pcodigo, below=of n145] (n146) {{\sffamily\bfseries\small 124\hspace{0.5em}Anexo do conteúdo aprovado no pedagógico (determinístico)}\hfill{\sffamily\tiny\color{white}\colorbox{unNavy}{CÓDIGO}}\\[1pt]{\tiny\ttfamily\color{unLavanda}A124AnexoPedagogico}\\[2pt]\campo{Faz}Cola no fim do pedagógico, literais, os tempos e cenários de aula, o podcast com o roteiro e o contrato de evidência do one-pager aprovado — o texto que o time pedagógico aprova.\\[1pt]\campo{Objetivo}Cola no fim de chave as seções do one-pager aprovado, literais (tools/conteudoaprovado.py).\\[1pt]\campo{Entra}\cod{one\_pager}\,\arq{produto/00-one-pager.md}\\[1pt]\campo{Sai}\cod{one\_pager\_pedagogico}\,\arq{one-pager-pedagogico.md}};
\draw[fluxo] (n145.south) -- (n146.north);
\node[pportao, below=of n146] (n147) {{\sffamily\bfseries\small 125\hspace{0.5em}Matriz de rastreabilidade (determinístico)}\hfill{\sffamily\tiny\color{white}\colorbox{unAzul!80!unNavy}{PORTÃO}}\\[1pt]{\tiny\ttfamily\color{unLavanda}A125Rastreabilidade}\\[2pt]\campo{Faz}Cada requisito do produto procurado em cada artefato de design e de dev, até o PRD e o pedagógico; o que se perdeu sai com o dono.\\[1pt]\campo{Objetivo}Determinístico: a matriz requisito do produto $\times$ artefato (tools/rastreabilidade.py).\\[1pt]\campo{Entra}\cod{spec\_missao}\,\arq{produto/05-spec-missao.md}, \cod{one\_pager}\,\arq{produto/00-one-pager.md}\\[1pt]\campo{Sai}\cod{rastreabilidade\_dev\_report}\,\arq{matriz-de-rastreabilidade.md}, \cod{rastreabilidade\_dev\_status}};
\draw[fluxo] (n146.south) -- (n147.north);
\node[pportao, below=of n147] (n148) {{\sffamily\bfseries\small 126\hspace{0.5em}Fidelidade ao refazer de dev}\hfill{\sffamily\tiny\color{white}\colorbox{unAzul!80!unNavy}{PORTÃO}}\\[1pt]{\tiny\ttfamily\color{unLavanda}A126FidelidadeDev}\\[2pt]\campo{Faz}Cada pedido de refazer feito ao dev (TRAVAS, AJUSTES, VALIDAÇÃO, frases entre aspas) conferido no handoff, no brief, no banco, no PRD e no pedagógico: o juiz cita, o código confere. Sem pedido de dev, não roda.  recusa o de acordo.\\[1pt]\campo{Objetivo}O que o PO pediu e o briefing fixou está nos documentos?\\[1pt]\campo{Entra}\cod{handoff\_dev}\,\arq{dev/01-handoff-dev.md}, \cod{brief\_assets}\,\arq{dev/01b-brief-assets.md}, \cod{banco\_conteudo}\,\arq{dev/01e-banco-conteudo.md}, \cod{prd\_final}\,\arq{PRD.md}, \cod{one\_pager\_pedagogico}\,\arq{one-pager-pedagogico.md}, \cod{divisao\_do\_pedido}, \cod{pedidos\_do\_po}, \cod{pedido\_correcao}, \cod{etapa\_da\_correcao}, +2, \textit{+3 de controle}\\[1pt]\campo{Sai}\cod{fidelidade\_dev\_report}\,\arq{dev/06-fidelidade-refazer-dev.md}, \cod{fidelidade\_dev\_status}};
\draw[fluxo] (n147.south) -- (n148.north);
\draw[loopfaixa] ([xshift=-0.25cm]n144.north west) -- ([xshift=-0.25cm]n145.south west);
\end{tikzpicture}
\end{center}
\newpage
\markboth{Transversal}{}
{\sffamily\bfseries\Large\color{unNavy}Transversal}\\[0.3em]
\begin{center}
\begin{tikzpicture}[node distance=0.42cm]
\node[pcodigo] (n149) {{\sffamily\bfseries\small 127\hspace{0.5em}Diagnóstico da execução (determinístico)}\hfill{\sffamily\tiny\color{white}\colorbox{unNavy}{CÓDIGO}}\\[1pt]{\tiny\ttfamily\color{unLavanda}A127Diagnostico}\\[2pt]\campo{Faz}Como cada portão se saiu, o que trava e de quem é cada conserto, o que saiu provisório e que documento faltou. Sem LLM — roda mesmo quando tudo o mais falhou.\\[1pt]\campo{Objetivo}Fecha a execução com o dossiê do que passou, do que passou com ressalva e do que travou.\\[1pt]\campo{Entra}\cod{spec\_missao}\,\arq{produto/05-spec-missao.md}, \cod{one\_pager}\,\arq{produto/00-one-pager.md}\\[1pt]\campo{Sai}\cod{diagnostico\_execucao}\,\arq{diagnostico-execucao.md}, \cod{diagnostico\_status}};
\node[pmodelo, below=of n149] (n150) {{\sffamily\bfseries\small 128\hspace{0.5em}Aprendizado}\hfill{\sffamily\tiny\color{white}\colorbox{unLavanda}{MODELO}}\\[1pt]{\tiny\ttfamily\color{unLavanda}A128Aprendizado}\\[2pt]\campo{Faz}A leitura em prosa do diagnóstico, mais catálogo, padrões aprendidos, lacunas com dono/classe e candidatos a critério.\\[1pt]\campo{Objetivo}Você tem duas funções, nesta ordem de importância:\\[1pt]\campo{Entra}\cod{ambiente\_simulado}\,\arq{design/00b-ambiente-simulado.md}, \cod{ambiente\_status}, \cod{anexo\_design\_spec}\,\arq{design/00-anexo-design-spec.md}, \cod{avaliacao\_pedagogica\_report}\,\arq{produto/02-avaliacao-pedagogica.md}, \cod{bncc\_report}\,\arq{produto/05b-conformidade-bncc.md}, \cod{bncc\_status}, \cod{brief\_oficio}\,\arq{design/01-briefing-oficio.md}, \cod{catalogo\_missoes}, \cod{design\_status}, +23, \textit{+2 de controle}\\[1pt]\campo{Sai}\cod{aprendizado\_registro}\,\arq{aprendizado.md}};
\draw[fluxo] (n149.south) -- (n150.north);
\node[pcodigo, below=of n150] (n151) {{\sffamily\bfseries\small 129\hspace{0.5em}Registro no catálogo (determinístico)}\hfill{\sffamily\tiny\color{white}\colorbox{unNavy}{CÓDIGO}}\\[1pt]{\tiny\ttfamily\color{unLavanda}A129RegistroCatalogo}\\[2pt]\campo{Faz}A memória entre execuções: registra esta missão no índice, mostra as mais próximas e acusa slug repetido. Sem LLM — status derivado dos gates, promover é ato humano.\\[1pt]\campo{Objetivo}Registra a missão desta execução na memória do gerador, ao fim da esteira.\\[1pt]\campo{Entra}\cod{spec\_missao}\,\arq{produto/05-spec-missao.md}, \cod{one\_pager}\,\arq{produto/00-one-pager.md}\\[1pt]\campo{Sai}\cod{ficha\_catalogo}\,\arq{catalogo-ficha.md}, \cod{catalogo\_status}};
\draw[fluxo] (n150.south) -- (n151.north);
\end{tikzpicture}
\end{center}
\newpage
\end{document}
